\documentclass[12pt,a4paper,reqno]{amsart}
%\documentclass[reqno,10pt]{amsart}
\usepackage[dvipsnames]{xcolor}
\usepackage{amsmath,amssymb,graphics,epsfig,color,enumerate,mathrsfs}
%%%%%%%%%%%%%%%%%%%%%%%%%%%%%%%%%%%%%%%%%%%%%%%%%%%
\usepackage{dsfont}  \usepackage{verbatim}
\textwidth= 14. cm
%\usepackage[color]{showkeys}
\definecolor{refkey}{gray}{.75}
\definecolor{labelkey}{gray}{.5}


\newtheorem{Theorem}{Theorem}[section]
\newtheorem{Fact}[Theorem]{Fact}
\newtheorem{TheoremA}{Theorem}
\newtheorem{Lemma}[Theorem]{Lemma}
\newtheorem{Proposition}[Theorem]{Proposition}
\newtheorem{Corollary}[Theorem]{Corollary}
\newtheorem{Remark}[Theorem]{Remark}
\newtheorem{Remarks}[Theorem]{Remarks}
\newtheorem{Exercise}{Exercise}
\newtheorem{Problem}[Theorem]{Problem}
\newtheorem{Claim}[Theorem]{Claim}
\newtheorem{Definition}[Theorem]{Definition}
\newtheorem{Warning}[Theorem]{Warning}
\newtheorem{Conjecture}{Conjecture}


  

\newcommand{\ra}{\rangle}
\newcommand{\la}{\langle}



 \definecolor{darkgreen}{rgb}{0,0.4,0}


\definecolor{light}{gray}{0.9}
\newcommand{\pecetta}[1]{
$\phantom .$
\smallskip
\par\noindent
\colorbox{light}{\begin{minipage}{15cm}#1\end{minipage}}
\smallskip
}

 
%%%%%%%%%%%%%%%%%%%%%%%%%%%%%%%%%%%%%%%%%%%%%%%%%%%%%%%%%%%%%%%%%%%%%%%%%%%%%%
%%%%%%%%%% Calligraphic letters
%%%%%%%%%%%%%%%%%%%%%%%%%%%%%%%%%%%%%%%%%%%%%%%%%%%%%%%%%%%%%%%%%%%%%%%%%%%%%%

\newcommand{\cA}{\ensuremath{\mathcal A}}
\newcommand{\cB}{\ensuremath{\mathcal B}}
\newcommand{\cC}{\ensuremath{\mathcal C}}
\newcommand{\cD}{\ensuremath{\mathcal D}}
\newcommand{\cE}{\ensuremath{\mathcal E}}
\newcommand{\cF}{\ensuremath{\mathcal F}}
\newcommand{\cG}{\ensuremath{\mathcal G}}
\newcommand{\cH}{\ensuremath{\mathcal H}}
\newcommand{\cI}{\ensuremath{\mathcal I}}
\newcommand{\cJ}{\ensuremath{\mathcal J}}
\newcommand{\cK}{\ensuremath{\mathcal K}}
\newcommand{\cL}{\ensuremath{\mathcal L}}
\newcommand{\cM}{\ensuremath{\mathcal M}}
\newcommand{\cN}{\ensuremath{\mathcal N}}
\newcommand{\cO}{\ensuremath{\mathcal O}}
\newcommand{\cP}{\ensuremath{\mathcal P}}
\newcommand{\cQ}{\ensuremath{\mathcal Q}}
\newcommand{\cR}{\ensuremath{\mathcal R}}
\newcommand{\cS}{\ensuremath{\mathcal S}}
\newcommand{\cT}{\ensuremath{\mathcal T}}
\newcommand{\cU}{\ensuremath{\mathcal U}}
\newcommand{\cV}{\ensuremath{\mathcal V}}
\newcommand{\cW}{\ensuremath{\mathcal W}}
\newcommand{\cX}{\ensuremath{\mathcal X}}
\newcommand{\cY}{\ensuremath{\mathcal Y}}
\newcommand{\cZ}{\ensuremath{\mathcal Z}}

%%%%%%%%%%%%%%%%%%%%%%%%%%%%%%%%%%%%%%%%%%%%%%%%%%%%%%%%%%%%%%%%%%%%%%%%%%%%%%
%%%%%%%%%%%% Blackboard bolds
%%%%%%%%%%%%%%%%%%%%%%%%%%%%%%%%%%%%%%%%%%%%%%%%%%%%%%%%%%%%%%%%%%%%%%%%%%%%%%

\newcommand{\bbA}{{\ensuremath{\mathbb A}} }
\newcommand{\bbB}{{\ensuremath{\mathbb B}} }
\newcommand{\bbC}{{\ensuremath{\mathbb C}} }
\newcommand{\bbD}{{\ensuremath{\mathbb D}} }
\newcommand{\bbE}{{\ensuremath{\mathbb E}} }
\newcommand{\bbF}{{\ensuremath{\mathbb F}} }
\newcommand{\bbG}{{\ensuremath{\mathbb G}} }
\newcommand{\bbH}{{\ensuremath{\mathbb H}} }
\newcommand{\bbI}{{\ensuremath{\mathbb I}} }
\newcommand{\bbJ}{{\ensuremath{\mathbb J}} }
\newcommand{\bbK}{{\ensuremath{\mathbb K}} }
\newcommand{\bbL}{{\ensuremath{\mathbb L}} }
\newcommand{\bbM}{{\ensuremath{\mathbb M}} }
\newcommand{\bbN}{{\ensuremath{\mathbb N}} }
\newcommand{\bbO}{{\ensuremath{\mathbb O}} }
\newcommand{\bbP}{{\ensuremath{\mathbb P}} }
\newcommand{\bbQ}{{\ensuremath{\mathbb Q}} }
\newcommand{\bbR}{{\ensuremath{\mathbb R}} }
\newcommand{\bbS}{{\ensuremath{\mathbb S}} }
\newcommand{\bbT}{{\ensuremath{\mathbb T}} }
\newcommand{\bbU}{{\ensuremath{\mathbb U}} }
\newcommand{\bbV}{{\ensuremath{\mathbb V}} }
\newcommand{\bbW}{{\ensuremath{\mathbb W}} }
\newcommand{\bbX}{{\ensuremath{\mathbb X}} }
\newcommand{\bbY}{{\ensuremath{\mathbb Y}} }
\newcommand{\bbZ}{{\ensuremath{\mathbb Z}} }
%\newcommand{\mathds}{1}  %%%%%%%%%%%%%%%%%%%%%%%%%%%%%%%%%%%%%%%%%%%%%%%%%%%%% x ale: da togliere

%INIZIO COMANDI AGGIUNTI DA ALFIE
\newcommand{\lrgh}{\longleftrightarrow}
\newcommand{\Rgh}{\Rightarrow}
\newcommand{\rgh}{\rightarrow}
%FINE COMANDI AGGIUNTI DA ALFIE

\newcommand{\mamma}{MA$[\b]$}
\newcommand{\mammapoisson}{MA$[\b]_{\rho,\nu}$}
\newcommand{\mammaL}{MA$[\b,\ell]$}
\newcommand{\mammaLpoisson}{MA$[\b,\ell]_{\rho,\nu}$}
\newcommand{\V}{V_{\b,\ell}^\o}

\newcommand{\giggio}{$\cG[\z, \b]$}

\newcommand{\giggioest}{$\cG[\z,\b](\o)$}

%\newcommand{\giggioest}{$\cG[\z,\b](\o)= \bigl( \hat \o, \cE[\z,\b](\o)\bigr)$}

\newcommand{\giggioeststar}{$\cG_*[\z,\b](\o)= \bigl( \hat \o_*, \cE_*[\z,\b](\o)\bigr)$}


\newcommand{\ringo}{(RN)$^{\b}_\ell[\o]$}

\newcommand{\ent}{{\rm Ent} }
\newcommand{\var}{\operatorname{Var}}
\newcommand{\dert}{\frac{{\rm d}}{{\rm d} t}}
\newcommand{\av}{{\rm Av} }
%\newcommand{\letters}{\renewcommand{\theequation}{\Alph{equation}}}
%\newcommand{\numbers}{\renewcommand{\theequation}{\arabic{equation}}}

%
%%%%%%%%%%%%%%%%%%%%%%%%% GRECO
%
\let\a=\alpha \let\b=\beta   \let\d=\delta  \let\e=\varepsilon
\let\f=\varphi \let\g=\gamma \let\h=\eta    \let\k=\kappa  \let\l=\lambda
\let\m=\mu   \let\n=\nu   \let\o=\omega    \let\p=\pi  \let\ph=\varphi
\let\r=\rho  \let\s=\sigma \let\t=\tau   \let\th=\vartheta
\let\y=\upsilon \let\x=\xi \let\z=\zeta
\let\D=\Delta \let\F=\Phi  \let\G=\Gamma  \let\L=\Lambda \let\The=\Theta
\let\O=\Omega \let\P=\Pi   \let\Ps=\Psi \let\Si=\Sigma \let\X=\Xi
\let\Y=\Upsilon%

%


%
\newcommand{\rosso}{\textcolor{black}} %red
\newcommand{\vvv}{\textcolor{black}} %green
\newcommand{\blu}{\textcolor{black}}
\newcommand{\orc}{\textcolor{black}} % cyan
%\newcommand{}{\textcolor{blue}}



\newcommand{\scrM}{ \ensuremath{ \mathscr{M}}}

%%%%%%%%%%%%%%%%%%%%%%%%%%%%%%%%%%
\newcommand{\club}{$\clubsuit$}
\newcommand{\ds}{\displaystyle}
\newcommand{\da}{\downarrow}

\newcommand{\ee}{{\rm e}}


\newcommand{\ddiv}{\text{div}}

\newcommand{\rclub}{\textcolor{red}{$\clubsuit$}}




\newcommand{\be}{\begin{equation}}
\newcommand{\en}{\end{equation}}

\newcommand{\gino}{${\rm PPP}[\rho]$, ${\rm PPP}[\rho,\nu]$}

\newcommand{\cv}{\rosso{\cV_\alpha}}
\newcommand{\cvp}{\rosso{\cV^+_\alpha}}
\newcommand{\vpz}{\vvv{\bbP^{\hspace{0.013cm}0} }}
\newcommand{\hvpz}{\vvv{{\hat \bbP}^{\hspace{0.013cm}0} }}
\newcommand{\bvpz}{\vvv{{\bar\bbP}^{\hspace{0.013cm}0} }}
\newcommand{\tvpz}{\vvv{{\tilde\bbP}^{\hspace{0.013cm}0} }}
\newcommand{\vqz}{\vvv{\bbQ^{\hspace{0.013cm}0} }}


\newcommand{\pn}{\rosso{{\hat\bbP}^{\hspace{0.03cm}n}}}


\author[A.~Faggionato]{Alessandra Faggionato}
\address{Alessandra Faggionato.
  Dipartimento di Matematica, Universit\`a di Roma `La Sapienza'
  P.le Aldo Moro 2, 00185 Roma, Italy}
\email{faggiona@mat.uniroma1.it}

%
%
%   BIBLIOGRAPHY
%
%
%   Kal, Lemma 12.1 pag.225
%
%   \cite[Exercise~6.4.2]{DV}   pag. 195
%
%   \cite{Kal}[Theorem~16.16]  pag. 316
%
%    \cite[Thm.~16.19]{Kal}   pag. 318



%\thanks{We acknowledge the support by  PRIN 20155PAWZB "Large Scale Random Structures"}


\title[Mott's  law]{Mott's law for the  Miller--Abrahams random resistor network and for  Mott's random walk}

\begin{document}

\begin{abstract}
  Mott's  variable range hopping (v.r.h.) is  the phonon-induced hopping of electrons in disordered solids (as doped semiconductors) within the regime of strong Anderson localization. It was introduced by N.~Mott  to explain the anomalous low temperature conductivity decay in dimension $d\geq 2$, corresponding  now to the  so called Mott's law. We provide a rigorous derivation of this Physics law for two effective models of Mott v.r.h.: the Miller-Abrahams random resistor network and Mott's random walk. We also determine the  constant multiplying the power of the inverse temperature in the exponent in Mott's law, which was an open problem also on a  heuristic level.  Our proof uses also our previous results of stochastic homogenization and percolation for the above models. The percolation  results, and therefore also our derivation of Mott's law, are exhaustive      when the energy marks have a fixed sign. For  energy marks with different signs we have derived  the low temperature  behavior of the critical conductance, while the proof of Mott's law  still works if    the expected percolation property of left-right crossings in the supercritical phase   is assumed as an Ansatz.


\smallskip

\noindent {\em Keywords}: marked simple point process, 
Poisson point process,    Mott's random walk,  Miller--Abrahams random resistor network, Mott's law, percolation, stochastic homogenization. 

\smallskip

\noindent{\em AMS 2010 Subject Classification}: 
%60K35,  %Interacting random processes; statistical mechanics type models; percolation theory
60G55, %point processes
82B43, % percolation
%82C05.  %Classical dynamic and nonequilibrium statistical mechanics (general)
%74Q05, %	Homogenization in equilibrium problems
82D30 %	Random media, disordered materials 




\end{abstract}


\maketitle



\section{Introduction} 


Mott's  variable range hopping is a mechanism of phonon--assisted electron transport taking place in amorphous solids (as doped semiconductors)  in the regime of strong Anderson localization \cite{AHL,MD,POF,Sa,SE}. It has been introduced by N.~Mott in order to explain the anomalous 
 decay of the conductivity  at low temperature \cite{Mott,Mott_Nob}.  
 
Calling $x$ the  impurity positions  in the doped semiconductor, the electron Hamiltonian has exponentially localized quantum eigenstates with localization centers $x$ and corresponding energy $E_{x}$  close to the Fermi level, set equal to zero in what follows. At low temperature phonons  induce transitions between the localized eigenstates, the rate of which can be calculated by the Fermi golden rule \cite{MA,SE}. In the simplification of spinless electrons,   the resulting rate for an electron to  hop from $x$ to the unoccupied site $y$ is then given by  (cf. \cite[Eq.~(3.6)]{AHL}) 
\be\label{tunnel} \exp\Big\{-\frac{2}{\g}  |x-y| - \beta \{ E_{x} -E_{y} \}_+\Big\} \,.\en
In \eqref{tunnel} $\g$ is the localization length, $\beta=1/k T $ is the inverse temperature and $\{a\}_+:= \max\{0,a\}$.


The above set \rosso{of impurity locations} can be modeled by a   random simple point process (SPP), marked by  random variables $E_x$ (called energy marks) which are usually taken i.i.d. with some common distribution $\nu$. The physically relevant distributions for inorganic media are of the form  $\nu(dE)= c\, |E|^{\alpha} dE$ with  finite support  $[-A,A]$  for some exponent $\a\geq 0$ \cite{SE}. Mott's law \cite{Mott,Mott_Nob,MD,SE} then  predicts  that, for $d\geq 2$,  the DC conductivity matrix $\s(\b)$ of the medium decays to zero when  $\b \to+ \infty$ as 
\be\label{mottino}
\s(\b) \approx  A(\b)\exp \bigl\{ -\k \, \beta ^\frac{\a+1}{\a+1+d}\bigr\}\,,
 \en
 where the prefactor matrix $A(\b)$ exhibits  a negligible $\b$--dependence (to simplify here the exposition, we restrict just now to isotropic media).
 Strictly speaking, Mott derived  the above asymptotics for $\a=0$, while  Efros and Shklovskii derived its extension. 
 %for $\a=d-1$. The cases $\a=0,d-1$ are the two physically relevant ones. 
  For $d=1$ the DC conductivity presents an Arrenhius--type decay  as $\b \to+ \infty$  for all $\a\geq 0$ \cite{Kur}, i.e.\
 \be\label{mottino1}
 \s(\b)\approx  A(\b)\exp \bigl\{ -\k \b\}\,.
 \en

Due to   localization one can treat the above electron conduction by a hopping process of classical particles (see also \cite{ABS,BRSW}), thus leading  anyway to a complicate simple exclusion process due to the  Pauli blocking. The reversible measure of the exclusion process is  the Fermi-Dirac distribution (i.e. the Bernoulli product probability measure  such that the probability of having a particle at $x$ is proportional to $e^{-\b (E_x-c)}$, $c\in \bbR$ being the chemical potential).
 Effective simplified models   are  given by the Miller--Abrahams (MA) random resistor network \cite{AHL,MA,POF,Sa,SE}   and  by Mott's random walk \cite{FSS}.
The MA random resistor network has nodes $x$ and,   between any  pair of nodes $x\not=y$, it has  an electrical filament of  conductivity 
\be\label{condu}
c_{x,y}:=\exp\Big\{ - \frac{2}{\gamma} |x-y| -\frac{\b}{2} ( |E_x|+ |E_y|+ |E_x-E_y|) \Big\}\,.
\en
Mott's random walk is the continuous--time random walk with \rosso{states given by the impurity locations} and probability rate for a jump from $x$ to $y$ given by \eqref{condu}. We point out that the r.h.s. of \eqref{condu}   corresponds to the leading term for $\b$ large of \eqref{tunnel} multiplied by the probability in the Fermi-Dirac distribution that $x$ and $y$ are, respectively,   occupied and unoccupied  by an electron 
(see \cite[Eq.~(3.7)]{AHL}). A more detailed discussion on the MA resistor network and Mott's random walk  is provided in Sections \ref{ssec_MA} and \ref{sec_mott_rw}.



The original  derivation of the  laws \eqref{mottino} and \eqref{mottino1} is rather heuristic.  More robust    arguments  have been proposed in the Physics  literature (see \cite{AHL,Po,POF,Sa,SE}), in particular for the case $d\geq 2$.  Rigorous bounds in agreement with \eqref{mottino} and \eqref{mottino1} have beed derived in \cite{CF1,Fhom2,FM,FSS} (see Section~\ref{sec_mott_rw}). The characterization of the constant $\k$ was an open problem, even at heuristic level (see \cite[Chapter~5]{POF}). 


\smallskip
Let us describe our main contributions. First  of all we point out that, by means of homogenization theory,  under very general assumptions in \cite{Fhom2}
we proved  that the conductivity matrix $\s(\b)$ of the MA random resistor network equals the asymptotic diffusion matrix $D(\b)$  of Mott's random walk times the intensity $\rho$ of the SSP  (i.e. the mean impurity density). Hence, all results presented below for $\s(\b)$ hold also for $D(\b)$.

 In Section \ref{MR_PPP} we  take $d\geq 2$ and  consider as  marked SPP  a homogeneous Poisson point process (shortly, PPP)  with intensity $\rho$ and i.i.d. energy marks   with common distribution $\nu$ which, in some neighborhood of the origin,
 \rosso{has density proportional to $|E|^\a$ or $E^\a \mathds{1}(E\geq 0)$, where $\a\geq 0$. In particular, for some $C_0>0$, near the origin $\nu$}
    has   the form $(\a+1) C_0^{-\a-1}2^{-1} |E|^\a dE$ (case $\nu\in \cv$)  or $(\a+1)C_0^{-\a-1} E^\a \mathds{1}(E\geq 0) dE$  (case $\nu\in \cvp$).
We first consider 
  the critical conductance $c_c(\beta, \rho,\nu)$. According to \cite{AHL}  $c_c(\beta, \rho,\nu)$ is such that the graph obtained from the MA resistor network by keeping only the filaments  $\{x,y\}$ 
with conductivity $c_{x,y}>C$ a.s. does not percolate if $C>c_c(\beta, \rho,\nu)$, while it  a.s. percolates if $c_c(\beta, \rho,\nu)>C$.  In Theorem \ref{teo1} and its Corollary \ref{luna} we show that the limit 
\[ \chi(\rho, \nu):=\lim _{\b\to +\infty} \b^{- \frac{\a+1}{\a+1+d} } \ln c_c(\b,\rho,\nu )\,\] exists and compute it. 
We then prove Mott's law \eqref{mottino} for $\nu\in \cvp$ (cf. Theorem \ref{teo1} and  Corollary  \ref{immenso}), showing the the constant $\k$ in \eqref{mottino} indeed equals $- \chi(\rho, \nu)$. It turns out that 
\be\label{eccolino}
-\k=\chi(\rho,\nu) =\bigl( {\l_c^+(\a)   C_0^{\a+1}/\r}\bigr)^{\frac{1}{\a+1+d}}\,,
\en
where $\l_c^+(\a)\in (0,+\infty) $ is a constant of percolation type  depending  only from $\a$ and $d$ (although the dependence on $d$ does not appear in the notation). It has the following characterization.  Consider the MA resistor network on a PPP with density $\l$ and  i.i.d. energy  marks   in $[0,1]$ having  distribution $\nu_{1,\a}^+$  with  density proportional to $E^\a$ on $[0,1]$. Afterwards keep only filaments with conductivity at least $e^{-1}$ at inverse temperature $\b=1$. Then  the resulting network a.s. percolates for $\l>\l_c^+(\a)$ and a.s. does not percolate for $\l<\l_c^+(\a)$.


\smallskip

Our proof of the above  Mott's law relies also  on our previous results in  percolation \cite{FagMim1,FagMim2} (obtained in collaboration with H.A.~Mimun) and stochastic  homogenization \cite{Fhom2}. Having expressed $\s(\b)$   in terms of the effective homogenized matrix in \cite{Fhom2}, we can use the variational characterization of the latter to obtain upper bounds on $\s(\b)$ via suitable test functions. On the other hand, being $\s(\b)$ the limit of the suitably rescaled conductivity of finite volume resistor networks as described in Section \ref{ssec_MA}, we can use Rayleigh's monotonicity law to lower bound 
$\s(\b)$ by thinning the finite volume resistor networks  when  having a good knowledge about the statistics of their  vertex-disjoint left-right (LR) crossings. The procedure is anyway different (and much refined) from the one in \cite{FM,FSS} and  now our upper and lower bounds on $\s(\b)$ asymptotically match as $\b\to+\infty$ (thanks also to suitable rescaling arguments). The test functions and thinned resistor networks used in \cite{FM,FSS}, respectively,  are of different nature, not optimal, while the ones presented here are inspired by the paradigm of Critical Path Analysis \cite{AHL} and require a detailed knowledge of a suitable percolation model investigated separately in   \cite{FagMim1,FagMim2}.
The overall derivation of Mott's law  has required (and motivated)   the  results derived in \cite{Fhom2,FagMim1,FagMim2} (although we treated there larger classes of models).


The analysis of LR  crossings in the supercritical phase provided  in \cite{FagMim2} uses the FKG inequality which is satisfied when the energy marks have a fixed sign (case $\nu \in \cvp$) but is violated otherwise (case $\nu\in \cv$), as discussed in Appendix \ref{app_FKG}. Once  the results of \cite{FagMim2} are extended  to the MA resistor network with energy marks of arbitrary signs, then our analysis here provides Mott's law \eqref{mottino} also for the case $\nu\in \cv$. Hence, in Corollary \ref{immenso} we obtain Mott's law for $\nu\in \cv$ under a suitable Ansatz on the LR crossings. In this case 
\eqref{eccolino} is valid with the constant $\l_c^+(\a)$ replaced by $\l_c(\a)$. The definition of $\l_c(\a)$ is the same of  $\l_c^+(\a)$  with exception that the energy marks have  value in $[-1,1]$  with  distribution $\nu_{1,\a}$ having  density proportional to $|E|^\a$ on $[-1,1]$. \rosso{The above Ansatz is used to derive a lower bound in \eqref{mottino} for $\nu\in \cV$, while the upper bound is proved without any Ansatz (see Theorems \ref{parataUB}  and \ref{parataLB}).} We point out that when  the intensity of the PPP is large enough  the above Ansatz  can be proved  by stochastic domination with the Poisson Boolean model (see Appendix \ref{guizzo} for this domination).We stress that also the spinless modeling common in the Physics literature is a simplification of the realistic model with spins.

\smallskip

Another main contribution  is presented in Section \ref{MR_univ} and comes from  the investigation of the universality in Mott's law as Physics law.  The derivation of the results   \rosso{commented above} holds in good part for generic marked SPPs, and not only for the ones obtained by marking a PPP by i.i.d. random variables. As the reader will see,  a central role is played by a  mechanism of thinning and rescaling which ultimately  captures the backbone of  filaments in the MA resistor network mainly contributing to the electron transport \rosso{(see the discussion at the end of  Section \ref{MR_univ} where an overview of the main strategy is provided)}. In particular, at inverse temperature $\b$, just  points $x$  with marks  $E_x$ in a suitable window $[ -E(\b), E(\b)]$ give the main contribution to transport.
In Theorem \ref{paradiso} we  show the following. Let us  start with a generic 
ergodic stationary SPP with i.i.d. energy marks having distribution $\nu$  which asymptotically behaves as $\nu^+_{1,\a}$ or $\nu_{1,\a}$ while approaching  the origin. Then \rosso{the rescaled set of marked points $\left( x/(\b E(\b)), E_x /E(\b)\right)$  with $|E_x|\leq  E(\b)$} 
%the above set of marked points $x$ satisfying  $|E_x|\leq  E(\b)$  converges (after space and energy rescaling) 
converges to a  PPP with some intensity  $\l$ and with  i.i.d. energy marks 
having distribution respectively  $\nu^+_{1,\a}$ and $\nu_{1,\a}$ (the same random geometric structure entering in the definition of $\l_c^+(\a)$ and $\l_c(\a)$).
% backbone of  filaments, under suitable rescaling, will converge to the random graph entering in the definition of the above constant $\l_c^+(\a)$ and which is built on a PPP with i.i.d. marks. 
Hence, also when starting from a  non Poissonian  distribution of impurities and from a more general energy distribution $\nu$, at low temperature   the above marked  PPP  emerges and the contribution to transport of the  details of the original distribution progressively  vanishes.


\smallskip


As next steps we aim to derive the scaling limit  \eqref{mottino1} for $d=1$ and prove  the above Ansatz on LR crossings for energy marks of different signs. These problems are the object  of our present ongoing work.



\smallskip



{\bf Outline of the paper}.
In Section \ref{ssec_marked_SPP} we recall some basic notions of marked simple point processes on $\bbR^d$. In Section \ref{palmato} we introduce the Palm distribution 
and the effective homogenized matrix. In Section \ref{ssec_MA} we describe the MA  random resistor network  and recall the scaling limit of its directional conductivity obtained in \cite{Fhom2}. In Section \ref{sec_mott_rw} we describe  Mott's random walk and recall some previous results  from \cite{CF1,Fhom1,Fhom2,FM,FSS}. In Section \ref{grafene} we  focus on the part of the MA resistor network with conductances lower bounded by a given threshold and introduce the critical conductance threshold. Sections \ref{MR_PPP} and \ref{MR_univ} contains our main results. 
%
%In the former we  focus of marked SPPs given by $\nu$--randomization of a PPP (i.e.  given by ${\rm PPP}[\rho,\nu]$)
%and  prove precise low temperature  scaling limits  for the critical conductance and  the infinite volume conductivity of the
%MA resistor network as well for the asymptotic diffusivity of Mott's random walk. 
%In the latter we focus on a mechanism of thinning and rescaling which  captures the backbone of  filaments in the MA resistor network mainly contributing to the electron transport and show how, starting with a generic stationary ergodic marked SPP $\hat\bbP_\nu$, it  leads to  ${\rm PPP}[\rho,\nu^+_{1,\a}]$ or ${\rm PPP}[\rho,\nu_{1,\a}]$ for suitable probabilities $\nu^+_{1,\a}$ and $\nu_{1,\a}$.
The other sections are devoted to the proofs, with exception of Appendix \ref{app_FKG} where we show the violation of FKG inequality when the energy marks are not of fixed sign.

We point out that several intermediate results are valid for all dimensions $d\geq 1$. We will state explicitly when we will restrict to $d\geq 2$.
%
%%%%%%%%%%%%%%%%%%%
\section{Some preliminaries on marked simple point processes on $\bbR^d$}\label{ssec_marked_SPP} 
In what follows, given a topological space $\cX$, we denote by $\cB(\cX)$ its Borel $\s$--algebra. 
%\rosso{We  work with the Euclidean space $\bbR^d$, $d\geq 2$.}
 Given $A\in \cB(\bbR^d)$ we denote by $\ell(A)$ the Lebesgue measure of $A$. 
Given a measure $\mathfrak{m}$ and a function $f$ defined on the same space, we write $\mathfrak{m}[f]$ for the integral of $f$ w.r.t.~$\mathfrak{m}$.






 
 We denote by $\hat \O$ the space of  locally finite subsets of $\bbR^d$. Trivially, elements of $\hat \O$ are countable sets. Moreover, each $\hat \o \in \hat \O$ can be identified with the atomic measure $\sum_{x\in \hat \o }\d_x$, which  gives finite mass to bounded Borel subsets of $\bbR^d$. This allows to endow $\hat \O$ with the standard metric defined on the space of boundedly finite measures on $\bbR^d$ (cf. \cite[Eq.~(A2.6.1) in Appendix~A2.6]{DV}). The precise definition of this metric is not necessary for the rest and  we do not recall it here, but we will recall some of its properties when needed. We just mention here that a sequence $(\hat\o_n)_{n\geq 1}$  converges to $\hat \o$ in $\hat \O$ if and only if $\lim_{n\to +\infty}\hat\o_n[f]= \hat \o[f]$ for all   continuous  functions $f: \bbR^d\to \bbR$ with bounded  support (see  \cite{DV}[Thm.~A2.6.II]). 
 
 Due to  the above interpretation of $\hat \o$ as atomic measure, 
  given  $A\subset \bbR^d$,  $\hat \o (A)$ equals  the cardinality of $\hat \o \cap A$.
  One can prove that the Borel $\s$--algebra $\cB(\hat\O)$ is then generated by the sets $  \{\hat \o \in \hat \O\,:\, \hat \o (A) =n\} $ with $A\in \cB(\bbR^d)$ and $n \in \bbN$ (combine Prop.~7.1.III and  Cor.~7.1.VI in \cite{DV}).  From now on, $\hat \O$ will be thought of  as a measure set with $\s$--algebra of  measurable sets given by $\cB(\hat \O)$. 
 
 
   Given $x\in\bbR^d$ and a  configuration $ \hat \o\in\hat \O$, we define the translated configuration $\t_x\hat \o$  as
\be\label{traslare}
\t_x\hat \o:= \{ x_i-x \}\; \text{ if }\; \hat \o=\{ x_i\}\,.
\en
\rosso{(above, and in the rest, we write $\{x_i\}$ for $\{x_i:i\in I\}$, omitting the index set).}
Each translation $\t_x:\hat\O\to \hat\O$ is a measurable map.
 A set $A\in \cB(\hat\O)$ is called \emph{translation invariant} if $A=\t_x A$ for all $x\in \bbR^d$.
 
 \begin{Definition}[SPP] \label{def_spp}
 A \emph{simple point process} (shortly, SPP) $\hat \xi$ on $\bbR^d$ is a random element of $\hat \O$. 
 $\hat \xi$ is called \emph{stationary} if its law $\hat \bbP$ satisfies $\hat \bbP( \t_x A)=\hat \bbP(A)$ for all $A\in \cB(\hat \O)$ and $x\in \bbR^d$. In this case, the \emph{intensity} of $\hat \xi$ is defined as $\rho:=\hat \bbE\bigl[ \hat\o([0,1]^d) \bigr]\in [0,+\infty]$, $\hat \bbE$ being the expectation w.r.t. $\hat \bbP$.
$\hat \xi$  is called \emph{ergodic} if it is stationary  and  $\hat \bbP(A)\in \{0,1\}$ for any translation invariant set $A\in \cB(\hat \O)$. 
\end{Definition}
For a stationary SPP with intensity $\rho$ it holds  $\hat \bbE\bigl[ \hat\o(A) \bigr]= \rho \ell(A)$ for any Borel set $A\in \cB (\bbR^d)$ \cite{DV}.
 
 
\medskip
In order to introduce the marked simple point processes,  
 we denote by $\O$ the space of subsets $\o=\{ (x_i, E_{x_i}) \}\subset\bbR^d\times \bbR$ such 
 that $\{x_i\}$ is a (countable) locally finite subset of $\bbR^d$ (i.e.~$\{x_i\}\in \hat \O$).
% \begin{itemize}
% \item[(a)] $\{x_i\}$ is a (countable) locally finite subset of $\bbR^d$ (i.e.~$\{x_i\}\in \hat \O$);
% \item[(b)] for any two points $(x_i, E_{x_i})$, $(x_j , E_j)$ in $\o$ it holds $x_i\not =x_j$.
% \end{itemize}
  The real number $E_{x_i}$ is called the \emph{energy mark} of the site $x_i$.
  % (by property (b) above the energy mark is unique). %Given $\o\in \O$ we denote by $\hat \o$ the spatial support of $\o$, i.e.
% \[
% \hat \o=\{x_i\} \;\text{ if } \; \o = \{(x_i, E_{x_i)}\}\,.
% \]
%The set $\o$ is usually identified with the atomic measure $\sum_{x\in \hat \o} \d_{(x,E_x)}$ (which we still denote by $\o$). In particular,
We will often  identify $\o=\{ (x_i, E_{x_i}) \} \in \O$ with the atomic measure $\sum_i \d_{(x_i,E_{x_i})}$. Note that, given  $A\subset \bbR^d\times \bbR$,   $\o(A)$  is the number of points of $\o$ in $A$.


If we denote by $\hat \O_{d+1}$ the metric space defined as $\hat \O$ above but with $\bbR^d$
replaced by $\bbR^{d+1}=\bbR^d\times\bbR$, one can prove that $\O$ is a Borel subset of $\hat \O_{d+1}$\footnote{$\O=\cap_{n=1}^\infty \cup_{m=1}^\infty B_{n,m}$ where $B_{n,m}$ is the set of 
 $\o \in \hat \O_{d+1}$ such   that $\o$ contains at most one point in the sets $(z,0) + [-1/m,1/m)^{d-1}\times \bbR$ for all $z\in [-n,n)^{d-1}\cap (\bbZ^{d-1}/m)$.}.
We consider $\O$ with the metric and topology induced from $\hat \O_{d+1}$. We  then get that  the Borel $\s$--field  $\cB(\O)$  is generated by the sets $\{ \o\in \O\,:\, \o(A)=n\}$, where    $A$ and $n$ vary respectively in $\cB( \bbR^d\times \bbR) $  and $\bbN$. Moreover,  a  sequence $(\o_n)_{n\geq 1}$ in $\O$ converges to $ \o\in \O$ if and only if $\lim_{n\to +\infty} \o_n[f]= \o[f]$ for all   continuous  functions $f: \bbR^d\times \bbR \to \bbR$ with bounded support.

%Again, by identifying $\o=\{ (x_i, E_{x_i}) \} \in \O$ with the atomic measure $\sum_i \d_{(x_i,E_{x_i})}$ (which is boundedly finite on $\bbR^{d+1})$, 
%the space $\O$
%  is naturally endowed with a metric  for which the Borel $\s$--field  $\cB(\O)$  is generated by the sets $\{ \o\in \O\,:\, \o(A)=n\}$, where    $A$ and $n$ vary respectively in $\cB( \bbR^d\times \bbR) $  and $\bbN$ \rosso{\cite{DV}}. Moreover, a  sequence $(\o_n)_{n\geq 1}$ converges to $ \o$ in $\O$ if and only if $\lim_{n\to +\infty} \o_n[f]= \o[f]$ for all   continuous  functions $f: \bbR^d\times \bbR \to \bbR$ with bounded support \rosso{\cite{DV}}. 
%  

    Given $x\in\bbR^d$ and  $ \o\in \O$, we define the translated configuration $\t_x\o$  as
\be\label{traslare_bis}
\t_x\o:= \{ (x_i-x, E_{x_i}) \}\; \text{ if }\; \o=\{ (x_i,E_{x_i})\}\,.
\en
 Finally, 
 we introduce the spatial projection
\be
\O \ni \o \mapsto \widehat{\o} \in \hat \O \text{ where } \widehat{\o}:=\{x_i\} \text{ if }
\o = \{(x_i, E_{x_i})\}\,.
\en
 In particular, we will frequently  write $\o=\{ (x, E_x)\,:\, x\in \widehat \o\}$. Both the translation  $\t_x:\O\to \O$ and the projection $\O\ni \o \mapsto \widehat{\o} \in \hat \O$ are measurable functions, as can be easily checked.
 %\footnote{The map $\t_x$ is the restriction to $\O$ of the translation along $(x,0)$ in $\hat\O_{d+1}$, which is measurable. The measurability of the projection follows from the trivial observation that $\{\o\in \O: \widehat{\o}(A)=n\}=\{\o \in \O: \o( A\times \bbR)=n\}$ for all $A\in \cB(\bbR^d)$.  }.
 

% \rosso{E' misurabile..? continua?.. }
\begin{Definition}[Marked SPP] \label{def_mspp}
A \emph{marked simple point process} (shortly, marked SPP) $\xi$  on $\bbR^d$ is a random element of  $\O$.  $\xi$ is called \emph{stationary} if its law $\bbP$ satisfies $\bbP( \t_x A)= \bbP(A)$ for all $A\in \cB(\O)$ and $x\in\bbR^d$. In this case, the intensity of $\xi$ is defined as $\rho:=\bbE\bigl[ \widehat\o([0,1]^d) \bigr]\in [0,+\infty]$, $\bbE$ being the expectation w.r.t. $\bbP$.
$\xi$  is called \emph{ergodic} if it is stationary  and if $\bbP(A)\in \{0,1\}$ for any translation invariant set $A\in \cB(\O)$ \rosso{(i.e.~such that $\t_x A=A$ for all $x\in \bbR^d$).} 
 \end{Definition}

\begin{Warning} In general, the symbol $\; \hat{}$ will refer to SPPs. Space of configurations, generic configurations, SPPs, laws of SPPs will 
usually be denoted respectively by $\hat\O$, $\hat \o$, $\hat \xi$, $\hat \bbP$.
Note that the spatial projection of $\o\in \O$ is denoted by $\widehat{\o}$, while a generic element of $\hat \O$ is denoted by $\hat{\o}$. The context will also allow  to easily distinguish between   $\widehat{\o}$ and $\hat \o$.
\end{Warning}


The marked SPPs we are interested to will be obtained  as
  $\nu$--randomization of a SPP \cite{Kal}. Let us recall  this procedure.



\begin{Definition}[$\nu$-randomization $\hat \bbP_\nu$] \label{def_randomization}
Consider   a SPP $\hat \xi$ on $\bbR^d$ with law  $\hat\bbP$ on $(\hat\O, \cB(\hat\O))$. Given a realization of $\hat \xi$, mark each point by i.i.d. random variables  (independently from all the rest) with common distribution $\nu$, and call $\xi$ the resulting marked SPP. 
We say that the law $\hat \bbP_\nu$ of 
  $\xi$ is the  $\nu$-randomization of $\hat \bbP$ and that   $\xi$ is the $\nu$--randomization of $\hat \xi$.  
 \end{Definition}
 
 We now recall the notion of $p$--thinning, which will be crucial in what follows:
 \begin{Definition}[$p$--thinning $\hat \bbP_p$] \label{def_thinning} Given a SPP  $\hat \xi$  on $\bbR^d$ with law $\hat \bbP$ and given $p\in [0,1]$, the $p$--thinning of $\hat \xi$ is the SPP  obtained  as follows. Given a realization $\hat\o$ of $\hat \xi$, for each point in $\hat \o$ one tosses  a biased coin with probability $p$ to give head: if the coin gives head the point is kept, otherwise it is deleted. We denote by $\hat\bbP_p$  the law of the $p$-thinning of $\hat\xi$. We will shortly say that $\hat\bbP_p$ is  the $p$-thinning of $\hat \bbP$.
 \end{Definition}


\subsection{Poisson point processes}
A special role in Mott's variable range hopping is played by Poisson point processes.

\begin{Definition} [\gino]
\label{def_PPP} We denote by ${\rm PPP}[\rho]$ the homogeneous Poisson point process  on $\bbR^d$ with intensity  $\rho$ and we denote by  ${\rm PPP}[\rho,\nu]$ the 
 $\nu$--randomization of ${\rm PPP}[\rho]$ (having automatically intensity $\rho$). We use the same symbols also for the associated laws, i.e.~${\rm PPP}[\rho](A)$ is the probability of the event $A\in \cB (\hat \O)$ for the homogeneous Poisson point process with intensity $\rho$, and similarly for  ${\rm PPP}[\rho,\nu](A)$ with $A\in \cB(\O)$.
 \end{Definition}
 %kuka, controllare che queste proprieta' siano derivate
  We  recall that a SPP $\hat \xi$ on $\bbR^d$ is a {\rm PPP}$[\rho]$ whenever
%\begin{itemize}
%\item[(i)]
  (i) for mutually disjoint Borel subsets $A_1,A_2, \dots, A_k$ in $ \bbR^d$, the random variables $\hat \xi(A_1),\hat \xi(A_2),\dots, \hat \xi(A_k)$ are independent;
  % \item[(ii)]
  (ii) for any bounded   set  $A\in \cB(  \bbR^d)$, $\hat \xi(A)$ is a Poisson random variable with parameter $\r\, \ell(A)$.
%\[ P\bigl(\hat \xi (A)= k\bigr)= \frac{e^{-\rho \ell(A)}  (\rho \ell(A) ) ^k }{ k!} \qquad k=0,1,2,\dots\,. \]
In particular, it holds  $ E[ \xi(A)] =\hat\bbE[\hat \o(A)]=\rho \ell(A)$ for any $A \in \cB(\bbR^d)$.
% \end{itemize}
 We  recall that {\rm PPP}$[\rho]$ and {\rm PPP}$[\rho,\nu]$ are stationary and ergodic \rosso{(see e.g.~\cite[Prop.~8.13]{LP}, whose proof can be generalized to the marked case)}.

The families of PPP's and marked PPP's are closed w.r.t. some fundamental operations, which play a central role in Mott's law:
\begin{Proposition}\label{cotechino}
The following holds:
\begin{itemize}
\item[(i)] The $p$--thinning of 
  ${\rm PPP}[\rho]$ is  given by 
${\rm PPP}[p \rho]$. 
\item[(ii)] By sampling $\hat\o\in \hat\O$ according to  ${\rm PPP}[\rho]$, the SPP obtained by rescaling each $x\in \hat \o $ into $x/\ell$ for some fixed $\ell>0$ has law ${\rm PPP}[\ell^d \rho]$.
\item[(iii)] Fixed $\g>0$, by sampling $\o\in \O$ according to  ${\rm PPP}[\rho,\nu]$,    the marked SPP $\o_\gamma:=\{(x,E_x)\,:\, |E_x|\leq \g \}$ has  law ${\rm PPP}[ p  \rho,\nu_\g]$, where $p:=\nu([-\g, \g])$ and $\nu_\g$ is the probability measure $\nu$ conditioned to the event $[-\g,\g]$.
    \end{itemize}
\end{Proposition}
\begin{proof} Properties (i) and (ii) are well known and follow easily from  the above characterization of ${\rm PPP}[\rho]$.
% kuka  \rosso{\club Mettere una referenza per ergodicita' e stazionarieta' del PPP}.
 Property (iii) follows from Lemma \ref{lemma_2021} in Section \ref{Tarvisio} and from Item (i).
\end{proof}
%%%%%%%%%%%%%%%%%%%%%%%%%%%%%%%%%%%%%%%%%
%
%
%
%
%%%%%%%%%%%%%%%%%%%%%%%%%%%%%%%%%%%%%%%%%
\section{Palm distribution and effective homogenized matrix $D(\b)$}\label{palmato}
In this section we introduce the effective homogenized matrix $D(\b)$. As discussed later,  it is  intimately  related to the conductivity properties of the MA random resistor network and to the diffusion properties of Mott's random walk. To define $D(\b)$ we  will  first need to introduce the Palm distribution.


In what follows, $\b$ is a positive number which, in the applications,  will correspond to the inverse temperature.
\begin{Definition}[Conductance field]
Given $\o\in \O$ we associate to each  \rosso{unordered} pair $\{x,y\}$ with $x,y\in \widehat{\o}$ and $x\not=y$ the conducivity 
\begin{equation}\label{conda}
c_{x,y}(\o,\b):=\exp\bigl\{ -  |x-y| -\b  ( |E_x|+ |E_y|+ |E_x-E_y|) \}\,.
\end{equation}
\end{Definition}
\rosso{Note that \eqref{conda} corresponds to \eqref{condu} by taking $\g:=2$ and by replacing $\b/2$ in \eqref{condu}  with $\b$. Trivially, by suitably rescaling distance and energy, the results obtained for \eqref{conda} can be extended to \eqref{condu}.}


Trivially, we have 
\be\label{balzetto}
|E_x|+ |E_y|+ |E_x-E_y| = \begin{cases} 
2 \max \big\{ |E_x|,|E_y|\big\} & \text{ if } E_x \cdot E_y \geq 0 \,,\\
2 |E_x-E_y| & \text{ if } E_x \cdot E_y \leq 0\,.
\end{cases}
\en

To simplify the notation below, it is convenient to set 
\[
%\be\label{bowl89}
c_{x,x}(\o,\b):=0 \qquad \forall \o \in \O\,,\; x\in \widehat{\o}\,.
%\en
\]



\begin{Definition}[Palm distribution]  \label{def_palm} Let $\bbP$ be the law of a stationary  marked SPP on $\bbR^d$ with finite and positive  intensity $\rho$. The Palm distribution $\vpz$ associated to $\bbP$  is defined  as  the unique probability measure on  $(\O,\cB(\O))$  concentrated on the Borel set $\O_0:=\{\o\in \O\,:\, 0 \in \widehat \o\}$ and given by 
\begin{equation}\label{marlena}
\vpz(A)= \frac{1}{m} \int _{\O}d \bbP(\o)\int_{[0,1]^d}  d\widehat{\o} (x) \mathds{1}_A(\t_x \o)\,, \qquad \forall A \in \cB( \O) \,.
 \end{equation}
\end{Definition}
\rosso{Above, and in the rest, we denote by $\mathds{1}_A$ or $\mathds{1}(A)$ (according to better readability) the characteristic function of the event $A$.}

Equation \eqref{marlena} is a special case of Campbell's formula. A more general  Campbell's formula is the following (cf.~\cite[Eq.~(12.2.4)]{DV}, \cite[Thm.~1.2.8]{FKAS}): for any  nonnegative Borel function $f: \bbR^d\times \O \to[0,\infty) $   and for any function $f\in L^1( \bbR^d\times \O, dx \times \vpz)$ it holds 
 \begin{equation}\label{campanello}
 \int_{\bbR^d}dx  \int _{\O_0} d\vpz ( \o) f(x, \o) =\frac{1}{m} \int _{\O}d \bbP(\o)\sum_{x  \in \widehat{\o}} f(x, \t_x \o) \,.
 \end{equation}
 
As discussed  in \cite[Chapter~12]{DV}, roughly $\vpz=\bbP(\cdot\,|\, 0\in \widehat \o)$. 


When $\bbP=${\rm PPP}$[\rho,\nu]$, $\vpz$ can be realized as follows. Sample $\o\in \O$ with  law $\bbP$ and sample, independently from $\o$, the number    $E\in \bbR $ with law $\nu$. Then 
the configuration $\o \cup \{ (0,E)\}$ has law $\vpz$  (see \cite[Chapter~12]{DV}, in addition note that $\bbP( 0\in \widehat{\o})=0$, hence the probability that  in this construction  the origin appears with two different marks is zero). 

If $\bbP$ is the $\nu$--randomization of $\hat \bbP$ (i.e.~$\bbP=\hat\bbP_\nu$) and $\hat \bbP$ is the law of a stationary SPP with finite and positive  intensity $\rho$, then it is simple to check 
% kuka, controllare la referenza
 that the Palm distribution $\vpz$ is the $\nu$--randomization of the Palm distribution $\hvpz$ associated to $\hat\bbP$ (i.e.~$\vpz=(\hvpz)_\nu$). \rosso{Indeed we recall that,}  similarly to \eqref{marlena}, $\hvpz$ is defined as the 
unique probability measure on  $(\hat\O,\cB(\hat \O))$  concentrated on the Borel set $\hat \O_0:=\{\hat \o\in \hat \O\,:\, 0 \in \hat \o\}$ and given by 
\begin{equation}\label{marlenabis}
\hvpz(A)= \frac{1}{m} \int _{\hat\O}d \hat \bbP(\hat \o)\int_{[0,1]^d}  d\hat{\o} (x) \mathds{1}_A(\t_x \hat \o)\,, \qquad \forall A \in \cB( \hat \O) \,.
 \end{equation}





\begin{Definition}[Effective homogenized  matrix] Let $\bbP$ be the law of a stationary  marked SPP on $\bbR^d$ with finite and positive  intensity $\rho$ and let $\vpz$ be the associated Palm distribution. Suppose that 
\be\label{2nd_moment}
\int_{\O_0} d\vpz(\o)\sum_{x\in \widehat \o}  c_{0,x}(\o,\b)|x|^2 <+\infty\,.
\en
Then
the \emph{effective homogenized  matrix}  $D(\b)$ is  
the unique $d\times d$ nonnegative symmetric matrix such that
 \begin{equation}\label{def_D}
 a \cdot D(\b) a =\inf _{ f\in L^\infty(\bbP_0) } \frac{1}{2}\int_{\O_0} d\vpz(\o)\sum_{x\in \widehat \o}  c_{0,x}(\o,\b) \left
 (a\cdot x - \nabla_x f (\o) 
\right)^2\,,
 \end{equation}
 where $\nabla_x f (\o) := f(\t_x \o) - f(\o)$.
\end{Definition}

Note that condition \eqref{2nd_moment} guarantees that the r.h.s. of \eqref{def_D} is finite (take $f\equiv 0$ there).



%%%%%%%%%%%%%%%%%%%%%%%%%%%%%%%%%%%%
%
%
%     MAIN ASSUMPTIONS
%
%
%%%%%%%%%%%%%%%%%%%%%%%%%%%%%%%%%%%%%
%
%\section{Main assumptions}
%
%\club \club aggiornare\\
%
%$\bbP$ stationary marked SPP with finite and positive intensity.
%\be\label{non_periodico}\bbP ( \o \in \O:  \t_x\o\not = \t_{x'} \o  \; \; \forall x\not =x' \text{ in }\bbR )=1\,.
%\en
%As proved in \rosso{\cite{FSS}} the above  property  implies that 
%\be\label{non_periodico_bis}
%\bbP_0 ( \o \in \O:  \t_x\o\not = \t_{x'} \o  \; \; \forall x\not =x' \text{ in } \widehat{\o})=1\,.
%\en
%
%
%
%
%We can state our main  assumptions  on the point process:
%%
%%\begin{Assumptions}
%%
%%\end{Assumptions}
%$\bbP$ is a marked SPP  given by the by $\nu$-randomization of a stationary SPP $\hat\bbP$ on $\bbR^d$ with finite and positive intensity $\rho$. We assume that $\bbP$ is ergodic and that \emph{non periodicity}
% (these last two properties  are automatic is $\nu$ is not degerate)
%Mettere moment conditions...





%%%%%%%%%%%%%%%%%%%%%%%%%%%%%%%%%%%%
%
%
%
%
%
%%%%%%%%%%%%%%%%%%%%%%%%%%%%%%%%%%%%%
\section{Miller--Abrahams random  resistor network}\label{ssec_MA}
In this section we recall some concepts and results from \cite{Fhom2} for what concerns the MA resistor network (we point out that  \cite{Fhom2} covers also many other resistor networks and uses idea of homogenization theory).

\smallskip

 In what follows $\b$ will denote the inverse temperature. Recall \eqref{conda}.
\begin{Definition}[Resistor network \mamma$(\o)$] \label{def_MA}
Given  $\b>0$ and   $\o \in \O$, the  Miller--Abrahams   resistor network MA$[\b](\o)$  has node set $\widehat \o$ and has an electric  filament between any distinct nodes $x\not =y $ in $\widehat \o$ with conductivity $c_{x,y}(\o,\b)$.
\end{Definition}

%Note that, to simplify the notation and without any loss of generality,  \eqref{conda} is obtained from \eqref{condu} by  taking  the localization length $\g$ equal to  $2$ and  by replacing   $\beta/2$ with  $\beta$ (keeping the name of inverse temperature).

%%%%%%%%%%%%%%%%%%%%%%%%%%%%%%%%%%%%%%%%%
%
%
%
%
%%%%%%%%%%%%%%%%%%%%%%%%%%%%%%%%%%%%%%%%%



\subsection{Finite volume  $\ell$--conductivity} \label{fumata}
Given $\ell>0$, we consider the stripe\footnote{We keep a 2-dimensional terminology, although $S_\ell$ lives in $\bbR^d$.}   $S_\ell:=\bbR\times (-\ell/2,\ell /2)^{d-1}$ and we decompose it as 
\be\label{colle1}
S_\ell=S_\ell^-\cup \L_\ell \cup S_\ell^+\,,\en
where 
\[S_\ell^-:=  \{x\in S_\ell\,:\, x_1\leq -\ell/2\}\,,\;  S_\ell^+:=  \{x\in S_\ell\,:\, x_1\geq \ell/2\}\,, \;
 \L_\ell:=(-\ell/2,\ell/2)^d\,.\]



\begin{figure}
\includegraphics[scale=0.20]{FISSO_mott.pdf}
%\includegraphics[scale=0.4]{new-fig2.pdf}
\caption{A portion of the resistor network \mammaL. The box and the stripe correspond to  $\L_\ell$ and   $S_\ell$, respectively. }\label{messicano1}
\end{figure}


For the rest, we assume that  for $\ell$ large enough  the configuration $\o$ satisfies the following property: $\widehat \o$  has nonempty intersection with 
the sets $S_\ell^-$,  $S_\ell^+$  and   $\L_\ell$.
We also assume that $\sum_{y\in \widehat\o }c_{x,y}(\o)<+\infty$ for all $x\in \widehat \o$.
If $\o$ is sampled with law $\bbP$ which is an ergodic stationary marked SPP with finite \rosso{and positive} intensity \rosso{such that 
\be\label{0_moment}
\int d\vpz(\o)\sum_{x\in \widehat \o}  c_{0,x}(\o,\b) <+\infty
\en
holds}  (as in our applications), the above properties hold $\bbP$--a.s. \rosso{(the first property follows by applying the ergodic theorem as in \cite[Prop.~4.2]{Fhom2}, while the boundedness of  $\sum_{y\in \widehat\o }c_{x,y}(\o)$ follows from \cite[Lemma~4.1]{Fhom2}). In the definition  of electrical potential and conductivity given below, we will  understand that the above properties  are  satisfied (in particular,  $\ell$ is taken large enough).}

 \begin{Remark}\label{chitarrina} \rosso{As proved in \cite[Section~5.4]{Fhom1} conditions \eqref{2nd_moment} and \eqref{0_moment} are satisfied whenever $ \bbE[ \widehat{\o} ([0,1]^d)^2 ]  <+\infty$.} 
\end{Remark}

\begin{Definition}[Resistor network \mammaL$(\o)$] \label{def_MAelle}
Given  $\b,\ell>0$ and  $\o \in \O$,
 the Miller--Abrahams  resistor network 
 \mammaL$(\o)$ on the box $\L_\ell$ has node set $ \widehat \o \cap S_\ell$  and   is obtained by assigning to    each  pair  $\{x,y\}$ with   $x \in  \widehat \o \cap \L_\ell$ and $y \in \widehat \o \cap S_\ell$ 
 an electrical filament of conductivity $c_{x,y}(\o,\b)$. 
% If $\o$ is sampled according to a {\rm PPP}$(\rho,\nu)$, we will write \mammaLpoisson$(\o)$  when we need to   point out the sampling procedure.
\end{Definition}
We refer to  Figure \ref{messicano1} for a portion of \mammaL. 


\smallskip

%
%It is simple to check   (see \cite{Fhom2}) that if $\bbP$ is the law of the \verde{$\nu$-randomization of  a stationary ergodic SPP with positive intensity}  the following property holds:   for $\bbP$--a.a. $\o$ the set $\widehat \o$ has nonempty intersection with 
%the sets $S_\ell^-$ and  $S_\ell^+$ for any $\ell>0$ and has  nonempty intersection with $\L_\ell$ for  $\ell$ large enough, i.e.~for any  $\ell  \geq \ell_0(\o)$. In the definition  of electrical potential and conductivity given below, we will  understand that the above property  is satisfied (in particular, that $\ell$ is large enough).
%
%
%
%\rosso{Per dare senso alla def sotto mi serve $\sum_{y\in \widehat{\o}} c_{x,y}(\o,\b)<+\infty$ for all $x\in \widehat{\o}$.}
\begin{Definition}[Electrical potential $\V$] \label{def_pot}  Given  $\b,\ell>0$ and  $\o \in \O$,
 we denote by  $\V$
  the \emph{electrical potential} of the resistor network \mammaL$(\o)$  with values $1$ and $0$ on $S_\ell^- $ and $S_\ell ^+$, respectively. In particular, $\V$  is the unique function $\V:  \widehat{\o}\cap S_\ell \to [0,1]$ such that 
\[ {\sum} _{y\in \widehat \o \cap S_\ell} c_{x,y}(\o,\b) \left(  \rosso{\V (x)-\V (y)}\right) =0 \qquad \forall 
x\in \widehat \o \cap \L_\ell\,,
\]
and satisfying the boundary conditions
\[
\begin{cases}
 \V (x) = 1  & \text{ if } x \in \widehat \o \cap S_\ell^-\,,\\
\V (x) = 0  & \text{ if } x \in \widehat \o \cap S_\ell^+\,.
\end{cases}
\]
\end{Definition} Existence and uniqueness of the above electrical potential can be easily derived (\rosso{see~\cite[Section~2]{Fhom2} and note \cite[Remark~2.4]{Fhom2}}).
Since  the electrical potential is $1$ on the left and $0$ on the right, the electrical field points along the first direction. 
We recall that given two nodes $x,y$,
\be\label{corr_liabel} i^\o _{\b, \ell}(x,y):=c_{x,y}(\o,\b) \bigl(\V (x)- \V (y) \bigr)
\en  is the electrical current flowing from $x$ to $y$ under the electrical potential $\V$.
%%%%%%%%%%
\begin{Definition}[$\ell$--conductivity]  \label{def_ell_cond} Given  $\b,\ell>0$ and  $\o \in \O$,
we call $\s_\ell(\o,\b)$ the \emph{effective conductivity} of the resistor network  \mammaL{}  along the first direction under a unit  electrical potential difference.  In particular,  $\s_\ell (\o,\b)$ is given by (cf. \eqref{corr_liabel})
 \be\label{ide1}
\begin{split}
 \s_\ell(\o,\b)& :=  \sum _{x\in \widehat \o \cap S_\ell^-} \; \sum_{y \in  \widehat \o \cap \L_\ell}i^\o _{\b, \ell}(x,y) \,.
\end{split}
 \en
\end{Definition} 



We denote by 
$\bbB_\ell^\o$  the family of  pairs $\{x,y\}$  for which there is an electrical filament  in \mammaL$(\o)$.  In particular, $\bbB_\ell^\o$ is  the family of  pairs $\{x,y\}$ 
with $x\not =y$, $x,y\in \hat \o \cap S_{\ell}$ and $\{x,y\}\cap \L_\ell \not =\emptyset$.



As mentioned in \cite{Fhom2}, for any $\g \in [-\ell/2, \ell/2)$, 
 $\s_\ell(\o)$ equals the current flowing through the hyperplane $\{x\in \bbR^d\,:\, x_1=\g\}$:
\begin{equation}\label{eq_ailo0}
\s_\ell(\o,\b )=\sum _{\substack{x\in \hat \o \cap S_\ell:\\
 x_1\leq \g }} \; \sum_{\substack{ y \in  \hat \o \cap S_\ell:\\
 \{x,y\} \in \bbB^\o_\ell, \, y_1>\g
 }} i^\o _{\b, \ell}(x,y) \,.
 \end{equation}
 Moreover, $\s_\ell(\o,\b )$   also equals  the total dissipated energy:
\begin{equation}\label{eq_ailo}
\s_\ell(\o,\b)= \sum_{ \{x,y\} \in \bbB^\o_\ell } c_{x,y}(\o) \bigl( V^\o _{\b,\ell} (x) - V^\o _{\b,\ell} (y) \bigr)^2\,.
\end{equation}
%To justify \eqref{eq_ailo} we observe that, by collapsing all nodes in $\hat \o \cap S_\ell^-$ into a single node and similarly for $\hat \o \cap S_\ell^+$, one reduces to the same setting of  \cite[Section~1.3]{DS} where \eqref{eq_ailo} is proved.



We recall 
a useful variational  characterization of $\s_\ell(\o, \b) $, usually called \emph{Dirichlet principle}  (cf.~\cite[Eq.~\rosso{(24)} and \rosso{Lemma~5.2}]{Fhom2} and \rosso{\cite[Exercise~1.3.11]{DS}}):
\be\label{papero}
\s_\ell (\o,\b)=\inf \bigl\{ \cD(u) \,\big{|} \, u: \widehat \o \cap S_\ell \to [0,1]\,, u_{|\widehat \o \cap S_\ell^-} \equiv 1
\,, \; u_{|\widehat \o \cap S_\ell^+}\equiv 0 \bigr\}\,,
\en
where $\cD(u)$ is the Dirichlet form
\[ \cD(u) := \sum _{\{x,y\} \in \bbB^\o_\ell }c_{x,y}(\o,\b) \left( u(y)- u(x) \right)^2\,.
\]
%%%%%%%%%%%%%%%%%%%%%%%%%%%%%%

\subsection{Infinite volume effective conductivity}
The following result is a special application of  \cite{Fhom2}[\rosso{Corollary~2.7}].


\begin{Fact}[A. Faggionato  \cite{Fhom2}] \label{teo_eff_cond}
Let $\bbP$ be a stationary ergodic marked SPP  with finite and positive intensity $\rho$. Suppose in addition that 
\be\label{non_periodico}\bbP ( \o \in \O:  \t_x\o\not = \t_{x'} \o  \; \; \forall x\not =x' \text{ in }\rosso{\bbR^d})=1
\en
and that 
 $ \bbE[ \widehat{\o} ([0,1]^d)^2 ]  <+\infty$. 
If $e_1$ is an eigenvector of $D(\b)$, then 
\be\label{nanna}
\s(\b):=\lim _{\ell \to +\infty} \ell^{2-d} \s_\ell(\o,\b)=\rho  \rosso{D(\b)_{1,1}}\,.
\en
\end{Fact}

The hypotheses of Fact \ref{teo_eff_cond} assure the validity of Assumptions (A1),...,(A8) in \cite[Section~2]{Fhom2} (for (A7) and (A8) see \cite[Section~5.4]{Fhom1}). \rosso{We recall that the bound $ \bbE[ \widehat{\o} ([0,1]^d)^2 ]  <+\infty$ imply both \eqref{2nd_moment} and \eqref{0_moment} (see Remark~\ref{chitarrina}).}
As discussed in \cite[Section~2]{FSS}, condition 
\eqref{non_periodico}
 is equivalent to 
\be\label{non_periodico_bis}
\vpz ( \o \in \O:  \t_x\o\not = \t_{x'} \o  \; \; \forall x\not =x' \text{ in } \widehat{\o})=1\,.
\en
If $\bbP=\hat \bbP_{\nu}$ for some non-degenerate probability $\nu$, then \eqref{non_periodico} is automatically satisfied (cf.~\cite[Section~2]{FSS}).
If $\bbP$ is reflection invariant  then $D(\b)$ is diagonal, while if $\bbP$ is isotropic then $D(\b)$ is a multiple of the identity  matrix (cf.~\cite{demasi}[Thm.~4.6--(iii)]). In the above two cases, trivially $e_1$ is an eigenvector of $D(\b)$.
If $\bbP={\rm PPP}[\rho,\nu]$, then  all the  assumptions of Fact \ref{teo_eff_cond}  are satisfied  and $D(\b)$ is a multiple of the identity matrix by isotropy.




Due to Fact \ref{teo_eff_cond}, when  $e_1$ is an eigenvector of $D(\b)$, $\s(\b)=\rho\rosso{D(\b)_{1,1}}$ can be thought of as the effective infinite volume  rescaled conductivity along the first direction of the MA resistor network.

\rosso{Finally, although not used below,  we mention that \cite[Theorem~2.6]{Fhom2} gives the scaling limit of $\s_\ell(\o,\b)$ assuming more generally that $e_1\in {\rm Ker} (D(\b))$ or $e_1\in {\rm Ker}(D(\b))^\perp$.}


%without assuming that $e_1$ is an eigenvector of $D(\b)$ (one needs that $e_1\in {\rm Ker}\big(D(\b)\big)$ or $e_1\in {\rm Ker}\big(D(\b)\big)^\perp$).}
%\begin{Remark}\label{rem_isotropia}
%In our applications $\bbP=\hat\bbP_\mu$. If $\hat \bbP$ is left invariant by reflection along the coordinate axes or  is left invariant by $\bbZ^d$ rotation by $\pi/2$, then $D(\b)$ is diagonal {\rm \cite{demasi}[Thm.~4.6--(iii)]}.
%\end{Remark}

%%%%%%%%%%%%%%%%%%%%%%%%%%%%%%%%%%%%%%%%
%
%
%
%
%
%%%%%%%%%%%%%%%%%%%%%%%%%%%%%%%%%%%%%%%%


%\section{Parti spostate perche' pesanti}
%
%In what follows we  think of  a resistor network as a weighted non-oriented graph $\bbG=(\bbV, \bbE, c)$, 
%where $\bbV$,  $\bbE$ and $c:\bbE\to [0,+\infty)$ are respectively the vertex set, the edge set and the weight function. The vertexes of $\bbG$ are thought of as the nodes of the network and to each edge $\{x,y\}\in \bbE$ we assign an electrical filament with conductivity $c(\{x,y\})$.
%
%\begin{Remark} The resistor network   MA$[\b](\o)$ corresponds to the weighted non-oriented graph
%$(\widehat \o, \bbB(\widehat \o), c)$, where $\bbB(\widehat \o):= \{ \{x,y\}\,:\, x\not = y \text{ in } \widehat \o\}$ and 
%$c(\{x,y\}):= c_{x,y}(\o,\b)$.
%\end{Remark}
%
%\begin{Remark}
%We can think  of $ \text{\mammaL}(\o)$  as 
% a weighted non--oriented graph:
%   \be\label{biscottino}
%   \text{\mammaL}(\o)=( \widehat\o \cap S_\ell , \bbB_\ell(\widehat \o) , c)
%   \en
%   with edge set 
%\be\label{biscotto} \bbB_\ell(\widehat \o) :=\bigl\{\,\{x,y\}\,:\, x \in  \widehat \o \cap \L_\ell\,,\; y \in \widehat \o \cap S_\ell\,,\; x\not =y 
%\bigr\}
%\en
%  and edge weight  $c(\{x,y\}):=c_{x,y}(\o,\b)$.
%\end{Remark}



%%%%%%%%%%%%%%%%%%%%%%%%%%%%%%%%%%%%
%
%
%     MOTT'S RANDOM WALK 
%
%
%%%%%%%%%%%%%%%%%%%%%%%%%%%%%%%%%%%%%


\section{Mott's random walk}\label{sec_mott_rw}
Given $\o \in \O$, Mott's random walk (introduced in \cite{FSS}) is the continuous-time 
Markov chain with state space $\widehat\o$ and with jump rates $c_{x,y}(\o,\b)$: once arrived at the site $x\in \widehat\o$ the random walk waits an exponential time of rate $c_x(\o,\b):=\sum_{y\in \widehat\o}c_{x,y}(\o)$, afterwards it jumps to a different site $y\in \widehat\o$ with probability $c_{x,y}(\o,\b)/c_x(\o,\b)$ (recall that $c_{x,x}(\o,\b):=0$). As random walk in a random environment, Mott's random walk  belongs to the class of random conductance models \cite{Bis}.
 
 To have a well defined Markov chain we assume that $c_{x}(\o,\b)<+\infty$ for all $x\in \widehat\o$ and that no explosion occurs. Due to \cite[Lemma~3.5]{Fhom1} and the discussion 
in \cite[Section~5.4]{Fhom1} devoted to Mott's random walk, the above  conditions are satisfied for $\bbP$--a.a.~$\o$ if $\bbP$ is a stationary ergodic  marked SPP with finite and positive intensity $\rho$ satisfying \eqref{non_periodico} and $\bbE[ \widehat{\o} ([0,1]^d)^2 ]  <+\infty$. 


The effective homogenized matrix $D(\b)$ corresponds to the asymptotic diffusion matrix of Mott's random walk as formalized \rosso{Fact~\ref{tolmezzo}} below. We first introduce some notation.
 Given $\e>0$ we write $(P^\e _{\o ,t} )_{t \geq 0}$ for the $L^2(\mu^\e _\o)$--Markov semigroup associated to the random walk $(\e X^\o _{ \e^{-2} t} )_{t\geq 0}$ on $\e \widehat \o$, where $\mu^\e_\o := \e^d \sum _{x\in \widehat \o} \d_{\e x}$. Simply $P^\e _{\o ,t}f(x)$ is the expectation of $f ( \e X^\o _{ \e^{-2} t})$ when the random walk starts at $\e x$.
  Similarly we write $( P_t  )_{t \geq 0} $ for the  Markov semigroup on $L^2( \rho dx)$ associated to the  (possibly degenerate)  Brownian motion on $\bbR^d$  with diffusion matrix $2 D(\b)$. 
%  Given $\l>0$ we write  $R^\e _{\o ,\l}: L^2(\mu^\e _\o)\to L^2(\mu^\e _\o) $ for the $\l$--resolvent  associated to the random walk $\e X^\o _{ \e^{-2} t} $, i.e.  $R^\e _{\o ,\l}:= (\l -\bbL^\e _\o )^{-1} =\int_0^\infty e^{- \l s} P^\e _{\o ,s} ds $. 
% We write $ R_\l : L^2(m dx) \to L^2(m dx)$ for the $\l$--resolvent associated to the above Brownian motion on $\bbR^d$  with diffusion matrix $2 D$.
 The following fact is   a  very special application of \orc{a much more general result given by  \cite[Theorem~4.4]{Fhom1}}:
 % (which applies to a much larger class of random walks on simple point processes and also contains further results):
 \begin{Fact}[A.~Faggionato \cite{Fhom1}] \label{tolmezzo}
 Let $\bbP$ be a stationary ergodic marked SPP  with finite and positive intensity $\rho$. Suppose in addition that $\bbP$ satisfies \eqref{non_periodico}
and that 
 $ \bbE[ \widehat{\o} ([0,1]^d)^2 ]  <+\infty$. Then for $\bbP$--a.a.~$\o$, given   $t\geq 0$  and $f\in C_c(\bbR^d)$, as $\e\da 0$  it holds
$ L^2(\mu^\e _\o ) \ni P^\e_{\o,t} f \to P_t f \in L^2(\rho dx)$.
 \end{Fact}
 The precise  meaning of the   convergence $ P^\e_{\o,t} f \to P_t f$  is recalled in \cite[Section~3.5]{Fhom1}.
 The above assumptions on the law $\bbP$ are the same of the ones in   Fact \ref{teo_eff_cond} (with exception that here we do not require that $e_1$ is an eigenvector of $D(\b)$). Hence the comments following  Fact \ref{teo_eff_cond} can be applied to the present setting \rosso{as well}.
 
 Stronger form of CLTs or even  invariance principles (annealed or quenched) have been proved under stronger conditions (we stress that the basin of application of Fact \ref{tolmezzo} is \rosso{very large}). 
 We recall that the annealed invariance principle of Mott's random walk has been obtained in \cite{FSS}. The quenched invariance principle has been obtained in \cite{CF1} for $d=1$ and in \cite{CFP} for $d\geq 2$. In the above papers, the hypotheses also \rosso{imply} the non-degeneracy of the diffusion matrix. 
 
 Bounds on the diffusion matrix  in agreement with Mott's law have been obtained  under suitable assumptions in \cite[Thm.~1]{FM} and \cite[Thm.~1]{FSS}  for $d\geq 2$, while bounds 
 in agreement with an Arrhenius-type behavior  have been obtained in 
  \cite[Thm.~1.2]{CF1} for $d=1$. In particular \rosso{these} bounds are of the form 
  \begin{align*}
  & c_1 \exp \bigl\{ -c_1' \, \beta ^\frac{\a+1}{\a+d +1}\bigr\}\mathbb{I} \leq D(\b)\leq 
c_2 \exp \bigl\{ -c_2' \, \beta ^\frac{\a+1}{\a+d +1}\bigr\}\mathbb{I} \qquad &(d\geq 2)\,,\\
& c_1 \exp \bigl\{ -c_1' \, \beta  \bigr\} \leq D(\b)\leq  c_2 \exp \bigl\{ -c_2' \, \beta \}\qquad & (d=1)\,,
\end{align*}
for suitable $\b$--independent positive  constants $c_1,c_1',c_2,c_2'$.
%For $d= 1$  annealed and quenched  invariance principles  have been obtained in \cite[Thm.~1.1]{CP1}.  Again $D(\b)$ has a variational characterization and satisfies bounds in agreement with \eqref{mottino1} (see \cite[Thm.~1.2]{CP1}):
%\be \label{silenzioso2}
%c_1 \exp \bigl\{ -c_1' \, \beta  \bigr\} \leq D(\b)\leq 
%c_2 \exp \bigl\{ -c_2' \, \beta \}\,,
%\en
%for suitable $\b$--independent positive  constants $c_1,c_1',c_2,c_2'$.
\rosso{Under Fact \ref{teo_eff_cond}}, as pointed out in \cite[Corollary~3.1]{Fhom2},  the above  bounds imply similar bounds for  $\s(\b)/\rho$ due to \eqref{nanna}.
%%%%%%%%%%%%%%%%%%%%%%%%%%%%%%%%%%%%%%%%
%
%
%
%
%
%%%%%%%%%%%%%%%%%%%%%%%%%%%%%%%%%%%%%%%%

\section{Graph $\cG[\z,\b](\o)$ and  critical conductance for $d\geq 2$}\label{grafene}
We restrict to dimensions  $d\geq 2$ (this will be understood in what follows).

\medskip

The infinite volume MA resistor network, i.e.~\mamma$(\o)$ introduced in Definition \ref{def_MA}, corresponds to a weighted complete graph with vertex set $\widehat{\o}$. On the other hand, at small temperature, one expects that the  electrical filaments with a conductivity lower bounded by  a suitable $\b$--dependent threshold  give the main contribution to the electron transport, apart from a negligible term. It is therefore not surprising that a crucial role in our analysis will be played by the following  graph $\cG[\z,\b](\o)$:
 
 
\begin{Definition}[Graph \giggioest] \label{giggio_bello} Given $\o\in \O$, $\b>0$, $\z>0$, the  graph \giggio$(\o)$ has vertex set $\widehat \o$ and has edges $\{x,y\}$ with $x\not =y$ in $\widehat \o$ such that 
$c_{x,y}(\o,\b) \geq e^{-\z}$, i.e. 
\be \label{connectcont}
|x-y|+\beta(|E_x|+|E_y|+|E_x-E_y)|)\leq \z\,.
\en
\end{Definition}
Note that $\cG[\z,\b](\o)$  is the graph in $\bbR^d$ with vertex set $\widehat \o$ and supporting the family of filaments in the  resistor network MA$[\b](\o)$ with conductivity at least $e^{-\z}$.



\medskip

We say that the graph $\cG[\z,\b](\o) $ percolates if it has  some unbounded connected component and we set 
\be
\theta (\z,\b, \rho, \nu):= {\rm PPP}[\rho, \nu]  \Big( \cG [\z,\b] \text{ percolates} \Big)\,.
\en
We recall that, according to our notation, the above r.h.s. is the probability for ${\rm PPP}[\rho, \nu]  $ of the event that $\cG [\z,\b]$ percolates (see Definition \ref{def_PPP}). 

By ergodicity w.r.t.~spatial translations and since the existence of an unbounded connected component is a translationally  invariant event,  we have  that $\theta(\z,\b, \rho, \nu) \in \{0,1\}$. %By Lemma \ref{transitare} in Appendix \ref{proof_transitare}
\rosso{As derived in Appendix \ref{guizzo}}, since $d\geq 2$  there exists 
a critical threshold  $\z_c\in (0, +\infty) $  depending on $\b$, $\rho$ and  $\nu$ such that 
\be\label{critico}
\begin{cases}
\theta(\z,\b, \rho, \nu) =1 & \text{ if } \z > \z_c(\b,\rho,\nu)\,,\\
\theta(\z,\b, \rho, \nu)=0 & \text{ if } \z< \z_c(\b,\rho,\nu)\,.
\end{cases}
\en




Similarly to \cite[Section~IV.B]{AHL} for the critical conductance, we set the following definition:
\begin{Definition}[Critical \rosso{threshold} $\z$ and conductance] \label{criticone}    The value $\z_c(\b,\rho,\nu)$  in \eqref{critico}  and  the value 
\[c_c(\b,\rho,\nu):= \exp\{-  \z_c(\b,\rho,\nu) \}\]
  are  called, respectively,  \emph{critical \rosso{threshold} $\z$}   and  \emph{critical conductance} of the  MA random resistor network when $\bbP={\rm PPP}[\rho,\nu]$.
\end{Definition}
%From the physical viewpoint, the numbers in \eqref{connectcont} are pure numbers. Indeed,\rosso{$|x-y|$ is divided by  $\g/2$ in \eqref{condu} and the localization length $\g$ has been set equal to 2  length units. Therefore we do not call $\z_c$ for example as critical length.} 
 

%%%%%%%%%%%%%%%%%%%%%%%%%%%%%%%%
%%%%%%%%%%%%%%%%%%%%%%%%%%%%%%%%
\section{Main results for  ${\rm PPP}[\rho, \nu]$, $d\geq 2$}\label{MR_PPP}

\rosso{In this section we restrict to dimension  $d\geq 2$ without further mention}.


\begin{Definition}[\rosso{Probabilities $\nu^+ _{C_0,\a}$ and $\nu_{C_0,\a}$}]\label{almograve}
Given   $C_0>0$ and $\a\geq 0$, we introduce the probability measures $\nu^+ _{C_0,\a}$ and $\nu_{C_0,\a}$ with support $[0,C_0]$ and $[-C_0,C_0]$ respectively as follows:
  \begin{align}
&\nu^+ _{C_0,\a}(dE)=\frac{(\a+1)}{  C_0 ^{\a+1} }  E^\a \rosso{\mathds{1}_{[0,C_0]}(E)}  dE\,,\label{speciale2}\\
&\nu_{C_0,\a} (dE)=\frac{(\a+1)}{ 2 C_0 ^{\a+1} }  |E|^\a  \rosso{\mathds{1}_{[-C_0,C_0]}(E)}   dE\,.\label{speciale1}
\end{align}
\end{Definition}
Recall the graph $\cG[\z,\b]$ introduced in Definition \ref{giggio_bello}.  
%and denote by  ${\rm PPP}^0(\l, \nu)$  the Palm distribution of ${\rm PPP}(\l, \nu)$. The following fact introduces the critical intensities  $\l^+_c(\a)$ and $\l_c(\a)$ (see Appendix \ref{guizzo} for a derivation based on \cite[Prop.~2.2]{FagMim1}):}
%Due to  \cite[Prop.~2.2]{FagMim1}    a phase transition appears in the random graphs $\cG[1, 1]$ under ${\rm PPP}[\l,\nu^+_{1,\a}] $ and under ${\rm PPP}[\l,\nu_{1,\a}] $  by changing the intensity  $\l$.  This justifies the following definition:
%%%%%%%%%%%%%%%%%%%%%%%%%%%%%%%%%%%%%%%%%%%%%%%%%
\begin{Definition}[Critical intensities $\l^+_c(\a)$ and $\l_c(\a)$]
\label{fiordo} Given $\a\geq 0$ we  define $\l^+_c(\a) \in(0,+\infty)$  as the critical intensity such that   the graph $\cG[1, 1]$ under ${\rm PPP}[\l,\nu^+_{1,\a}] $ a.s. percolates  if $\l >\l^+_c(\a)$, while  it 
a.s.   does not percolate if $\l<\l^+_c(\a)$.
We define  similarly  $\l _c(\a)\in(0,+\infty)$  by replacing  $\nu^+_{1,\a}$ with $\nu_{1,\a}$.
  \end{Definition}
\rosso{The existence and non-triviality of the above critical intensities  $\l^+_c(\a)$ and  $\l_c(\a)$ can be easily proved by stochastic domination (in particular,  by comparing the random graph $\cG[1,1]$ with suitable  Poisson Boolean models as in  \cite{FagMim1,FagMim2}).  We sketch the derivation in   Appendix \ref{guizzo}, where we also derive another useful characterization of the critical intensities in terms of the Palm distribution and the connected component of the origin}.

%%%%%%%%%%%%%%%%%%%%%%%%%%%%%%%%%%%%%%%%%%%%%%%%%%%%%%%%%%%%%%%%%%%%%%%%%%%%%%%%%%%%%%%%%%%%%%%%%%%%%%%%%%%%%%%%%%%%%%%%%%%%%%%%%%%%%%%%%%%%%%%%%%%%%%%%%%%%%%%%%%%%%%
%\begin{TheoremA} \label{teo1} \rosso{si puo' generalzzare}
%Let $C_0>0$ and $\a\geq 0$.  
% Then the  following holds:
%\begin{itemize}
%\item 
%For  $ \beta >   \big( \l^+_c(\a) /\rho\big) ^{ 1/d} C_0^{-1}$ (hence at low temperature), the critical length of the resistor network
%${\rm MA}(\b)$  under $\bbP(\rho, \nu^+_{C_0,\a})$  is given by 
%\be
%\label{zetacritico+}
%\z_c(\b,\rho, \nu^+_{C_0,\a} )= \bigl( {\l^+_c(\a)/\r}\bigr)^{\frac{1}{\a+1+d}}
%  (  \b C_0) ^{\frac{\a+1}{\a+1+d}}\,.
%\end{equation}
%\item  For  $ \beta >   \big( \l_c(\a) /\rho\big) ^{ 1/d} C_0^{-1}$ (hence at low temperature), the critical length of the
%resistor network
%${\rm MA}(\b)$  under $\bbP(\rho, \nu_{C_0,\a})$  is given by 
%\be
%\label{zetacritico}
%\z_c(\b,\rho, \nu_{C_0,\a} )= \bigl( {\l_c(\a)/\r}\bigr)^{\frac{1}{\a+1+d}}
%  (  \b C_0) ^{\frac{\a+1}{\a+1+d}}\,.
%\end{equation}
%\end{itemize}
%\end{TheoremA}
%%%%%%%%%%%%%%%%%%%%%%%%%%%%%%%%%%%%%%%%%%%%%%%%%%%%%%%%%%%%%%%%%%%%%%%%%%%%%%%%%%%%%%%%%%%%%%%%%%%%%%%%%%%%%%%%%%%%%%%%%%%%%%%%%%%%%%%%%%%%%%%%%%%%%%%%%%%%%%%%%%%%%%%%%%%%%%%%%%%%%%%%%%%%%%%%%%%%%%%%%%%%%%%%%%%%%%%%%%%%%%%%%%%%

%In the next results we treat two fundamental cases:
\begin{Definition}[Classes $\cvp$ and $\cv$, critical density $\l_c^*(\a)$] \rosso{Let $\a\geq 0$.}
 \label{monte_pisciano} 
$\;\; $
\begin{itemize}
\item We define $\cvp$ as the family of probability measures $\nu$ such that $\nu$ restricted to a neighborhood of the origin is proportional to $E^\a \mathds{1}(E\geq 0) dE$, i.e.~for some $\e,C_0>0$   we have   $\nu(B)=\nu^+_{C_0,\a}(B) $ for each Borel set $B\subset[-\e,\e]$.
\item  We define $\cv$ as the family of  probability measures $\nu$ such that $\nu$ restricted to a   neighborhood of the origin is proportional to $|E|^\a  dE$,  i.e.~for some  $\e,C_0>0$  we have   $\nu(B)=\nu_{C_0,\a}(B) $ for each Borel set $B\subset[-\e,\e]$.
\item When dealing with $\nu \in \cvp$ or $\nu \in \cv$ we set
\begin{equation}\label{nuziale}
\l_c^*(\a):=
\begin{cases}
\l^+_c(\a) & \text{ if } \nu\in \cvp ,\\
\l_c(\a) & \text{ if } \nu \in \cv\,.\\
\end{cases}
\end{equation}
\end{itemize}

\end{Definition}


%\begin{Definition}[Cases (A) and (B) for $\nu$] \label{monte_pisciano} 
%We consider the following fundamental cases for the measure $\nu$ of the energy marks:
%\begin{itemize}
%\item[(A)] The measure $\nu$ restricted to a neighbourhood of the origin is proportional to $E^\a \mathds{1}(E\geq 0) dE$ for some $\a\geq 0$. In particular,  for suitable constants $\e,C_0>0$ and $\a\geq 0$, we have   $\nu(B)=\nu^+_{C_0,\a}(B) $ for each Borel set $B\subset[-\e,\e]$.
%\item[(B)] The measure $\nu$ restricted to a neighbourhood of the origin is proportional to $|E|^\a \mathds{1}(E\geq 0) dE$ for some $\a\geq 0$. In particular,  for suitable constants $\e,C_0>0$ and $\a\geq 0$, we have   $\nu(B)=\nu_{C_0,\a}(B) $ for each Borel set $B\subset[-\e,\e]$.
%\end{itemize}
%Moreover, we set
%\begin{equation}
%\l_c^*(\a):=
%\begin{cases}
%\l^+_c(\a) & \text{ in Case } {\rm (A)}\,,\\
%\l_c(\a) & \text{ in Case } {\rm (B)}\,.\\
%\end{cases}
%\end{equation}
%\end{Definition}



\begin{Warning}\label{pericolo} When we assume  that $\nu\in \cvp$  or $\nu\in \cv$,  the constants $C_0$ and $\e$ will  be defined automatically  as the ones appearing in Definition \ref{monte_pisciano}. If $\nu=\nu^+ _{C_0,\a}$ or $\nu= \nu _{C_0,\a}$, then \rosso{we take}  $\e=C_0$.
\end{Warning}
%%%%%%%%%%%%%%%%%%%%%%%%%%%%%%%%%%%%%%%%%%%%%%%%%%%%%%%%%%%%%%%%%%%%%%%%%%%%%%%%%%%%%%%%%%%%%%%%%%%%%%%%%%%%%%%%%%%%%%%%%%%%%%%%%%%%%%%%%%%%%%%%%%%%%%%%%%%%%%%%%%%%%
Recall Definition \ref{criticone} and in particular that  $c_c(\b,\rho,\nu):= \exp\{-  \z_c(\b,\rho,\nu) \}$.
\begin{Theorem}[Critical \rosso{threshold} $\z$ and conductance   of ${\rm MA}\text{$[\b]$}$] \label{teo1}  \rosso{Assume that $\nu \in \cvp\cup\cv$}. Then 
 the critical \rosso{threshold} $\z$ and the critical conductance  of the resistor network
${\rm MA}[\b](\o)$  with $\o$ sampled according to ${\rm PPP}[\rho, \nu]$  are given by 
\begin{align}
\z_c(\b,\rho, \nu )& = \bigl( {\l^*_c(\a)/\r}\bigr)^{\frac{1}{\a+1+d}}
  (  \b C_0) ^{\frac{\a+1}{\a+1+d}}\,, \label{zetacritico}\\
c_c(\b,\rho, \nu )& = \exp\big\{- \bigl( {\l^*_c(\a)/\r}\bigr)^{\frac{1}{\a+1+d}}
  (  \b C_0) ^{\frac{\a+1}{\a+1+d}}\big\}\, ,\label{ccritico}
\end{align}
if  $\beta$ is large enough. More precisely, \eqref{zetacritico} and \eqref{ccritico} hold  whenever the  r.h.s.  of \eqref{zetacritico}  \orc{is smaller than} $\min\{C_0, \e\}\b$. 
\end{Theorem}
%%%%%%%%%%%%%%%%%%%%%%%%%%%%%%%%%%%%%%%%%%%%%%%%%%%%%%%%%%%%%%%%%%%%%%%%%%%%%%%%%%%%%%%%%%%%%%%%%%%%%%%%%%%%%%%%%%%%%%%%%%%%%%%%%%%%%%%%%%%%%%%%%%%%%%%%%%%%%%%%%%%%%%%%%%%%%%%%%%%%%%%%%%%%%%%%%%%%%%%%%%%%%%%%%%%%%%%%%%%%%%%%%%%


The proof of Theorem \ref{teo1} is provided in Section \ref{sec_proof_teo1}. 
%Since the critical conductance $c_c(\b,\rho,\nu)$ is defined as $c_c(\b,\rho,\nu):= \exp\{-  \z_c(\b,\rho,\nu) \}$,  
Theorem \ref{teo1} immediately implies the following:
%%%%%%%%%%%%%%%%%%%%%%%%%%%%%%%%%%%%%%%%%%%%%
\begin{Corollary}[Value $\chi(\rho,\nu)$]\label{luna}
Suppose that $\nu \in \cvp\cup \cv$ and consider the resistor network
${\rm MA}[\b](\o)$  with $\o$ sampled according to ${\rm PPP}[\rho, \nu]$. 
Then 
%the asymptotic (rescaled) critical conductance $c_c(\b,\rho,\nu)  $ defined as 
%\[
%c_c(\rho, \nu):=\lim _{\b\to +\infty} \b^{- \frac{1+\a}{d+1+\a} } \ln c_c(\b,\rho,\nu )
%\]
%is given by 
%Since the critical conductance $c_c(\b,\rho,\nu)$ is defined as $c_c(\b,\rho,\nu):= e^{-  \z_c(\b,\rho,\nu) }$,  Theorem \ref{teo1} implies in particular that
 \be\label{premium}
\chi(\rho, \nu):=\lim _{\b\to +\infty} \b^{- \frac{\rosso{\a+1}}{\rosso{\a+1+d}} } \ln c_c(\b,\rho,\nu )=  - \bigl( {\l^*_c(\a)   C_0^{\a+1}/\r}\bigr)^{\frac{1}{\a+1+d}}
     \,.
\en
\end{Corollary}
%%%%%%%%%%%%%%%%%%%%%%%%%%%%%%%%%%%%%%%%%%%%%%
\smallskip


We now present our results for Mott's law. For the next result recall  the notational Warning \ref{pericolo}. We point out that, since the PPP is isotropic, $D(\b)$ is a multiple of the identity and $e_1$ is always an eigenvector of $D(\b)$. In particular \eqref{nanna} is satisfied.
% kuka Commentare che anche le altre assunzioni del teorema che porta a \eqref{nanna} sono soddisfatte

%%%%%%%%%%%%%%%%%%%%%%%%%%%%%%%%%%%%%%%%%%%%%
  \begin{Theorem}[Upper bound in Mott's law] \label{parataUB} Suppose that  $\nu\in \cvp\cup \cv$ and take $\bbP={\rm PPP}[\rho,\nu]$.
Then we have 
 \begin{align}
 &  \limsup _{\b\to +\infty} \b^{- \rosso{\frac{\a+1}{\a+1+d}} } \ln  D(\b)_{1,1}  \leq
 %  -  \left( \l_c^*(\a)   C_0^{\a+1}/\rho)\right)^{\frac{1}{\a+1+d}}  =
  \chi (\rho, \nu)\,, \label{preUB}\\
&  \limsup _{\b\to +\infty} \b^{- \rosso{\frac{\a+1}{\a+1+d}} } \ln \s(\b)\leq   
%  -  \left( \l_c^*(\a)   C_0^{\a+1}/\rho)\right)^{\frac{1}{\a+1+d}}  = 
\chi (\rho, \nu)
\,.\label{UB}
  \end{align}
 % where the second bound holds whenever $e_1$ is an eigenvector of $D(\b)$.
\end{Theorem}
The proof of Theorem \ref{parataUB} is provided in Section \ref{sec_UB}. We point out that \eqref{UB} is an immediate consequence of \eqref{preUB} since we have  \eqref{nanna}. The strategy to prove \eqref{preUB} will be to use the variational formula \eqref{def_D} for $D(\b)$ for clever choices of the test function $f$ (partially inspired by the proof of \cite[Theorem~3.12]{PR} and several scaling considerations discussed in the next sections).
% kuka: controllare teorema 3.12 in \cite{PR}

\smallskip

We now move to the lower bound in Mott's law.
To this aim, first we need some definitions.
Recall first the definition of $S_\ell, S_\ell^-, S_\ell^+, \L_\ell$ given at the beginning of Section \ref{fumata}.
\begin{Definition}\label{def_LR} Let $\bbG=(\bbV,\bbE)$ be a graph with $\bbV\subset \bbR^d$.
A left-right (LR) crossing  of  the box $\L_\ell$ in the graph  $\bbG$    is any  sequence of distinct points $x_1$, $x_2$,$\dots$, $x_n \in \bbV$  with $n\geq 3$  such that 
%\begin{itemize}
%\item 
$\{x_i ,x_{i+1}\} \in  \bbE$ for all $i=1,2,\dots, n-1$;
%\item 
$x_1 \in S_\ell^{-} $ and   $x_n \in S_\ell^+ $;
%\item 
$x_2, x_3,\dots, x_{n-1}\in \L_\ell$.
%\end{itemize}
We also define $\cN_\ell(\bbG)$ as the  maximal number of  vertex-disjoint LR crossings of $\L_\ell$ in $\bbG$.
\end{Definition}
\rosso{By the above notation $\bbG=(\bbV,\bbE)$, we mean that $\bbV$ is the vertex set and $\bbE$ is the edge set, with $\bbE\subset \big\{ \,\{x,y\} \,:\, x,y\in \bbV,\;x\not=y \text{ in }\bbV\big\}$.}


\begin{Definition}\label{good_LR_distr}
Let $\bbG=(\bbV,\bbE)$ be a random  graph  with $\bbV\subset \bbR^d$.
We say that $\bbG$ has a   \emph{good LR crossing  distribution} if   a.s. there exists 
 $c>0$ such that  $\cN_\ell(\bbG)\geq c\, \ell^{d-1}$ for $\ell $ large enough.
% it holds $\sum_{\ell=1}^\infty P\left( \cN_\ell(\bbG)\leq c\, \ell^{d-1}\right)<+\infty$  for some   constant $c>0$.
 \end{Definition}
We point out that the constant $c$ in the above definition can be random.

The following fact \rosso{(valid for $d\geq 2$)} is obtained by applying   \cite[Thm.~1]{FagMim2} and by using arguments similar to the ones in  \cite[Section~8]{FagMim2} to check the hypotheses. We detail its derivation in Appendix \ref{pinco}. We mention that \cite{FagMim2} covers a large class of graphs built on PPP's, including the   graph $\cG[1,1]$. 
%%%%%%%%%%%%%%%%%%%%%%%%%%%%%%%%%%%%%
\begin{Fact}[A.~Faggionato \& A.H.~Mimun \cite{FagMim2} and Appendix \ref{pinco}] \label{cortocircuito}
%Let $\bbP_*={\rm PPP}(\l, \nu^+_{1,\a})$ with 
 Fixed
 $\l > \l^+_c(\a)$ there exist  constants $c_1,c_2>0$ depending on $\l$ \rosso{and} $\a$ such that 
 \be
{\rm PPP}[\l , \nu^+_{1,\a}] \left( \cN_\ell (\cG[1,1]) \geq  c_1 \ell^{d-1}\right)\geq 1- \exp\left \{- c_2 \,\ell ^{d-1}\right\}\,,
\en
for $\ell$ large enough.
In particular, by Borel-Cantelli lemma, the random graph $\cG[1,1](\o)$ with $\o$ sampled according to 
${\rm PPP}[\l , \nu^+_{1,\a}] $ has a good LR crossing distribution.
\end{Fact}

\rosso{The above fact on LR crossings is one of the main tools in the proof of the following result:}
  \begin{Theorem}[Lower bound in Mott's law] \label{parataLB} Suppose that  $\nu\in \cvp$ and take $\bbP={\rm PPP}[\rho,\nu]$.
Then we have 
 \begin{align}
 &  \liminf _{\b\to +\infty} \b^{- \rosso{\frac{\a+1}{\a+1+d}} } \ln D(\b)_{1,1}  \geq  \chi(\rho, \nu)\,, \label{preLB}\\
&  \liminf _{\b\to +\infty} \b^{- \rosso{\frac{\a+1}{\a+1+d}} } \ln \s(\b)\geq    \chi (\rho, \nu)
\,.\label{LB}
  \end{align}
 % where the second bound holds whenever $e_1$ is an eigenvector of $D(\b)$.
  
  If in addition  the random graph $\cG[1,1](\o)$ with $\o$ sampled according to ${\rm PPP}[ \l, \nu_{1,\alpha}] $ has a good LR crossing distribution  for all $\l>\l_c(\a)$, then 
  \eqref{preLB} and \eqref{LB} hold also for  $\nu\in \cv$ and $\bbP={\rm PPP}[\rho,\nu]$.  \end{Theorem}
The proof of Theorem \ref{parataUB} is provided in Section \ref{sec_LB}. Due to \eqref{nanna},  \eqref{LB} implies \eqref{preLB}. We will  prove \eqref{LB} by means of  Rayleigh monotonicity law for resistor networks.


\medskip

Theorems \ref{parataUB} and \ref{parataLB} have the following immediate consequence:
\begin{Corollary}[Mott's law] \label{immenso}
Suppose that  $\nu\in \cvp$ and take $\bbP={\rm PPP}[\rho,\nu]$.
Then we have 
 \[  \lim _{\b\to +\infty} \b^{- \rosso{\frac{\a+1}{\a+1+d}} } \ln D(\b)_{1,1}  =     \lim _{\b\to +\infty} \b^{- \rosso{\frac{\a+1}{\a+1+d}} } \ln \s(\b)=    \chi (\rho, \nu)\,.
\]
  The same results hold for $\nu\in \cv$ and $\bbP={\rm PPP}[\rho,\nu]$ if 
the random graph $\cG[1,1](\o)$ with $\o$ sampled according to ${\rm PPP}[ \l, \nu_{1,\alpha}] $ has a good LR crossing distribution  for all  $\l>\l_c(\a)$.
\end{Corollary}

\begin{Remark} The property that the constant $\chi(\rho,\nu)$ appears in all the above asymptotics confirms the validity of the paradigm of critical path analysis (CPA) in our context (see \cite{Ji} for other CPA  rigorous results in   a different model of transport in disordered media).
\end{Remark}
%%%%%%%%%%%%%%%%%%%%%%%%%%%%%%%%%%%%%%%%%%%%%%%%%%%%%%%%%%%%%%%%%%%%%%%%%%%%%%%%%%%%%%%%%%%%%%%%%%%%%%%%%%%%%%%%%%%%%%%%%%%%%%%%%%%%%%%%%%%%%%%%%%%%%%%%%%%%%%%%%%%%%%%%%%%%%%%%%%%%%%%%%%%%%%%%%%%%%%%%%%%%%%%%%%%%%%%%%%%%%%%%%%%%%%%%%%%%%%%%%%%%%%%%%%%%%%%%%%%%%%%%%%%%%%%%%%%%%%%%%%%%%%%%%%%%%%%%%%%%%%%%%%%%%%%%%%%%%%%%%%%%

\section{Main results concerning Universality}\label{MR_univ}


 The results presented in this section hold for all dimensions $d\geq 1$.
 
 \medskip 
  \blu{The derivation of the results presented  in the previous section holds in good part for generic marked SPPs. As the reader will see, in the derivation of the  upper  and lower bounds in agreement with Mott's law  (cf.~Sections \ref{sec_UB} and \ref{sec_LB}), a central role is played by a  mechanism of thinning and rescaling which ultimately  captures the backbone of  filaments in the MA resistor network mainly contributing to the electron transport. This mechanism is described by the marked SPP \eqref{gandhi} in Theorem \ref{paradiso} below, which will frequently appear in the next sections. We will give more detail at the end of this section. }
  
 
 \blu{We show here  that, also  when starting with a marked SPP not of the form ${\rm PPP}[\rho,\nu]$, this mechanism asymptotically leads to ${\rm PPP}[\rho,\nu^+_{1,\a}]$ or ${\rm PPP}[\rho,\nu_{1,\a}]$, thus explaining the universality behind Mott's law as Physics law and the central role played by the  PPP.}
%Indeed, a suitable quantitative version of Theorem \ref{paradiso} below would allow to extend the proofs in  Sections \ref{sec_UB} and \ref{sec_LB} (and therefore the results in the previous section) to more general marked SPPs.}



% \smallskip
   %Recall the probability  measures  $\nu^+_{1,\a}$ and $\nu_{1,\a}$ defined as in  \eqref{speciale2} and \eqref{speciale1} with $C_0=1$. 
 \begin{Definition}[Probability measure $\nu_{\star, \g}$] \label{stella_stellina}
Given   a probability measure $\nu$ on $\bbR$ having $0$ in its support and given $\g>0$, we define  $\nu_{\star, \g}$  as the probability measure with support inside $[-1,1]$ \rosso{such that}
 \begin{equation}
 \nu_{\star, \g}(B):= \frac{\nu(\g B \cap [-\g, \g])}{\nu([-\g, \g])}\,, \qquad B \in \cB (\bbR)\,.
 \end{equation}
 \end{Definition}
 Note that $ \nu _{\star,\g} ( \cdot )= P\big(\,  X/\g \in \cdot \,\big{|}\, |X/\g |\leq 1\,\big) $, where $X$ is a random variable with distribution  $\nu$.
\begin{Theorem}\label{paradiso} Consider a stationary ergodic SPP  on $\bbR^d$ with law  $\hat \bbP$  and  intensity $\rho \in (0,+\infty) $ and consider  a probability measure $\nu$ on $\bbR$ with $0$ in its support such that 
\begin{itemize}
\item[(i)] $\nu_{\star, \g}$  converges \rosso{weakly}  to $\nu^+_{1,\a}$ or to  $\nu_{1,\a}$ as $\g \da 0$;
%\be \label{vichinghi50} \lim_{\ell \da 0}\nu_{\star, \g}=\nu^+_{1,\a} \quad \text{  or }\quad \lim_{\ell \da 0}\nu_{\star, \g}=\nu_{1,\a}\,; \en
\item[(ii)]    the following limit exists, finite and  positive:
 \be\label{vichinghi100}
C_*:=\lim _{\g \da 0} \frac{\nu([-\g, \g] )}{ \g ^{\a+1}}\in (0,+\infty)\,.
\en
\end{itemize}
Fix $\l>0$ and set 
\be\label{canditi105} \ell(\b):= \bigl(\l /(C_* \rho)\bigr)^{\frac{1}{\a+1+d}} \b ^{\frac{\a+1}{\a+1+d}}\,.
\en
Then, when  $\o $ is  sampled according to   $\hat{\bbP}_\nu$, the rescaled marked SPP \be\label{gandhi}
\left\{ \Big ( \frac{x}{\ell(\b)}, \frac{\b}{\ell(\b)} E_x\Big) \,: \, x\in\widehat \o  \,, |E_x| \leq \frac{\ell(\b)}{\b} \right \}
\en 
 converges \rosso{weakly} to   
\be\label{ghibaudo}
\begin{cases}
{\rm PPP}[\l, \nu^+_{1, \a}] & \text{ if } \lim_{\g\da 0}\nu_{\star, \g}=\nu^+_{1,\a}\,,\\
{\rm PPP}[\l, \nu_{1, \a}] & \text{ if }\lim_{\g \da 0}\nu_{\star, \g}=\nu_{1,\a}\,.
\end{cases}
\en
%\verde{In the special case $\bbP={\rm PPP}[\rho,\nu]$ and $\nu= \nu^+_{1,\a}$ or   $\nu=\nu_{1,\a}$  the rescaled marked SPP \eqref{gandhi} is given by \club \club \club} 
\end{Theorem}
The proof of the above theorem is given in Section \ref{sec_dim_paradiso}.  \rosso{As stated in Proposition \ref{paradiso_PPP} below, if $\nu \in \cvp$ or $\nu\in\cv$, then conditions (i) and (ii) of the above theorem are automatically satisfied. Moreover, for} the marked SPPs treated in Section \ref{MR_PPP}, the limit point in the weak convergence of \eqref{gandhi} described above is achieved already for $\b$ large enough
(recall the notational Warning \ref{pericolo}):
%%%%%%%%%%%%%%%%%%%%%%%%%%%%%%%%

\begin{Proposition}\label{paradiso_PPP} Let $\nu\in \cvp\cup\cv$.
  Then, for    $0<\g \leq \min\{C_0,\e\} $, it holds 
\begin{itemize}
\item[(i)]    $\nu_{\star, \g} = \nu^+_{1,\a}$ if $\nu \in \cvp$ and   $\nu _{\star, \g}=\nu_{1,\a}$ if $\nu\in \cv$; 
\item[(ii)] $\nu([-\g, \g])= \g^{\a+1}/C_0^{\a+1}$, in particular   \eqref{vichinghi100} holds with $C_*:= C_0 ^{-\a-1}$\,.
\end{itemize}
 Given $\rho, \l>0$ define  $\ell(\beta)$ as in \eqref{canditi105}  with  $C_*:= C_0 ^{-\a-1}$, i.e.
\be\label{canditi200}
 \ell(\b):= \bigl(\l /\rho\bigr)^{\frac{1}{\a+1+d}}(C_0 \b) ^{\frac{\a+1}{\a+1+d}}\,.
\en
 Then, for $\beta$ large \rosso{(in particular, when $\ell(\b)/\b\leq \min\{C_0,\e\}$)}, the law of \eqref{gandhi} 
with $\o$ sampled according to ${\rm PPP}[\rho,  \nu]$
 equals ${\rm PPP}[\l, \nu^+_{1, \a}] $ if $\nu\in \cvp$  and ${\rm PPP}[\l, \nu_{1, \a}] $  if $\nu\in \cv$. \end{Proposition}



%%%%%%%%%%%%%%%%%%%%%%
%\begin{Proposition}\label{paradiso_PPP} Suppose $\nu=\nu^+_{C_0,\a}$   (Case A) or $\nu=\nu_{C_0,\a}$  (Case B). Then, for    $0<\g \leq C_0$, it holds 
%\begin{itemize}
%\item[(i)]    $\nu_{\star, \g} = \nu^+_{1,\a}$ in Case A and   $\nu _{\star, \g}=\nu_{1,\a}$ in Case B; 
%\item[(ii)] $\nu[-\g, \g]= \g^{\a+1}/C_0^{\a+1}$, in particular   \eqref{vichinghi100} holds with $C_*:= C_0 ^{-\a-1}$\,.
%\end{itemize}
% Given $\rho, \l>0$ define  $\ell(\beta)$ as in \eqref{canditi105}  with  $C_*:= C_0 ^{-\a-1}$, i.e.
%\be\label{canditi200}
% \ell(\b):= \bigl(\l /\rho\bigr)^{\frac{1}{\a+1+d}}(C_0 \b) ^{\frac{\a+1}{\a+1+d}}\,.
%\en
% Then, for $\beta$ large, the law of \eqref{gandhi} 
%with $\o$ sampled according to ${\rm PPP}[\rho,  \nu]$
% equals ${\rm PPP}[\l, \nu^+_{1, \a}] $ in Case A  and ${\rm PPP}[\l, \nu_{1, \a}] $  in Case B. \end{Proposition}
%%%%%%%%%%%%%%%
The proof of the above proposition  is given in Section \ref{Trento}.





\subsection{Universality and strategy of the proofs: an overview}\label{lachiave}
 We now  explain more in detail how   Theorem \ref{paradiso} is related with universality of Mott's law, giving also an overview of the strategy in the derivation of Theorems \ref{parataUB} and \ref{parataLB}. 
 
 Let us first observe  that, given $E(\b)$ and $\z(\b)$, by Rayleigh monotonicity theorem we decrease the infinite volume conductivity $\s (\b) $ if we erase filaments in the MA resistor network,   in particular if we keep nodes $x$ with $E_x\in [-E(\b),E(\b)]$  and put a filament of conductance $c_{x,y}(\o,\b)$  for any survived nodes $x\not =y$ with  $c_{x,y}(\o,\b) \geq e^{-\z(\b)}$.   Note that 
  this thinned  resistor network is simply the weighted graph $\cG[\z(\b), \b] (\o _{E(\b)})$ with  weights  $c_{x,y}(\o,\b)$. We further decrease the infinite volume conductivity if we  replace the weights $c_{x,y}(\o,\b)$ by $e^{-\z(\b)}$.  We now restrict to the case $E(\b)=\z(\b)/\b$. Then, due to \orc{Lemma} \ref{fadeev_new} below,  the map $x\mapsto x/\z(\b)$ gives an isomorphism of the graph $\cG[\z(\b), \b] (\o _{E(\b) })$ with the graph $\cG[1,1](\o_{\z(\b), \b})$, where $\o_{\z(\b), \b}$ is defined as \eqref{gandhi} with $\z(\b)$ instead of $\ell(\b)$. Hence, by the above rescaling, we get a lower bound of $\s(\b)$ as  $e^{-\z(\b)}$ times a suitable functional of the graph $\cG[1,1](\o_{\z(\b), \b})$ (this is formalized in Lemma \ref{lasagne25} below).  This arguments hold in full generality, and do not restrict to $\nu$-randomized PPPs. 
  
  On the other hand, by using the variational formula \eqref{def_D} and Fact \ref{teo_eff_cond}, we can upper bound $D(\b)_{1,1}$ and therefore $\s(\b)$ in terms of  suitable expectations of  functions of the graph $\cG[\z(\b), \b] (\o)  $. See in particular Lemma \ref{santo}, where $C[\z,\b]$ is the cluster containing the origin in $\cG[\z(\b),\b) (\o)$. On the other hand, in Lemma \ref{serafico} we show that the above expectations can be written in terms of expectations of  functions of the graph $\cG[\z(\b), \b] (\o_{E(\b)})$.  Again these functions can be thought in terms of  $\cG[1,1](\o_{\z(\b), \b})$ by the above isomorphism and 
  again this derivation does not restrict to $\nu$-randomized PPPs. 
  

  The idea is then to apply the above machinery by taking $\z(\b)$ as $\ell(\b)$ of Theorem \ref{paradiso} with $\l <\l_c^*(\a)$ for the lower bound of $\s(\b)$ and with $\l>\l^*_c(\a)$ for the upper bound of $\s(\b)$, when the marked SPP $\bbP_\nu$  is as in Theorem \ref{paradiso} and $\l^*_c(\a)$ equals $\l^+_c(\a)$ or $\l_c(\a)$
 according to the limiting measure in Item (i) of the above theorem. 
  
   In both cases, Theorem \ref{paradiso} implies that the marked SPP $\o_{\z(\b), \b}$ (with $\o$ sampled by $\bbP_\nu$)
   converges to ${\rm PPP}[\l, \nu^+_{1, \a}]$  or  
${\rm PPP}[\l, \nu_{1, \a}]$. Hence somehow  the   graph $\cG[1,1](\o_{\z(\b), \b})$ 
suitably converges to the graph $\cG[1,1]$ read on the above marked PPP. This convergence and the above observations are in agreement with the experimented universality of Mott's law. 
 From a rigorous viewpoint the conclusion of the proof  of Mott  law in this general case would need some quantitate estimate on the speed of convergence in Theorem \ref{paradiso} (at least for some observables), which is trivial when one starts from  ${\rm PPP}[\rho,\nu]$ with $\nu\in \cv\cup \cvp$ (cf.~\orc{Proposition \ref{paradiso_PPP}}).  The resulting lower and upper bounds match better and better as soon as $\l$ approaches $\l_c^*(\a)$ (from below and from above, respectively).  
 



%
%
%\verde{
%\begin{Remark} Take $0\leq \g \leq C_0$. If $\nu=\nu^+_{C_0,\a}$  or $\nu=\nu_{C_0,\a}$, then   it holds   $\nu_{\star, \g} = \nu^+_{1,\a}$ and   $\nu _{\star, \g}=\nu_{1,\a}$, respectively (apply Lemma \ref{fissato} below with $\z:= \g \b$). Moreover, 
%in both cases, we have $\nu[-\g, \g]= \g^{\a+1}/C_0^{\a+1}$. Hence  condition \eqref{vichinghi100} holds with $C_*= C_0 ^{-\a-1}$. 
%\end{Remark}
%}
%
%
%%%%%%%%%
%
%
%\subsection{Arrenhius scaling of  conductivity  for  nearest-neighbor hopping} 
%
%\noindent
%$\bbZ^d$ \\
%Non-oriented edges $\bbE^d$. \\
%$\bigl( E_{b}\,:\, b \in \bbE^d\bigr)$ stationary ergodic field defined on some probability space $(\O, \cF,\bbP)$. $E_b \in \bbR$. Without loss, at cost to replace $\bbP$ with the law of the field, we think of $\bbP$ as a probability measure on $\O=\bbR^{\bbE^d}$. We write $\o=\bigl( E_{b}\,:\, b \in \bbE^d\bigr)$ for the generic element of $\O$. \\
%If $b=\{x,y\}$, then we write $E_{x,y}$ for $E_b$. Note that $E_{x,y}=E_{y,x}$.\\
%$c_{x,y}(\o,\b)= e^{-\b E_{x,y}(\o) } $\\
%Given $\ell \in \bbN$ even, we set $B_\ell:=[-\ell/2, \ell/2]^d\cap \bbZ^d$. 
%We also  consider the resistor network with nodes $B_\ell $ and electrical filaments between nearest--neighbor points $\{x,y\}\in \bbE^d$ with $x,y\in B_\ell$ having conductivity $c_{x,y}(\o, \b)$.  We write $\s_\ell (\o, \b)$ for the conductivity along the first direction, i.e. the total electrical flow along the first direction when setting the electrical potential equal to $0$ on $\{x\in B_\ell\,:\, x_1 =-\ell\}$ and equal to $1$ on $\{x\in B_\ell\,:\, x_1 = \ell\}$.
%
%
%%%%%%%%%%%%%%%%%%%%%%
%\begin{Definition}
% We define the effective diffusion matrix $D$ as the unique $d\times d$ symmetric matrix such that 
% \be\label{letargo1}
%  a \cdot Da =\inf _{ f\in L^\infty(\bbP) } \frac{1}{2} \sum_{e: |e|=1} \int d\bbP(\o) c_{0,e}(\o,\b)  \left
% (a\cdot e - \nabla_e  f (\o) 
%\right)^2\,,
%\en
%    where $\nabla_e  f (\o) := f(\t_e \o) - f(\o)$.
% \end{Definition}
% \club mettere condizione affinche' sia ben definita
% 
% \medskip
% 
%By  using 2-scale convergence as in \cite{Fhom2} one can prove the following:
% \begin{Fact}
%  Suppose that $\bbE[  c_{0,e_i}(\o,\b)^2 ]<+\infty$  for all $i=1,2,\dots, d$. Then, for  $\bbP$--a.a. $\o$, it holds
% \be \label{paloma} \lim_{\ell \to +\infty} \ell ^{2-d}\s_{\ell} (\o,\b) = D_{1,1}(\b) \,.
% \en
% \end{Fact}
% Physically $E_{x,y}$ corresponds to an energy barrier associated to the edge $\{x,y\}$.  It is therefore natural to think of the random variables $E_{x,y}$'s as uniformly lower bounded. In this case, the bound  $\bbE[  c_{0,e_i}(\o,\b)^2 ]<+\infty$ is satisfied for all $\b$.
% 
%  If the law of the the random field $\bigl( E_b\,:\, b \in \bbE^d\bigr)$ is invariant by coordinate exchange, then $D(\b)= d(\b) \bbI$ and  the r.h.s. of \eqref{paloma} can be replaced by the effective diffusion coefficient $d(\b)$.
%
%
%\begin{TheoremA}\label{teo3} Let $p_c\in (0,1)$ be the critical probability for Bernoulli bond percolation.
%Suppose that the field $\left(E_b\,:\, b \in \bbE^d\right)$ is given by i.i.d. random variables  bounded from below (i.e. there exists $c_0\in \bbR$ with $\bbP( E_b \geq c_0)=1$). Call $p(\theta):= \bbP( E_b \leq \theta)$, let $\theta^* := \inf\{ \theta \in \bbR\,:\, p(\theta) >p_c\}$ and  $\theta_* := \sup\{ \theta \in \bbR\,:\, p(\theta) <p_c\}$ (note that $\theta^* \geq \theta_*$).
%Then we have  
%\begin{align}
%& \liminf_{\b \to +\infty} \frac{1}{\b} \ln D_{1,1}(\b) \geq - \theta^*\,,\label{gelato1}\\
%& \limsup _{\b \to +\infty} \frac{1}{\b} \ln D_{1,1}(\b) \leq - \theta_*\,.\label{gelato2}
%\end{align}
%In particular, if there exists $\theta_c$ with $p(\theta)<p_c$ for $\theta <\theta_c$ and $p(\theta) >p_c$ for $\theta> \theta_c$ (e.g.  if $p(\theta_c)=p_c$ and $p$ is strictly monotone at $\theta_c$),
%  then \be\label{gelato3}
% \lim_{\b \to +\infty} \frac{1}{\b} \ln D_{1,1}(\b)= - \theta_c\,.
% \en
%\end{TheoremA}
%The proof of Theorem \ref{teo3} is given in Section \ref{proof_teo3} together with other results.
%
%\rosso{$p(\cdot)$ is weakly monotone and right--continuous.  If  $\theta_*<\theta$ then $p(\theta)\geq p_c$. If $\theta < \theta^*$, then  $p(\theta)\leq  p_c$. In particular, if  $\theta_*<\theta^*$, the $p(\theta)=p_c$ for all $\theta \in (\theta_*,\theta^*)$. By right--continuity $p(\theta)=p_c$ for all $\theta \in [\theta_*,\theta^*)$. It could be as well that $p(\theta)=p_c$ for all $\theta \in [\theta_*,\theta^*]$. For example, suppose that $E_b$ has value $1$ with probability $p_c$ and has value $2+y$ with probability $(1-p_c) \mathds{1}(y\in (0,1)) dy$. In this case $\theta_*=1$, $\theta ^*=2$ and $p(\theta)=p_c$ for all $\theta \in [1,2]$. \club }
%
%
%%%%%%%%




\section{Proof of Theorem \ref{paradiso} and Proposition \ref{paradiso_PPP}}\label{Tarvisio}


The results presented in this section hold for all dimensions $d\geq 1$.

\medskip

%%%%%%%%%%%%%%%%%%%%%%%%%%%%%%%%%%%
At the end of this section we give the proof of Theorem \ref{paradiso} and Proposition \ref{paradiso_PPP}. We first collect several preliminary results. To this aim 
recall Definition \ref{stella_stellina} of $\nu_{\star,\g}$.
\begin{Definition}[Constant $\nu(\g)$ and probability measure $\nu_\g$] \label{san_benedetto}
Let  $\nu$ be   a  probability measure on $\bbR$ and let  $\g\geq 0$.  We set 
\begin{equation}\label{alloro}
\nu(\g):=  \nu \bigl([ -\g , \g]\bigr)\in [0,1]\,.
 \end{equation}
 If $0$ belongs to the support of $\nu$, given $\g>0$ we define the  probability measure   $\nu_{ \g} $ as
 \be
 \nu _{\g} ( B ) =\frac{\nu ( B \cap [-\g, \g])}{\nu([-\g, \g])} \,, \qquad B \in \cB (\bbR) \,.\label{sette1}
 \en
\end{Definition}
Note that the  probability measures  $\nu_{ \g} $ and  $\nu_{\star,\g}$ have  support contained in   $[ -\g , \g]$ and $[-1,1]$, respectively. Moreover, it holds 
\begin{align}
& \nu _{\g} ( B )= P\big(\,   X \in B \,\big{|}\, |X |\leq \g \,\big)\,,\label{sette1_bis}\\
&   \nu _{\star,\g} ( B )= P\big(\,  (X/\g) \in B \,\big{|}\, |X/\g |\leq 1\,\big) =\nu_\g(\g B) \,,\label{sette2}
 \end{align}
where $X$ is a random variable with distribution  $\nu$ and $B\in \cB(\bbR)$.




\begin{Definition}[$ \widehat\o_\g$, $\o_\g$] \label{porto_dascoli}
Given $\o \in \O$ and $\g >0$ we define 
\begin{align}
& \widehat\o_\g:=\{x\in \widehat \o\,:\, |E_x| \leq \g\}\,,\label{flautino1}\\
& \o_\g :=\{ (x, E_x)\,:\, x\in \widehat{\o} \,, \; |E_x| \leq \g\}\,. \label{flautino2}
\end{align}
\end{Definition}
Note that $\widehat{\,\o_\g\, } = {\widehat{\o}}_\g$. 

Recall Definitions \ref{def_randomization}  and \ref{def_thinning} of $\nu$--randomization and $p$--thinning, respectively. The above  \rosso{configurations}   $\widehat\o_\g$ and $\o_\g$ are obtained from  $\widehat\o$ and $\o$ by a kind of energy-based thinning.  We now clarify the relation between the above energy-based thinning and the random one introduced in Definition  \ref{def_thinning}. 
%%%%%%%%%%%%%%%%%%%%%%%%%%%%%%%%%%%%%%%%%%%%%%%%%%%%%%%%%%%%%%%%%%%%%%%%%%%%%%%%%%%%%%%%%%%%%%%%%%%%%%%%%%%%%%%%%%%%%%%%%%%%%%%%%%%%%%%%%%%%%%%%%%%%%%%%%%%%%%%%

\begin{Lemma}\label{lemma_2021} Let $\g>0$,  let $\hat\bbP$ be  the law of a SPP  on $\bbR^d$   and let $\nu$ be a probability measure on $\bbR$.
If $\o$ is sampled according to $\hat \bbP_\nu$, then $\o_\g $ has law $(\hat \bbP_p)_{\nu_\g}$ where $p:=\nu(\g)\rosso{=\nu([-\g,\g])}$.\end{Lemma}
%%%%%%%%%%%%%%%%%%%%%%%%%%%%%%%%%%%%%%%%%%%%%%%%%%%%%%%%%%%%%%%%%%%%%%%%%%%%%%%%%%%%%%%%%%%%%%%%%%%%%%%%%%%%%%%%%%%%%%%%%%%%%%%%%%%%%%%%%%%%%%%%%%%%%%%%%%%%%%%%%%%%%%%%%%%%%%%%
\begin{proof} 
% The claim is very intuitive and the proof is a trivial exercise, hence we only sketch it.
By  \cite[Lemma~12.1]{Kal} we just need to prove that, 
given  any measurable function $f: \bbR^d\times\bbR\to [0,+\infty)$,  the random variable $\sum_{x\in \widehat{\o}_\g }f(x, E_x)$ with $\o$ sampled according to $\hat \bbP_\nu$ has the same law of the random variable $\sum_{x\in \widehat{\o} }f(x, E_x)$ with 
 $\o$ sampled according to $(\hat \bbP_p)_{\nu_\g}$.  \rosso{Moving to the Laplace transform}, it is enough to show that,  for any $f$ as above, it holds %the corresponding Laplace transforms coincide:
% 
% 
% On the other hand, by  \cite[Section~\rosso{XIII.4}]{Fel2}, two  random variables $X$ and $Y$ with values in $[0,+\infty]$ have the same law if and only if they have the same Laplace transform, i.e. $E[e^{-\g X}]=E[e^{-\g Y}]$ for all $\g \geq 0$. As a byproduct of  the above two general results, we just need to prove that 
 \be\label{gente} \int_\O d \hat \bbP_\nu(\o)
 e^{- \sum_{x\in \widehat \o_\g} f(x,E_x)}  =\int _\O d (\hat \bbP_p)_{\nu_\g} (\o) e^{-\sum_{x\in \widehat \o} f(x,E_x)} \,.
 \en
%for any  measurable function $f:\bbR^d\times \bbR\to [0,+\infty)$.  
\rosso{ Since $e^{- \sum_{x\in \widehat \o_\g} f(x,E_x)}=\prod_{x\in \widehat \o} e^{- f(x,E_x) \mathds{1}( |E_x| \leq \g )  }$}, it is easy to rewrite \eqref{gente}  as \begin{multline*}
\int _{\hat \O} d \hat \bbP(\hat \o) \prod_{x\in \hat \o} \Big[ \nu( [ -\g, \g]^c)+\int_{[-\g, \g]} \nu( dE)  e^{- f(x,E) }\Big]\\ = \int _{\hat \O} d \hat \bbP(\hat \o) \prod_{x\in \hat \o}     \Big[ 1-p+ p\int_{[-\g, \g]} \nu_\g ( dE)  e^{- f(x,E) }   \Big]  \,.
\end{multline*}
The claim follows from the identity  $p=\nu([-\g,\g])$ and the definition of $\nu_\g$.
%
%To this aim we first observe  that 
%\be\label{int1}
%\begin{split}
%& \int_\O d \hat \bbP_\nu(\o)
% e^{- \sum_{x\in \widehat \o_\g} f(x,E_x)}  =
% \int_\O d \hat \bbP_\nu(\o)
% \prod_{x\in \widehat \o} e^{- f(x,E_x) \mathds{1}( |E_x| \leq \g )  }\\
% &=
% \int _{\hat \O} d \hat \bbP(\hat \o) \prod_{x\in \hat \o} \Big[ \nu( [ -\g, \g]^c)+\int_{[-\g, \g]} \nu( dE)  e^{- f(x,E) }\Big]\,.
% \end{split}
%\en
%Let $\mathfrak{m}:= p \d_1 +(1-p) \d_0$. We recall that the $p$--thinning of a SPP  with law $\hat \bbP$ can be obtained by marking the points $x\in \hat \o $   by independent Bernoulli  random variables $t_{x}$  with parameter $p$ (independently from all the rest) and afterwards keeping only the points $x$ such that  $t_x=1$. Hence
%we have 
%\begin{equation*}
%\begin{split}
%&\int _\O d (\hat \bbP_p)_{\nu_\g} (\o) e^{-\sum_{x\in \widehat \o} f(x,E_x)} =
% \int _\O d (\hat \bbP_p)_{\nu_\g} (\o)
% \prod_{x\in \widehat \o} e^{- f(x,E_x)  }  \\
% &=  
% \int _{\hat \O}  d \hat \bbP_p(\hat \o) \prod_{x\in \hat \o} \Big[ \int_{[-\g, \g]} \nu_\g( dE)  e^{- f(x,E) } \Big]  \\
% &=
% \int _{\hat \O} d \hat \bbP_p(\hat \o) \prod_{x\in \hat \o} \Big[p^{-1} \int_{[-\g, \g]} \nu( dE)  e^{- f(x,E) } \Big]\\
% &=
% \int _{\hat \O} d \hat \bbP(\hat \o) \prod_{x\in \hat \o}  \int _{\{0,1\}}d \mathfrak{m} (t_x) \Big\{
%\mathds{1}(t_x=0) +\frac{\mathds{1}(t_x=1)}{p} \int_{[-\g, \g]} \nu( dE)  e^{- f(x,E) } \Big \} \\
%&=
% \int _{\hat \O} d \hat \bbP(\hat \o) \prod_{x\in \hat \o}     \Big[ 1-p+ \int_{[-\g, \g]} \nu( dE)  e^{- f(x,E) }   \Big] \,.
%  \end{split}
%\end{equation*}
%As $\nu( [ -\g, \g]^c)=1-p$, 
% by comparing the last expression above with the last expression  in   \eqref{int1} we get  \eqref{gente}.
\end{proof}
%%%%%%%%%%%%%%%%%%%%%%%%%%%%%%%%%%%%%%%%%%%%%%%%%%%%%%%%%%%%%%%%%%%%%%%%%%%%%%%%%%%%%%%%%%%%%%%%%%%%%%%%%%%%%%%%%%%%%%%%%%%%%%%%%%%%%%%%%%%%%%%%%%%%%%%%%%%%%%%%%%%%%%%%%%%%%%%%%%%%%%%%%%%%%%%%%%%%%%%%%%%%%%%%%%%%%%%%%%
%%%%%%%%%%%%%%%%%%%%%%%%%%%%%%%%%%%%%%%%%%
\begin{Lemma}\label{weak_conv}
Consider a sequence of SPPs on $\bbR^d$  with laws $\pn$ and a sequence of probability measures $\mu_n$ on $\bbR$. Suppose that there exists a SPP   with law $\hat\bbP$ and a  probability measure $\mu$ on $\bbR$ such that  
$ \pn \,\Rightarrow \, \hat \bbP$  and $\mu_n \,\Rightarrow \, \mu$ as $n\to +\infty$.
Suppose also that $\int_{\hat \O}  d \hat \bbP(\hat \o ) \hat \o (\partial A)=0$ for any box $A\subset \bbR^d$.
Then $
 ( \pn)_{\mu_n} \,\Rightarrow \, \hat \bbP_\mu$ as $n\to +\infty$.
\end{Lemma}


Above  the symbol $\Rightarrow$ refers to weak convergence of probability measures \rosso{\cite{Bi}. When $\mu_n= \mu$ for all $ n$, the above result follows from \cite[Exercise~15,\,Chp.~16]{Kal}.}


 \begin{Remark}\label{condimento}The condition $\int_{\hat \O}  d \hat \bbP(\hat \o ) \hat \o (\partial A)=0$  for boxes $A\subset \bbR^d$  is satisfied in particular by any stationary $\bbP$ with finite intensity, as the Lebesgue measure gives zero mass to $(d-1)$--dimensional sets. 
 \end{Remark}

% We also point out that Lemma \ref{weak_conv} with $\mu_n\equiv mu$ and without any condition on $\partial A$ corresponds to \cite[Exercise~..]{Kal}.
%%%%%%%%%%%%%%%%%
\begin{proof} Below we identify the configurations $\hat\o\in \hat\O$ and $\o\in \O$ with the atomic measures $\sum_{x\in \hat\o} \d_x$ and $\sum_{x\in \widehat \o} \d_{(x,E_x)}$ respectively, keeping the same names $\hat\o$ and $ \o$. In particular, $\o[f]=\sum_{x\in \widehat{\o} }f(x,E_x)$. Moreover, we denote by $f\in C^+_c(\bbR^{d+1})$ the family of 
  continuous functions $f\geq 0$ on $\bbR^{d+1}$ with compact support.


 A marked SPP on $\bbR^d$ is in particular a SPP  on $\bbR^{d+1}$, thus allowing to apply \cite{Kal}[Thm.~16.16]. 
 As a byproduct  with the properties of Laplace transform,
 %(cf. \cite[Section~\rosso{XIII.4}]{Fel2}), 
  to prove that $(\pn)_{\mu_n} \,\Rightarrow \,  \hat \bbP_\mu$ it is enough to show \rosso{for any $f\in C_c^+(\bbR^{d+1})$} that 
 \be\label{prati} \int d(\pn)_{\mu_n} (\o) e^{- \o [f]} \to  \int d\hat \bbP _{\mu} (\o) e^{- \o [f]}  \;\; \text{ as }n\to+\infty\,.
 \en
 
 
  
 \rosso{Let 
$\cG$ be the class of functions  
$g$ of the 
 form $g=\sum_{i=1}^N c_i \mathds{1}_{A_i \times B_i}$, where all  $c_i$'s are positive  numbers, 
 all $A_i\times B_i$ are pairwise disjoint, 
  all $A_i $  are boxes of the form $\prod_{k=1}^d [a_k,b_k)$ and all $B_i$ are intervals of $\bbR$ whose extremes have zero $\mu$-measure. For any $\e>0$ one can find $f^{\pm }_\e\in \cG$  with support included in the neighborhood of the support of $f$ with radius $1$, such that $f^-_\e\leq f \leq f^+_\e$ and  $\|f- f^\pm _\e\|_\infty \to 0$ as $\e\da 0$. 
  Hence, by dominated convergence, $\int d \hat \bbP_{\mu} (\o) e^{- \o [f^\pm_\e]}\to \int d \hat \bbP_{\mu} (\o) e^{- \o [f]}$ as $\e \da 0$, and then one 
  easily derive  \eqref{prati} with    $f\in C_c^+(\bbR^{d+1}) $  once getting it for any
  $f\in \cG$.
 }
 
  
 
% We denote by  $\G$ the countable set $ \{x\in \bbR\,:\, \mu(\{x\})>0\}$.
%  For each $\e>0$ we fix a partition $\cS_\e$ of $\bbR^d\times \bbR$  in boxes of the form $A\times B$, where $A\subset \bbR^d$ is \rosso{a} box of the form $\prod _{i=1}^d [k_i \e ,k_i \e +\e)$  with $k_i \in \bbZ$ and $B \subset  \bbR$ is an  interval of the form $[a,b)$ with $a,b \not \in \G$ and $b-a\leq \e$. For each set $U\subset \bbR^d \times \bbR$ we denote by  $c^+(U) $ and $c^-(U) $ the supremum and the infimum of $f$ on $U$ respectively. Finally we set
% $ f_\e^\pm (x,E):= \sum_{A\times B \in \cS_\e} c^\pm (A\times B) \mathds{1}_{A\times B}(x,E)$.
% Then  $f_\e^\pm$  has compact support and is a finite linear combination of characteristic functions of sets in $\cS_\e$. Moreover,  $ 0\leq f_\e^- \leq f\leq f_\e^+$ and   $\lim_{\e \da 0} \|f_\e^\pm -f \|_\infty = 0$.
%\rosso{We first  suppose that we have already proved
%\eqref{prati} with $f^\pm_\e$ instead of $f$ and show how to conclude}. 
% Since $ e^{- \o [f^+_\e]}\leq  e^{- \o [f]} \leq  e^{- \o [f^-_\e]}$, we would get 
% \begin{equation}\label{altino}
% \begin{split}
% \int d\hat \bbP_{\mu} (\o) e^{- \o [f^+_\e]}  & \leq \liminf_{n\to +\infty} \int d(\pn)_{\mu_n} (\o) e^{- \o [f]}\\
% & \leq  \limsup_{n\to +\infty} \int d(\pn)_{\mu_n} (\o) e^{- \o [f]} \leq\int d \hat \bbP_{\mu} (\o) e^{- \o [f^-_\e]}\,.
% \end{split}
% \end{equation}
% Now we take the limit $\e\da 0$.  Denote by $\D$ the neighborhood of radius $1$ of the support of $f$. For $\e$ small enough  $f^+_\e$ has support in $\D$.   Hence $e^{- \o [f^+_\e]}= e^{-\sum_{z\in \o \cap \D} f^+_\e(z)}  $ and therefore  $e^{- \o [f^+_\e]}$ converges to $e^{- \o [f]}$ as $\e\da 0$ for each $\o$ (use that $\lim_{\e \da 0} \|f_\e^\pm -f \|_\infty = 0$).
%  Since $e^{- \o [f^+_\e]}$ is upper bounded by $1$, by dominated convergence the first term in \eqref{altino} converges to $\int d\hat \bbP_{\mu} (\o) e^{- \o [f]}$ as $\e \da 0$. The same holds for the last term in \eqref{altino} by similar arguments.  We then derive from \eqref{altino} that \eqref{prati} must be satisfied.
% 
% 
% Up to now we have shown that it enough to prove \eqref{prati} with $f$ of the form $f=\sum_{i=1}^N c_i \mathds{1}_{A_i \times B_i}$, where all  $c_i$'s are positive  numbers, all   $A_i \times B_i$'s belong to  $ \cS_\e$ for some $\e>0$ and are pairwise disjoint.
 %  where $\{A_i\times B_i\}_{i=1}^N$ are disjoint sets.
 %Note that there can  be some repetitions in  $A_1,A_2,\dots, A_N$. 

Let us now   prove \eqref{prati}
for any given $f=\sum_{i=1}^N c_i \mathds{1}_{A_i \times B_i}$ in $\cG$. At cost to split each  $A_i$ into smaller boxes, 
  we assume that $A_1,A_2, \dots, A_k$ are pairwise disjoint, while $A_{k+1},\dots, A_N$ equal some sets between $A_1,A_2, \dots, A_k$. Hence 
  we can rearrange the sum to write $f$ in  the form 
% \[
% f(x,E)= \sum_{i=1}^k \mathds{1}_{A_i}(x)   \sum _{j\in J_i} c_{i,j} \mathds{1} _{D_j}(E) = : \sum_{i=1}^k \mathds{1}_{A_i}(x)  g_i(E)\,,
% \] 
$
 f(x,E)=\sum_{i=1}^k \mathds{1}_{A_i}(x)  g_i(E)$, where $ g_i (E) := \sum _{\substack{j: 1\leq j \leq N\\
 A_j=A_i}} c_j \mathds{1}_{B_j} (E)$.
 %where  $D_j$ is an interval with no extreme  in $\G$, $J_i$ is a finite index set and the intervals $D_j$ with $j\in J_i$  are pairwise  disjoint.  
Then  we have 
\be\label{sole}
\begin{split}
 \int d(\pn)_{\mu_n} (\o) e^{- \o [f]} 
 %= \int d \hat\bbP_n (\hat \o)\prod _{x\in \hat \o } \mu_n[ e^{- f(x, \cdot) }] 
% & =   \int d \pn (\hat \o)\exp\Big\{ \sum_{x\in \hat \o}  \ln \mu_n[ e^{-f(x, \cdot)}] \Big\}\\
& =
 \int d \pn (\hat \o)\exp\Big\{ \sum_{i=1}^k \hat \o (A_i)  \ln \mu_n[ e^{-g_i}] \Big\}\,.
%\int d \hat\bbP_n (\hat \o)\exp\Big\{ \hat \o \bigl[\ln \mu_n[ e^{- h g}] \bigl] \Big\}\,,
\end{split}
\en
For each $i=1,\dots, k$, \rosso{$g_i:\bbR\to \bbR$ is a piecewise-constant function, with a finite family of   discontinuity points having zero $\mu$-mass. Hence   $\mu_n [ e^{-g_i}]\to \mu[e^{-g_i}]$  as $n\to \infty$.} 
% the discontinuity points of    $h_i :=e^{-g_i}:\bbR\to (0,1]$ belong to 
%$\rosso{D}:=\cup _{j=1}^N \partial B_j$. Since  $\mu(\rosso{D})=0$ by  definition of $\cS_\e$ and since $\mu_n\Rightarrow \mu$,  the mapping theorem (cf.~\cite[Thm.~2.7]{Bi}) implies that  $\mu_n \circ h_i^{-1} \Rightarrow \mu \circ h_i^{-1} $ as probability measures on $(0,1]$. Since the identity map on $(0,1]$ is bounded and continuous, we conclude that 
% $\mu_n [ e^{-g_i}]\to \mu[e^{-g_i}]$ \rosso{as $n\to \infty$}. 
Hence, fixed $\d>0$, for $n$ large enough we have
%$ \ln \mu[ e^{-g_i}] -\d \leq   \ln \mu_n[ e^{-g_i}] \leq  \{ \ln \mu[ e^{-g_i}] +\d \}_-\leq 0 $, 
%where $\{a\}_-:=\min\{a,0\}$. This observation and \eqref{sole} imply that 
\begin{multline}\label{fulmini}
 \int d \pn (\hat \o)\exp\Big\{ \sum_{i=1}^k \hat \o (A_i)  (\ln \mu [ e^{-g_i}] -\d) \Big\}
\leq 
\int d(\pn)_{\mu_n} (\o) e^{- \o [f]} 
\\ \leq 
 \int d \pn (\hat \o)\exp\Big\{ \sum_{i=1}^k \hat \o (A_i)  \{ \ln \mu[ e^{-g_i}] +\d \}_- \Big\}\,.
\end{multline}
For any  box $A_i$ it holds  $\hat \o (\partial A_i)=0$ for $\hat \bbP$--a.a.~$\hat \o$ as 
  $\int d \hat \bbP(\hat \o ) \hat \o (\partial A_i)=0$ by assumption.  Since in addition the sets $A_i$ are relatively compact, 
 we can apply \cite[Thm.~16.16]{Kal} getting that  the vector $(\hat \o (A_1), \hat \o (A_2), \dots , \hat \o (A_k))$ with $\hat \o$ sampled according to $\pn$  converges in distribution to the same vector with $\hat \o$  sampled now according to $\hat\bbP$. \rosso{This implies the convergence of expectations of bounded continuous functions applied to the random  vectors. Hence we get}
% \rosso{Since the function 
% %Given non-positive numbers $b_1,b_2,\dots, b_k$,  the function
%  $\bbR^k \ni z \mapsto e^{  \sum_{i=1}^k b_i z_i}  \in \bbR$ is continuous and bounded when $b_i\leq 0$ for all $i$, 
%  % In particular,  since the coefficients $\ln \mu [ e^{-g_i}] -\d$ and $ \{ \ln \mu[ e^{-g_i}] +\d \}_-$ are non-positive, by definition of weak convergence of random vectors in $\bbR^k$ we get
  \begin{align}
 & \lim_{n\to +\infty} \text{l.h.s. of \eqref{fulmini}}= \int d \hat\bbP (\hat \o)\exp\big\{ \sum_{i=1}^k \hat \o (A_i)  (\ln \mu [ e^{-g_i}] -\d) \big\}\,, \label{merenda1}\\
  & \lim_{n\to +\infty} \text{r.h.s. of \eqref{fulmini}}= \int d \hat\bbP (\hat \o)\exp\big\{ \sum_{i=1}^k \hat \o (A_i)  \{ \ln \mu[ e^{-g_i}] +\d \}_- \big\}\,.\label{merenda2}
\end{align}
%The integrands in the r.h.s. of \eqref{merenda1} and \eqref{merenda2} are upper bounded by $1$ and both converge to $\exp\Big\{ \sum_{i=1}^k \hat \o (A_i)  \ln \mu [ e^{-g_i}] \Big\} $ as $\d\da 0$. Hence, 
By dominated convergence,  the limits as $\d\da 0$ of \eqref{merenda1} and \eqref{merenda2} both equal
\[ \int d \hat\bbP (\hat \o)\exp\big\{ \sum_{i=1}^k \hat \o (A_i)   \ln \mu[ e^{-g_i}]  \big\}=\int d\hat \bbP_{\mu} (\o) e^{- \o [f]} \,.
\]
 By the above observations, taking the limit $n\to +\infty$ and afterwards the limit  $\d \da 0$ in \eqref{fulmini}, we get \eqref{prati}.
% $$ ****$$
% Let $f(x,E):=h(x) g(E)$, $h$ and $g$ being  nonnegative continuous functions with on $\bbR^d$ and $\bbR$ respectively, $h$ having compact support.
% Below we will write $\D$ for the neighborhood with unit  radius of the compact support of $h$. 
% We identify the configurations $\hat\o\in \hat\O$ and $\o\in \O$ with the atomic measures $\sum_{x\in \hat\o} \d_x$ and $\sum_{x\in \hat \o} \d_{(x,E_x)}$ respectively, keeping the same names $\hat\o, \o$. In particular, $\o[f]=\sum_{x\in \widehat{\o} }f(x,E_x)$.
%As in  \cite[Lemma 12.2]{Kal}--(iii) we have 
%\be
%\begin{split}
%& \int d(\hat \bbP_n)_{\mu_n} (\o) e^{- \o [f]} 
%= \int d \hat\bbP_n (\hat \o)\prod _{x\in \hat \o } \mu_n[ e^{- h(x) g }] =\\
%&   \int d \hat\bbP_n (\hat \o)\exp\Big\{ \sum_{x\in \hat \o}  \ln \mu_n[ e^{- h(x) g}] \Big\}=
%\int d \hat\bbP_n (\hat \o)\exp\Big\{ \hat \o \bigl[\ln \mu_n[ e^{- h g}] \bigl] \Big\}\,,
%\end{split}
%\en
%where the compact forms $\mu_n[ e^{- h(x) g }] $ and  $\hat \o \bigl[\ln \mu_n[ e^{- h g}]\bigr]$ mean respectively $\int _{\bbR} d\mu_n(E) e^{-h(x) g(E)} $ and  $\sum_{x\in \hat \o} \ln \mu_n[ e^{- h(x) g}]$.
%
%
%%In the above notation, we think of $\hat \o$ as $\sum_i \d_{x_i}$. Moreover, when writing $\hat \o \bigl[\ln \mu_n[ e^{- h g}] \bigl]$, we think that $\mu_n$ integrates among $E$ appearing in $g(E)$, while $\hat \o$ sums among the points $x_i$ appearing in $h(x_i)$.
%
%For each $x\in \bbR^d$ we have $\mu_n[ e^{- h(x) g(\cdot)}]  \to \mu[ e^{- h(x) g(\cdot)}]$ as $E\mapsto e^{- h(x) g(E)}$ is a bounded continuous function and $\mu_n \Rightarrow \mu$. Since
%$h$ is continuous with bounded support, given $\e>0$ we can find 
%disjoint  boxes  $A_1,A_2, \dots , A_N\subset \D$   of the form $\prod_{i=1}^d[c_i,d_i)$
%and values $a_1, \dots, a_N\geq \e $ such that   $h(x)=0$ for all $x\not \in \cup_{k=1}^N A_k$ and 
%$h(x) \in [a_k-\e , a_k +\e]$ for all $x\in A_k$. 
%%As $h,g\geq 0$,  given $x\in A_k$  we have 
%%\begin{multline}\label{stupidino}
%%0\leq  \{h(x)-2\e\}_+ g(E) \leq   (a_k-\e) g(E) \leq  h(x) g(E) \\\leq ( a_k+\e )g(E) \leq (h(x)+2\e) g(E) \,.
%%$\end{multline}
%Since $\mu_n \Rightarrow \mu$,  we have $\lim_{n\to \infty}  \ln \mu_n[ e^{- (a_k\pm \e)  g}]=
% \ln \mu_n[ e^{- (a_k\pm \e)  g}]\leq 0 $. %Note that also $\ln \mu_n[ e^{- (a_k\pm \e)  g}]$ is nonnegative. 
% Hence  
%we can  fix $n_0$ large enough such that for any $n\geq n_0$ and any $k\in \{1,\dots,N\}$ it holds
%\[
% \ln \mu[ e^{- (a_k+\e)  g}]-\e\leq 
% \ln \mu_n[ e^{- (a_k+\e)  g}] \leq  \ln \mu_n[ e^{- (a_k-\e)  g}] \leq \ln \mu[ e^{- (a_k-\e)  g}] +\e\,.
% \]
% As $a_k \pm \e \geq 0$ and $g\geq 0$, the first above three terms  are upper bounded by $0$. Hence, we can modify the last one getting
%  \be\label{lenticchia}
% \ln \mu[ e^{- (a_k+\e)  g}]-\e\leq 
% \ln \mu_n[ e^{- (a_k+\e)  g}] \leq  \ln \mu_n[ e^{- (a_k-\e)  g}] \leq  \{ \ln \mu[ e^{- (a_k-\e)  g}] +\e\}_-\,.
% \en
% Hence, for $n\geq n_0$, using the last above bound we have 
%\be\label{wurstel}
%\begin{split}
%& \int d \hat\bbP_n (\hat \o)\exp\Big\{ \hat \o \bigl[\ln \mu_n[ e^{- h g}] \bigl] \Big\}\\
%& \leq 
%\int d \hat\bbP_n (\hat \o)\exp\Big\{  \sum_{k=1}^N 
%\hat \o (A_k)  \ln \mu_n[ e^{- (a_k-\e)  g}] \Big\}\\
%& \leq 
%\int d \hat\bbP_n (\hat \o)\exp\Big\{  \sum_{k=1}^N  \hat \o(A_k)
%  \bigl\{  \ln \mu[ e^{- (a_k-\e)  g}] +\e\bigr\}_-   \Big\}\,.
%\end{split}
% \en
%For any  box $A_k$ it holds  $\hat \o (A_k)=0$ for $\hat \bbP$--a.a.~$\hat \o$ as 
%  $\int d \hat \bbP(\hat \o ) \hat \o (\partial A_k)=0$ by assumption.  Since in addition the sets $A_k$ are relatively compact, 
% we can apply \cite[Thm.~16.16]{Kal} getting that  the vector $(\hat \o (A_1), \hat \o (A_2), \dots , \hat \o (A_n))$ with $\hat \o$ sampled according to $\hat \bbP_n$  converges in distribution to the same vector with $\hat \o$  sampled now according to $\hat\bbP$. The above weak convergence implies the convergence of the associated  moment generating functions. \rosso{In particular, since all numbers $\{  \ln \mu[ e^{- (a_k-\e)  g}] +\e\}_- $ are non-positive, from \eqref{wurstel}}  we get 
% \be\label{wurstel1}
%\begin{split}
%& \limsup_{n\to+ \infty} \int d \hat\bbP_n (\hat \o)\exp\Big\{ \hat \o \bigl[\ln \mu_n[ e^{- h g}] \bigl] \Big\} \\
%&\leq  \int d \hat\bbP  (\hat \o)\exp\Big\{  \sum_{k=1}^N 
%\hat\o(A_k) \{  \ln \mu[ e^{- (a_k-\e)  g}]+\e\}_-   \Big\}\,.
%\end{split}
%\en 
%We recall that  $\D$ is the $1$--neighborhood of the support of $h$ and all sets $A_k$ are disjoint and included in $\D$. 
%\rosso{Let now $x\in A_k$. As $g\geq 0$ and $h(x) \leq a_k+\e$, we have  $(h(x)-2\e) g(E) \leq   (a_k-\e) g(E)$}. This implies that 
%\begin{equation*}
%\begin{split}
% \ln \mu[ e^{- (a_k-\e)  g}]  
%\leq \ln \mu[ e^{- (h(x)-2\e) g}] \leq  \ln \mu[ e^{- h(x) g+ 2 \e \| g\|_\infty }] =  \ln \mu[ e^{- h(x) g }]+ 2 \e \| g\|_\infty  \,.
%\end{split}
%\end{equation*}
%The first term is also upper bounded by $0$.
%Hence from \eqref{wurstel1} we get 
% \be\label{wurstel2}
%\begin{split}
%& \limsup_{n\to +\infty} \int d \hat\bbP_n (\hat \o)\exp\big\{ \hat \o \bigl[\ln \mu_n[ e^{- h g}] \bigl] \big\} \\
%&\leq  \int d \hat\bbP  (\hat \o)\exp\big\{ \hat \o \bigl[  \{  \ln \mu[ e^{- h  g}] +\e (2 \| g\|_\infty +1) {\mathds{1}}_{\D}  \}_-\bigl]   \big\} \,.
%\end{split}
%\en 
%The exponential  in  the r.h.s. of \eqref{wurstel2}  is upper bounded by $1$.  Hence, by  taking  the limit $\e \da 0$ and  applying stochastic domination, the $\e$--term in the r.h.s. of \eqref{wurstel2} disappears:
%\be\label{wurstel100}
%\limsup_{n\uparrow +\infty} \int d \hat\bbP_n (\hat \o)\exp\big\{ \hat \o \bigl[\ln \mu_n[ e^{- h g}] \bigl] \big\} \leq  \int d \hat\bbP  (\hat \o)\exp\big\{ \hat \o \bigl[   \ln \mu[ e^{- h  g}]\bigl]   \big\} \,.
%\en
%We can argue in the similar way to deal with the liminf instead. To this aim one has to use the first bound in \eqref{lenticchia} to get   for $n\geq n_0$ that 
%\begin{equation*}
%\begin{split}
%&  \int d \hat\bbP_n (\hat \o)\exp\Big\{ \hat \o \bigl[\ln \mu_n[ e^{- h g}] \bigl] \Big\}\geq \\
%& 
%\int d \hat\bbP_n (\hat \o)\exp\Big\{  \sum_{k=1}^N  \hat \o(A_k)
%  (  \ln \mu[ e^{- (a_k+\e)  g}] -\e)  \Big\}\,.
% \end{split}
% \end{equation*}
% Since the coefficients $\ln \mu[ e^{- (a_k+\e)  g}] -\e$ are negative, we can again apply the convergence of Laplace transforms as before getting that 
%\begin{equation}\label{wurstel3}
%\begin{split}
%& \liminf_{n \to +\infty}  \int d \hat\bbP_n (\hat \o)\exp\Big\{ \hat \o \bigl[\ln \mu_n[ e^{- h g}] \bigl] \Big\}
% \\
%&\geq  \int d \hat\bbP (\hat \o)\exp\Big\{  \sum_{k=1}^N  \hat \o(A_k)
%  (  \ln \mu[ e^{- (a_k+\e)  g}] -\e)  \Big\}
%\,.
%\end{split}
% \end{equation}
%Given $x\in A_k$ we have $ a_k -\e \leq h(x)$, i.e. $a_k+\e \leq h(x)+2\e$. Since in addition $g\geq 0$  we can estimate $  \ln \mu[ e^{- (a_k+\e)  g}] \geq \ln \mu[e^{-(h(x)+2 \e) g}]\geq  \ln \mu[e^{-h(x)g-2 \e\| g\|_\infty}]=-2 \e\| g\|_\infty + \ln \mu[e^{-h(x)g}]$. This allow to lower bound the r.h.s. of \eqref{wurstel3} with
%\be \label{wurstel4}
%\int d \hat\bbP (\hat \o)\exp\Big\{ 
%\hat \o \bigl[\ln \mu_n[ e^{- h g}] - \e( 1+ 2 \| g\|_\infty \mathds{1}_\D) \bigl]\Big\}
%\en
%Again one has to apply the dominated convergence theorem when taking the limit $\e\da 0$ in \eqref{wurstel3} and \eqref{wurstel4}. By combining the resulting estimate with \eqref{wurstel100} we get 
%\begin{equation*}
%\begin{split}
%\lim_{n\uparrow \infty} \int d(\hat \bbP_n)_{\mu_n} (\o) e^{- \o [f]} &   =
%\lim_{n\uparrow \infty} \int d \hat\bbP_n (\hat \o)\exp\big\{ \hat \o \bigl[\ln \mu_n[ e^{- h g}] \bigl] \big\}\\
%& =  \int d \hat\bbP  (\hat \o)\exp\big\{ \hat \o \bigl[    \ln \mu[ e^{- h  g}] \bigl]   \big\} = \int d(\hat  \bbP)_{\mu} (\o) e^{- \o [f]}
% \end{split}
% \end{equation*}
% for $f(x,E)= h(x) g(E)$, for $h,g$ bounded  continuous nonnegative functions $\bbR^d$ and $\bbR$, $h$ having bounded support.
%   \club bisogna fare vedere che queste funzioni test sono abbastanza (see the covering semiring p.139 DV2)
  \end{proof}

\bigskip


\begin{Definition}[Homothety $\Psi_z$] \label{def_zaino}
Given $z >0$ we define $\Psi_z: \bbR^d\to \bbR^d$ as the map $\Psi_z (x) := x/z$.
\end{Definition}



If we have a SPP $\hat{\xi}$ on $\bbR^d$ with law $\hat\bbP$,  we write $\hat \bbP\circ \Psi_z ^{-1}$  for the law of $\psi_z(\hat{\xi})=\{x/z\,:\, x\in \hat\xi\}$:
$\bigl(\hat \bbP\circ \Psi_z ^{-1}\bigr)(B):= \hat \bbP\bigl( \Psi_z^{-1} (B)\bigr) $ for all $B \in \cB (\hat\O)$.

%c
%%%%%%%%%%%%%%%%%%%%%%%%%%%%%%%%%%%%%%%%%%%%%%
\begin{Lemma}\label{la_bomba}
Fix $c>0$ and, for each $p\in (0,1)$, fix a number $z(p)>0$ and  a probability measure $\mu^{(p)}$ on $\bbR^d$. Suppose that 
(i)  $\lim_{p\da 0} p\, z(p)^d=c$, 
(ii) $\mu^{(p)}\Rightarrow  \mu$ as $p \downarrow 0$.
Then, for any  stationary  ergodic   SPP  on $\bbR^d$ with law $\hat \bbP$ and finite  intensity $\rho$,  it holds 
$
\bigl( \hat \bbP_p \circ \Psi_{z(p)}  ^{-1}\bigr)_{\mu^{(p)}} \Longrightarrow {\rm PPP}[c \rho, \mu]$ as $p\downarrow 0$.
\end{Lemma}
Let us describe  by worlds the SPP with law $( \hat \bbP_p \circ \Psi_{z(p)}  ^{-1})_{\mu^{(p)}}$. It can be obtained from  a SPP $\hat \xi$ on $\bbR^d$  with law $\hat\bbP$ by  making the following transformations: 
%\begin{itemize}
%\item[1)]  
(1) We thin $\hat \xi$ by  keeping each point with probability $p$ independently from all the rest. The resulting SPP has law $\hat \bbP_p $.
%\item[2)]
(2)  We rescale the SPP  produced in the previous step by mapping each point $x$ into $x/z(p)$. The resulting SPP has law $\hat \bbP_p \circ \Psi_{z(p)}  ^{-1}$.
%\item[3)]
(3)  We mark each point of the SPP  produced in the previous step by i.i.d. random variables  (independently from all the rest) with common distribution $\mu^{(p)}$.
%\end{itemize}
%We point out that Steps 2) and 3) can be inverted: one could first mark the  SPP produced in 1) by  i.i.d. random variables  (independently from all the rest) with common distribution $\mu^{(p)}$, and afterwards rescale the resulting SPP by mapping $x$ into $x/z(p)$.
%%%%%%%%%%%%%%%%%%%%%%%%%%%%%%%%%%
\begin{proof}[Proof of Lemma \ref{la_bomba}]
%Since $\hat\bbP$ is stationary with finite intensity, the same holds for $ \hat \bbP_p \circ \Psi_{z(p)}  ^{-1}$. As a consequence, for any boxed $A\subset \bbR^d$,  $\hat \o(\partial A)=0$    for $ \hat \bbP_p \circ \Psi_{z(p)}  ^{-1}$--a.a.~$\hat \o\in \hat \O$. This property will allow us to apply Lemma \ref{weak_conv}.  
Since {\rm PPP}$[c \rho,\mu]$ is the $\mu$--randomization of {\rm PPP}$[c\rho]$\rosso{,} since 
$\mu^{(p)}\Rightarrow  \mu$ \rosso{and due to Remark \ref{condimento}, we can apply   Lemma \ref{weak_conv}. As a consequence,} it is enough to prove that  $ \hat \bbP_p \circ \Psi_{z(p)}  ^{-1}\Rightarrow {\rm PPP}[c \rho]$.
 Let us call $\hat \xi$ the SPP  with law $\hat\bbP$ (shortly, $\hat \xi \stackrel{\cL}{\sim} \hat \bbP$).
 Then, we define $\rosso{\hat \xi}^{(p)}$ as the SPP  $\rosso{\hat \xi}^{(p)}:= \Psi_{z(p) } (\hat \xi)$.  Note that $\rosso{\hat \xi}^{(p)} \stackrel{\cL}{\sim}  \hat \bbP\circ \Psi_{z(p)}  ^{-1}$. We finally define $\rosso{\hat \z}^{(p)}$ as the $p$--thinning of $\rosso{\hat \xi}^{(p)}$. \rosso{Then}
$\rosso{\hat  \z}^{(p)} \stackrel{\cL}{\sim}  (\hat \bbP\circ \Psi_{z(p)}  ^{-1})_p=\hat \bbP_p\circ \Psi_{z(p)}  ^{-1}$. \rosso{Hence, we need to prove that the SPP ${\hat\z}^{(p)}$  converges \rosso{weakly} to  {\rm PPP}$[c \rho]$.}
 %We identify the SPP $\rosso{\hat \xi}^{(p)}$ with a random measure on $\bbR^d$.
 As  {\rm PPP}$[c \rho]$ corresponds to a Cox process directed by $c  \rho dx$ (cf.~\cite[Chp.~12]{Kal}), by  \cite[Thm.~16.19]{Kal} it is enough to prove that the random measure  $p\rosso{\hat  \xi^{(p)} }$ (defined as $[p\rosso{\hat \xi}^{(p)} ](B):=p \cdot \rosso{\hat \xi}^{(p)}(B)$)  converges \rosso{weakly}  to the deterministic measure $c \rho dx$.  Due to \cite[Thm.~16.16]{Kal}  it is enough to prove that, fixed a continuous function $f:\bbR^d \to \bbR$ with compact support, the random variable \rosso{$p \sum _{x \in \rosso{\hat \xi}^{(p)} } f(x)$}  
% \[p \eta^{(p)} [f]= p \int f( x ) d \eta^{(p)} (x)  = p \int f( x/z(p) ) d \xi (x) \]
%\be\label{zeffiro90}p \sum _{x \in \rosso{\hat \xi}^{(p)} } f(x) = p \sum_{x\in \hat \xi}  f( x/z(p) ) \en
  converges \rosso{weakly} to $c \rho \int f(x) dx$ as $p\downarrow 0$. Setting $\e=\e(p):= 1/z(p)$, we have 
 %\[ p \int f( x/z(p) ) d \xi (x)= ( p /\e(p)^{-d})  \e(p)^d   \int f(\e(p) x) d \xi(x)\,.  \] 
 \[p \sum _{x \in \rosso{\hat \xi}^{(p)} } f(x) =
 p \sum_{x\in \hat \xi}  f( x/z(p) )=( p /\e(p)^{-d})  \e(p)^d   \sum_{x\in \hat \xi}  f(\e(p) x) 
\]
 Note that $\e (p) \da 0$ as $p\da 0$ since  $p z(p) ^d\to c>0$ as $p\da 0$. 
 As $\hat \xi$ is an ergodic stationary SPP,   $\e^d  \sum_{x\in \hat \xi} f(\e x)$ converges a.s. (and therefore weakly) to $\rho \int f(x) dx$  as $\e\da 0$ (see e.g.~\cite[Prop.~3.1]{Fhom1},  \cite{T}). Since  $p /\e(p)^{-d}=p z(p) ^d\to c$ as $p\downarrow 0$,  $p \sum _{x \in \rosso{\hat \xi}^{(p)} } f(x) $  converges  \rosso{weakly} to $c \rho \int f(x) dx$ as $p\downarrow 0$.
 \end{proof}
 
 

 Recall \eqref{alloro},  \eqref{sette1}, \eqref{sette2} and \eqref{flautino2}.
 The following technical fact will be frequently used in the rest:
 \begin{Lemma}
\label{fadeev_250}  Let  $\z,\b,\rho >0$, let $\nu$ be a  probability measure   on $\bbR$ and let $\hat\bbP$ be a  SPP.
   Sample 
 $\o$  according to   $\hat \bbP_\nu$ and    set  $p:=\nu( \z/\b)$. 
 Then
$\o_{\z/\b}$ has law $(\hat \bbP_p)_{\nu_{\z/\b}} $ 
and 
  \be\label{box_pokemon}
  \o_{\z,\b}:=\left\{ \big (x/\z, \b E_x/\z\big) \,: \,  x\in \widehat{\o}_ {\z/\b} \right \}
\en
  has law  $(\hat \bbP_p\circ \Psi_{\z}^{-1})_{\nu_{\star,  \z/\b}}$.  In particular, if    $\hat \bbP={\rm PPP}[\rho]$ and $\hat \bbP_\nu={\rm PPP}[\rho,\nu]$, then 
 \[ 
  (\hat \bbP_p)_{\nu_{\z/\b}} ={\rm PPP}[\rho \nu(\z/\b) ,\nu_{\z/\b}]\qquad 
   (\hat \bbP_p\circ \Psi_{\z}^{-1})_{\nu_{\star,  \z/\b}}= {\rm PPP}[\rho \nu(\z/\b)\z^d   ,\nu_{\star,  \z/\b}] \,.
  \]
\end{Lemma}
Note that $\o_{\z,\b}$ equals \eqref{gandhi} when $\z=\ell(\b)$.
\begin{proof}
  By sampling $\o$ according to $\hat\bbP_\nu$, $\o_{\z/\b}$ has law $(\hat \bbP_p)_{\nu_{\z/\b}} $  by    Lemma \ref{lemma_2021}.
The fact that $\o_{\z,\b}$ has law $(\hat \bbP_p\circ \Psi_{\z}^{-1})_{\nu_{\star,  \z/\b}}$ follows  in particular from the observation that  if   the random variable  $E$ has law  $\nu$, then conditioning to the event  $|E|\leq \z/\b$ the random variable 
$\b E /\z $ has law $\nu_{\star,  \z/\b}$ (see  \eqref{sette1_bis} and  \eqref{sette2}).
The conclusion concerning marked PPP's follows from 
 Proposition \ref{cotechino}. 
\end{proof}


\subsection{Proof of Theorem  \ref{paradiso}}\label{sec_dim_paradiso} We set
\be \label{meta}
p:=\nu( \ell(\b)/\b)\,, \qquad \g := \ell(\b)/\b\,, \qquad z = \ell (\b)\,.
\en
%By Lemma \ref{lemma_2021}, when sampling $\o$ with law $\hat \bbP_\nu$,  the SPP 
%\[ \o_\g=\o_{\ell(\b)/\b}=
%\bigl\{ (x, E_x)   \,: \, x\in\widehat \o  \,, |E_x| \leq  \ell(\b) /\b \bigr \}
%\]
%has law $( \hat \bbP _p) _{\nu _\g} $.
%We then deduce that
By Lemma \ref{fadeev_250}
  \eqref{gandhi} has law $(\hat \bbP_p \circ \Psi_z^{-1} ) _{\nu_{\star,\g}}$.





We now want to apply Lemma \ref{la_bomba}. 
Firstly we   note that, by \eqref{vichinghi100} and \eqref{canditi105},   for $\b \to +\infty$
\[ 
p z^d= \nu\bigl( \ell(\b)/\b\bigr) \ell(\b)^d\asymp C_*  \bigl(\ell(\b)/\b) ^{\a+1} \ell(\b)^d =C_* \ell(\b)^{\a+1+d} \b ^{-\a-1}=\l/ \rho \,.
\]
As a consequence $\lim_{\b \to +\infty }p z^d = \l/\rho$.
% (compare with Item (i) in Lemma \ref{la_bomba}).  
  Secondly  we observe that 
  \[
  \g:=\ell(\b)/\b= \bigl(\l /(C_* \rho)\bigr)^{\frac{1}{\a+1+d}} \b ^{-\frac{d}{\a+1+d}}
  \] goes to zero when $\b \to+\infty$. \rosso{Then,} by \eqref{vichinghi100},  
  $p:=\nu( \ell(\b)/\b)=\nu(\g)$  goes to zero as  $\b \to +\infty$. \rosso{Moreover,} \rosso{by Item (i) of Theorem \ref{paradiso}}, $\nu_{\star,\g}$ converges to $\nu^+_{1,\a}$ or $\nu_{1,\a}$ as $\b \to +\infty$. \rosso{Due to the above observations}, the weak convergence of \eqref{gandhi} to \eqref{ghibaudo}  follows from Lemma \ref{la_bomba}. 


%%%%%%%%%
\subsection{Proof of  Proposition \ref{paradiso_PPP}}\label{Trento}
%We start with a simple preliminary fact:
%
%\begin{Lemma}\label{fissato}Given $\a\geq 0$ and given $\g>0$ the following holds (recall the notational Warning \ref{pericolo}):
%\begin{itemize}
%\item[(i)] If    $\nu\in \cV^+$ and $\g\leq \min \{C_0, \e\} $, then 
%$\nu_{\star,\g} =\nu^+_{1,\a}$ and $   \nu(\g)    =   \left(  \g/ C_0  \right) ^{\a+1} $.
%\item[(ii)] If    $\nu\in \cV$ and $\g\leq \min \{C_0, \e\} $, then 
%$\nu_{\star,\g} =\nu_{1,\a}$ and $   \nu(\g)    =   \left(  \g/ C_0  \right) ^{\a+1} $.
%\end{itemize}
%\end{Lemma}
%%%%%%%%%%%%%%%%%%%%%%%%%%%%%%%%%%%%%%%%%%%%%%%%%%%%%%%%%%%%%%%%%%%%%%%%%%%%%%%%%%%%%%%%%%%%%%%%%%%%%%%%%%%%%%%%%%%%%%%%%%%%%%%%%%%%%%%%%%%%%%%%%%%%%%%%%%%%%%%%%%%%%%%%%%%%%%%%%%%%%%%%%%%%%%%%%%%%%%%%%%%%%%%%%%%%%%%%%%%%%%%%%
%\begin{proof}
%We give the proof of Item (i), the proof of  Item (ii) is similar. 

We prove Items (i) and (ii) only for $\nu \in \cvp$ as the case $\nu \in \cv$ is similar.
 We first prove our claim for $\nu=\nu_{C_0,\a}^+$. In this case  $\e=C_0$, hence the upper bound  $\g \leq \min\{C_0, \e\}$ reads $\g\leq C_0$.
 Let $X$ be a random variable with law $\nu$. For $0 \leq t \leq C_0$ we have $\nu(t)=\rosso{\nu([-t,t])=}(t/C_0)^{\a+1}$. Then, 
given $0\leq u \leq 1$ and  using that  $\g \leq C_0$, we get 
 \begin{equation*}
 \nu_{\star,\g} ([0,u])= P\bigl(  X/\g\in [0,u]\,\big{|}\, X/\g \leq 1\bigr) =  \nu ( u\g)/\nu(\g)   =  u^{\a+1}=\nu^+_{1,\a} (  [0,u])\,.
  \end{equation*}
  This proves  that $  \nu_{\star,\g} =\nu^+_{1,\a}$.  Trivially, $\nu([-\g, \g])= \nu(\g)=(\g/C_0)^{\a+1}$.   Items (i) and (ii) are  then proved for $\nu=  \nu_{C_0,\a}^+$.
 We now move to a generic $\nu\in \cvp$. The value $\nu(\g)$ and the probability  $\nu_{\star,  \g}$ depend only on the measure $\nu$ restricted to $[-\g,\g]$  which equals the measure
 $\nu_{C_0,\a}^+$ restricted to $[-\g,\g]$ (as $\g\leq \e$). Hence  
    we have  $\nu(\g)=\nu'(\g)$  and  $\nu_{\star,  \g}=\nu'_{\star,  \g}$ where $\nu':=\nu_{C_0,\a}^+$.  The conclusion then follows from the special case 
 $\nu=\nu_{C_0,\a}^+$ treated above since $\g\leq C_0$.
 
 
  

\smallskip
We move to the proof of  the final statement.
If one starts with $\o$ sampled by  ${\rm PPP}[\rho, \nu]$, then by Lemma \ref{fadeev_250}
%
% To this aim recall first  Proposition \ref{cotechino}\rosso{-(iii)}, \eqref{sette1} and \eqref{sette1_bis}.
%If one starts with $\o$ sampled according to  ${\rm PPP}[\rho, \nu]$, then $\o_{\g}$ with $\g:=\ell(\b)/\b$ has law
%${\rm PPP}\left[\,\rho \nu(\g)  , \nu_\g \,\right]$. If we then rescale points in  $\o_{\g}$ by $\Psi_{\ell(\b)}$ (i.e.~$x\mapsto x/\ell(\b)$)  by keeping the same energy marks, the resulting marked SPP has law   ${\rm PPP}\left[\rho \nu(\g) \ell(\b)^d  , \nu_\g\,\right]$. If afterwards we replace the energy mark $E_x$ with $E_x/\g$ we get a marked SPP with law $
%{\rm PPP}\left[\,\rho \nu( \g ) \ell(\b)^d  , \nu_{*,\g} \,\right]$. In particular
%
%
 the law of \eqref{gandhi} is $
{\rm PPP}\left[\,\rho \nu( \g ) \ell(\b)^d  , \nu_{*,\g} \,\right]$.
As soon as  $\b$ is  large enough to assure that  $\g\leq \min\{C_0,\e\}$,  by Item (ii) we have 
\[
\rho \,\nu( \g ) \ell(\b)^d = \frac{\rho \g^{\a+1}}{C_0^{\a+1}} \ell(\b)^d= \frac{\rho}{ C_0^{\a+1}} \ell(\b)^{\a+1+d} \b^{-\a-1}= \frac{\rho}{ C_0^{\a+1}}   (\l/(C_*\rho) )=\l\,.\]
As a consequence, $
{\rm PPP}\left[\,\rho \nu( \g ) \ell(\b)^d  , \nu_{*,\g} \,\right]=
{\rm PPP}\left[\,\l , \nu_{*,\g} \,\right]$. Item (i) allows now to conclude.



%%%%%%%%%%%%%%%%%%%%%%%%%%%%%%%%%%


\section{Rescaling of  $\cG[\z,\b]$ and proof of Theorem \ref{teo1}}\label{escher}

The results presented in the first part of this section hold for all dimensions $d\geq 1$. In Subsection  \ref{sec_proof_teo1} we will restrict to $d\geq 2$.
 
 \smallskip
 
In this section we investigate the effect of rescaling on the random graph $\cG[\z,\b]$ introduced in  Definition \ref{giggio_bello}. We recall that, given $\o\in \O$, 
$\cG[\z,\b](\o)$  has vertex set $\widehat{\o}$  and edges given by pairs $\{x,y\}$ with $x\not =y$ satisfying 
\eqref{connectcont}, i.e.~$|x-y|+\beta(|E_x|+|E_y|+|E_x-E_y)|)\leq \z$. Afterwards we will apply these results to derive  Theorem \ref{teo1}.

\smallskip



For the next \orc{lemma} recall Definition \ref{stella_stellina}  of  $\nu_{\star,\g}$, Definition   \ref{san_benedetto} of  $\nu(\g)$ and  $\nu_\g$ and  
recall Definition \ref{porto_dascoli} of  $\widehat{\o}_\g$ and $\o_\g$.

%%%%%%%%%%%%%%%%%%%%%%%%%%
%%%%%%%%%%%%%%%%%%%%%%%%%%
\begin{Lemma}
\label{fadeev_new}  Let  $\z,\b,\rho >0$  and let $\nu$ be a  probability measure   on $\bbR$.  Then the  following holds:
\begin{itemize}
\item[(i)] For any $\o\in \O$
the graph $\cG[\z, \b](\o)$  is  the  union of the graph
$\cG[\z,\b](\o_{\z/\b})$   and the  singletons $\{x\}$ with  $x\in \widehat \o  \setminus \widehat \o_{\z/\b}$, which are isolated points in $\cG[\z, \b](\o)$. 
\item[(ii)]   
Calling $\Psi_\z \cG[\z,\b] (\o_{\z/\b})   $  the image of $ \cG[\z,\b] (\o_{\z/\b}) $ under the map  $\Psi_\z:x\mapsto x/\z$, the graph $\Psi_\z \cG[\z,\b] (\o_{\z/\b}) $  equals the graph $\cG[1,1](\o_{\z,\b})$ where  
  $\o_{\z,\b}:=\left\{ \big (x/\z, \b E_x/\z\big) \,: \,  x\in \widehat{\o}_ {\z/\b} \right \}
  $
  as in \eqref{box_pokemon}.
 \end{itemize}
\end{Lemma}
% Note that $\o_{\z, \b}$ equals   \eqref{gandhi}   when $\z=\ell(\b)$.
\begin{proof}
We start with Item (i).
Due to \eqref{connectcont}, if $ |E_x|>\z/\b$, then $x$ is an isolated point in the graph $\cG[\z,\b](\o) $. The conclusion is then immediate. Item
(ii) follows from the observation that \eqref{connectcont} is equivalent to 
$
\big|x/\z -y/\z\big |+\left(\big|\b E_x/\z \big|+\big|\b E_y/\z\big|+\big| \b E_x/\z-\b E_y/\z\big|\right)\leq 1.$
\end{proof}

%%%%%%%%%%%%%%%%%%%%%%%%%%%%%%%%%%%%%%%%%%%%%%%%%%%%%%%%%%%%%%%%%%%%%%%%%%%%%%%%%%%%%%%%%%%%%%
%%%%%%%%%%%%%%%%%%%%%%%%%%
%\begin{Proposition}
%\label{fadeev_new}  Let  $\z,\b,\rho >0$  and let $\nu$ be a  probability measure   on $\bbR$.
% Then the  following holds:
%\begin{itemize}
%\item[(i)] For any $\o\in \O$
%the graph $\cG[\z, \b](\o)$  is  the  union of the graph
%$\cG[\z,\b](\o_{\z/\b})$   and the  singletons $\{x\}$ with  $x\in \widehat \o  \setminus \widehat \o_{\z/\b}$, which are isolated points in $\cG[\z, \b](\o)$. 
%\item[(ii)] The graph $\cG[\z,\b](\o_{\z/\b})$  with 
% $\o$  sampled according to   $\hat \bbP_\nu$ has the same distribution of the graph  
%$\cG[\z,\b](\o)$ with $\o$ sampled according to   $(\hat \bbP_p)_{\nu_{\z/\b}} $ where  $p:=\nu( \z/\b)$. In particular, if    $\hat \bbP_\nu={\rm PPP}[\rho,\nu]$, then $(\hat \bbP_p)_{\nu_{\z/\b}} ={\rm PPP}[\rho \nu(\z/\b) ,\nu_{\z/\b}]$.
%\item[(iii)]   
%Calling $\Psi_\z \cG[\z,\b] (\o_{\z/\b})   $  the image of $ \cG[\z,\b] (\o_{\z/\b}) $ under the map  $\Psi_\z:x\mapsto x/\z$, the graph $\Psi_\z \cG[\z,\b] (\o_{\z/\b}) $  equals the graph $\cG[1,1](\o_{\z,\b})$ where  
%\be\label{box_pokemon}
%  \o_{\z,\b}:=\left\{ \big (x/\z, \b E_x/\z\big) \,: \,  x\in \widehat{\o}_ {\z/\b} \right \}
%\,.
%\en
%\item[(iv)] The graph $\Psi_\z \cG[\z,\b] (\o_{\z/\b}) $ \rosso{equals the graph $\cG[1,1](\o_{\z,\b})$}
% with 
% $\o$  sampled according to   $\hat \bbP_\nu$ has the same distribution of the graph $\cG[1,1](\o)$ with $\o$ sampled according to $(\hat \bbP_p\circ \Psi_{\z}^{-1})_{\nu_{\star,  \z/\b}}$ where $p:=\nu(\z/\b)$. In particular, if    $\hat \bbP_\nu={\rm PPP}[\rho,\nu]$, then 
% $(\hat \bbP_p\circ \Psi_{\z}^{-1})_{\nu_{\star,  \z/\b}}= {\rm PPP}[\rho \nu(\z/\b)\z^d   ,\nu_{\star,  \z/\b}] $.
% \end{itemize}
%\end{Proposition}
% Note that $\o_{\z, \b}$ equals   \eqref{gandhi} 
%when $\z=\ell(\b)$.
%%%%%%%%%%%%%%%%%%%%%%%%%%%%%%%%%%%%%%%%%%%%%%%%%%%%%%%%%%%%%%%%%%%%%%%%%%%%%%%%%%%%%%%%%%%%%%%%%%%%%%%%%%%%%%%%%%%%%%%%%%%%%%%%%%%%%%%%%%%%%%%%%%%%%%%%%%%%%%%%%%%%%%%%%%%%%%%%%%%%%%%%%%%%%%%%%%%%%%%%%%%%%%%%%%%%%%%%%%%%%%
%\begin{proof}
%We start with Item (i).
%Due to \eqref{connectcont}, if $ |E_x|>\z/\b$, then $x$ is an isolated point in the graph $\cG[\z,\b](\o) $. The conclusion is then immediate.
%Item (ii) follows from Lemma \ref{lemma_2021} and  Proposition \ref{cotechino}. 
%Let us prove Item (iii). 
%%Due to Item (iii) we just need to show that
%%\be \Psi_\z\,  \cG[\z,\b]_{ \rho \nu(\z/\b) ,\nu_{\z/\b}}(\o) \stackrel{\cL}{=}   \cG[1,1]_{\rho \nu(\z/\b)\z^d   ,\nu_{\star,  \z/\b}}(\o)\,.   \en
%Let $\o=\{(x, E_x)\,:\, x\in \widehat \o \}$.
% We have $\{x,y\}\in  \cG[\z,\b](\o)$ if and only if \eqref{connectcont} holds, i.e. 
% \begin{equation}
%\frac{|x-y|}{\z}+\frac{\beta}{\z}\left(|E_x|+|E_y|+|E_x-E_y|\right)\leq 1.
%\end{equation}
%Hence, if $\o=\{(x, E_x)\,:\, x\in \widehat \o \}$, 
%$\Psi_\z\,  \cG[\z,\b](\o)$ has vertex set $\{x'\,:\, x\in \widehat \o \}$ where $x':= x/\z$ and edges  given by the unordered pairs $\{x',y'\}$ with $x'\not =y'$  such that 
%\begin{equation}
%|x'-y'|+|A_{x'}|+|A_{y'}|+|A_{x'}-A_{y'}| \leq 1\,, 
%\end{equation}
%where we have set $A_{x'}:=(\b/\z)E_x$. This leads to Item (iii).
%
%
%Let us prove  Item (iv). We have already observed that  $\o_{\z/\b}$ has law  $(\hat \bbP_p)_{\nu_{\z/\b}} $ if $\o$ is sampled according to $\hat\bbP_\nu$. Then the first part of Item (iv) follows from Item (iii) and a rescaling. The second part follows from Proposition \ref{cotechino}.
%% We observe that, if  $\hat \o $ is sampled according to   ${\rm PPP}[\rho \nu(\z/\b)]$, then  $\{x'\,:\, x\in \widehat \o \}$   has law   ${\rm PPP}[\rho \nu(\z/\b)\z^d ]$. Since $A_{x'}$  is  given by $(\b/\z)E_x$ and $E_x$ has distribution $\nu_{\z/\b}$, we get  that $A_{x'}$ has distribution  $\nu_{\star,  \z/\b}$.  Finally the independence of the $A_{x'}$s follows from the independence of the $E_x$s. 
%\end{proof}


%%%%%%%%%%%%%%%%%%%%%%%%%%%%%%%%%%%%%%%%
%\begin{Proposition}
%\label{fadeev_new}  Let  $\z,\b >0$  and let $\nu$ be a  probability  on $\bbR$. Then the  following holds:
%\begin{itemize}
%\item[(i)]
%The graph $\cG[\z, \b](\o)$   can be written as the disjoint union of the graph
%$\cG_*[\z,\b](\o)= \bigl( \hat \o_*, \cE_*[\z,\b]\bigr)$   and the  singletons $\{x\}$ with  $x\in \hat \o  \setminus \hat \o_*$ (which are isolated points in $\cG[\z, \b](\o)$). 
%\item[(ii)] The graph $\cG_*[\z,\b](\o)$  with 
% $\o$  sampled according to  the $\nu$--randomization of $\hat \bbP$ has the same distribution of the graph  
%$\cG[\z,\b](\o')$ with $\o'$ sampled according to   the $\nu_{\z/\b}$--randomization of $\hat \bbP_p$ with thinning parameter $p:=\nu([-\z/\b, \z/\b])$.
%\item[(iii)] It holds 
%\be
%\cG_*[\z,\b]_{\rho,\nu}\stackrel{\cL}{=} \cG[\z,\b]_{\l(\z/\b,\rho,\nu)  ,\nu_{\z/\b}}\,.
%\en
%\item[(iv)] Calling  $\Psi: \bbR^d\to \bbR^d$ the bijection $\Psi(x):= x/\z$, the pushed forward graph  
%$\Psi_* \cG_*[\z,\b]$ (i.e. the image of $\cG_*[\z,\b]$ under $\Psi$) satisfies 
%  \be\label{fuocone}
%  \Psi_*  \cG_*[\z,\b]_{\rho,\nu} \stackrel{\cL}{=} 
%   \cG_*[1,1]_{\l(\z/\b, \rho,\nu)\z^d   ,\nu_{\star,  \z/\b}}\,.
%   \en
%\end{itemize}
%\end{Proposition}


%%%%%%%
%
%\begin{Proposition}[Effect of rescaling when $\bbP=PPP(\rho,\nu)$]
%\label{fadeev}  Let  $\z,\b,\r >0$  and let $\nu$ be a  probability  on $\bbR$. Then the  following holds:
%\begin{itemize}
%\item[(i)]
%The graph $\cG=(\xi,\cE):=\cG[\z, \b,\rho ,\nu] $   can be written as the disjoint union of the graph  $\tilde \cG=\bigl(\tilde \xi, \tilde \cE\bigr)$, where 
% \begin{align*}
% & \tilde \xi:= \{ x\in \xi\,:\, |E_x| \leq \z/\beta\}\,,\\
% &  \tilde{\cE}:=\big\{ \{x,y\} \in \cE\,:\, x,y \in \tilde \xi\}\,,
%  \end{align*} 
% and the  singletons $\{x\}$ with  $x\in \xi \setminus \tilde{\xi}$ (which are isolated points in $\cG$). 
%  Moreover, the graphs $\tilde \cG$ and $\cG \left[\z, \b,\l(\z/\b,\rho,\nu)  ,\nu_{\z/\b}\right ] $ have the same  laws:
%  \be\label{fuoco} \tilde \cG\stackrel{\cL}{=} \cG \left[\z, \b,\l(\z/\b,\rho,\nu)  ,\nu_{\z/\b}\right ]\,.\en
%\item[(ii)] Calling  $\Psi: \bbR^d\to \bbR^d$ the bijection $\Psi(x):= x/\z$, the pushed forward graph  
%%$\Psi_*\cG$ is the disjoint union of the set of isolated points $\Psi (\xi\setminus \tilde\xi)$ and the graph 
%$\Psi_* \tilde \cG$ (i.e. the image of $\tilde \cG$ under $\Psi$) satisfies 
%  \be\label{fuocone}
%  \Psi_* \tilde \cG \stackrel{\cL}{=} 
%   \cG\left [1, 1,\l(\z/\b, \rho,\nu)\z^d   ,\nu_{\star,  \z/\b}\right]\,.
%   \en
%\end{itemize}
%\end{Proposition}
%%%%%%%%%%%%%%%
%\begin{proof} We start with Item (i).
%Due to \eqref{connectcont}, if $ |E_x|>\z/\b$, then $x$ is an isolated point in the graph $\cG $. Hence $\cG$ is  the disjoint union of  the points  $x\in \xi\setminus \tilde \xi$, which are isolated,  and the graph $\tilde \cG$.
%
%
% We prove \eqref{fuoco}.
%The  vertex set  $\tilde \xi$ is obtained by  thinning $\xi$. In particular, a point $x$ in $\xi$ belongs to    $\tilde \xi$ independently from the other points  with probability 
%$ \nu \bigl([ -\z/\b , \z/\b]\bigr)$.
%Hence, $\tilde \xi$ is a homogeneous PPP with density  $\l(\z/\b, \rho,\nu)$.
%Moreover, if we know that  $x\in \tilde \xi$, the energy mark $E_x$ has distribution $\nu_{\z/\b}$.  The independence of $(E_x)_{x\in \tilde \xi}$ conditioned to $\tilde \xi$ is conserved.
%%Hence, $\tilde \cG= \cG[\z, \b,\r(\z,\b)  ,\nu_{\z,\b} ]$.
%
%We move to Item (ii).  We rewrite  \eqref{connectcont} as 
%\begin{equation}
%\frac{|x-y|}{\z}+\frac{\beta}{\z}\left(|E_x|+|E_y|+|E_x-E_y|\right)\leq 1.
%\end{equation}
%Hence $\cG':= \Psi_*\tilde \cG$ has vertex set $\xi'= \{x/\z\,:\, x\in \tilde \xi\}$ and edges  given by the unordered pairs $\{x',y'\}$ with $x'\not =y'$ in $ \xi'$ such that 
%\begin{equation}
%|x'-y'|+|A_{x'}|+|A_{y'}|+|A_{x'}-A_{y'}| \leq 1\,, 
%\end{equation}
%where we have set $A_{x'}:=(\b/\z)E_x$ if $x'=x/\z$. Since $\xi'$ is a PPP with density $\l(\z/\b, \rho,\nu) \z^d$ and since $A_{x'}$  is 
%given by $(\b/\z)E_x$ and $E_x$ has distribution $\nu_{\z/\b}$, we get  \eqref{fuocone}.
%%that $\cG'= \cG[1, 1,\r(\z,\b)\z^d   ,\nu_{\star,  \z,\b}]$.
%\end{proof}
%%%%%%%%%%%%%%
An immediate consequence of   \orc{Lemmas  \ref{fadeev_250} and \ref{fadeev_new}} is the following fact:
\begin{Corollary}\label{funghetto} We have
%\begin{multline}\label{olla}
\[
{\rm PPP}[\rho,
\nu]\big(\, \cG[\z, \b] \text{ percolates}\,\big)
%= {\rm PPP} [\rho \nu(\z/\b) ,\nu_{\z/\b}]  \Big(\, \cG[\z,\b]\text{ percolates}\, \Big)\\
= {\rm PPP}[\rho \nu(\z/\b)\z^d   ,\nu_{\star,  \z/\b}]
 \big( \, \cG[1,1]\text{ percolates}\, \big)\,.
 \]
%\end{multline}
\end{Corollary}
For $\nu$ given by  \eqref{speciale2} or \eqref{speciale1} or in general for $\nu\in \cvp \cup  \cv$,  the above r.h.s. can be made more explicit due to the following fact, which is an immediate consequence of Items (i) and (ii) in Proposition \ref{paradiso_PPP} applied with \rosso{$\g= \z/\b$} (hence we omit the proof). \rosso{Recall  the notational Warning \ref{pericolo}.}
%%%%%%%%%%%%%%%%%%%%%%%%%%%%%%%%%%%%%%%%%%
\begin{Lemma}\label{fissato} \rosso{Let  $\a\geq 0$, $\nu \in \cvp\cup\cv$ and  let $\z,\b>0$ be  such that $\z\leq \min \{C_0, \e\}\b $. Then $\nu_{\star,\z/\b} $ equals $\nu^+_{1,\a}$ when $\nu\in \cvp$ and equals $\nu_{1,\a}$ when $\nu\in \cv$. Moreover, $  \rho \nu(\z/\b)\z^d    = \rho  \left(  \b C_0  \right) ^{-(\a+1)} \z^{\a+1+d}$.}
\end{Lemma}

%%%%%%%%%%%%%%%%%%%%%%%%%%%%%%%%%%%%%%%%%%%%%%%%%%%%%%%%%%%%%%%%%%%%%%%%%%%%%%%%%%%%%%%%%%%%%%%%%%%%%%%%%%%%%%%%%%%%%%%%%%%%%%%%%%%%%%%%%%%%%%%%%%%%%%%%%%%%%%%%%%%%%%%%%%%%%%%%%%%%%%%%%%%%%%%%%%%%%%%%%%%%%%%%%%%%%%%%%%%%%%%%%
%\begin{proof}
%We give the proof of Item (i), the proof of  Item (ii) is similar. 
%
% We first prove our claim for $\nu=\nu_{C_0,\a}^+$. In this case we can take $\e=C_0$, hence the upper bound on  $\z$ reads $\z\leq C_0\b$.
% Let $X$ be a random variable with law $\nu$.
%Given $0\leq u \leq 1$ and  using that  $\z/\b \leq C_0$, we get 
% \begin{equation*}
% \begin{split}
% \nu_{\star,\z/\b} ([0,u])&= P\bigl(  (\b/\z)X\in [0,u]\,\big{|}\, (\b/\z)X \leq 1\bigr) =\frac{ P\bigl(  (\b/\z)X\in [0,u]\bigr)}{ P\bigl(  (\b/\z)X\leq 1\bigr) }\\
% &= \frac{ \nu ([0,(\z/\b) u])}{ \nu( [0, \z/\b] )}
%  =  \frac{\left(  \z u / \b C_0  \right) ^{\a+1}    }{ \left(  \z /\b C_0  \right) ^{\a+1}  } =  u^{\a+1}=\nu^+_{1,\a} (  [0,u])\,.
%  \end{split}
%  \end{equation*}
%  This proves  that $  \nu_{\star,\z/\b} =\nu^+_{1,\a}$.  Using that  $
% \nu(\z/\b)=  \left(  \z  / \b C_0  \right) ^{\a+1}
%  $,
%  we get the other identity in Item (i).
% 
% We now move to a generic $\nu\in \cV^+$. Since the value $\nu(\z/\b)$ and the probability  $\nu_{\star,  \z/\b}$ depend only on the measure $\nu$ restricted to $[-\z/\b,\z/\b]$ and since $\z/\b \leq \e$,   we have  $\nu(\z/\b)=\nu'(\z/\b)$  and  $\nu_{\star,  \z/\b}=\nu'_{\star,  \z/\b}$ where $\nu':=\nu_{C_0,\a}^+$.  The conclusion then follows from the special case 
% $\nu=\nu_{C_0,\a}^+$ treated above.
%   \end{proof}
%%%%%%%%%%%%%%%%%%%%%%%%%%%%%
%%%%%%%%%%%%%%%%%%%%%%%%%%%%%
%%%%%%%%%%%%%%%%%%%%%%%%%%%%%
%%%%%%%%%%%%%%%%%%%%%%%%%%%%%
%%%%%%%%%%%%%%%%%%%%%%%%%%%%%

\subsection{Proof of Theorem \ref{teo1}}\label{sec_proof_teo1} We take $d\geq 2$. 
We prove the theorem  for $\nu\in \cvp$. The   case $\nu\in \cv$  is similar.
To simplify the notation we call $\eta$ the r.h.s. of \eqref{zetacritico}.
We need to show  for  \orc{$\eta<  \min\{C_0 , \e\}  \b $}  that 
\begin{itemize}
\item[(1)]  if $\z<\eta$ then ${\rm PPP}[\rho,\nu] \left(\, \cG[\z, \b] \text{ percolates}\, \right)=0$, 
\item[(2)] if $\z>\eta$ then ${\rm PPP}[\rho,\nu]\left(\, \cG[\z, \b]\text{ percolates}\, \right)=1$.  
\end{itemize}
In Item (1)  we have \orc{$\z <  \min\{C_0 , \e\} \beta$} as $\eta <   \min\{C_0 , \e\} \beta$ by assumption. 
In Item (2) we can   restrict to $\z\in \left(\eta, \min\{C_0 , \e\} \beta\right]$ as this  set  is nonempty and the function  $\z\to  {\rm PPP}[\rho,\nu]\bigl( \cG[\z, \b] \text{ percolates}\, \bigr)$ does not decrease. Hence  in both Items (1) and (2)  we can restrict  $\z \leq \min\{C_0 , \e\} \beta $. Then, as a byproduct of   Corollary \ref{funghetto} and Lemma \ref{fissato}, we have 
  \begin{multline*}
  {\rm PPP} [\rho,\nu] \bigl( \cG[\z, \b] \text{ percolates}\, \bigr)
  = 
 {\rm  PPP}[\rho \nu(\z/\b) \z^d ,\nu_{\star,  \z/\b}] \bigl(  \cG[1, 1]  \text{ percolates}\,\bigr )\\
  ={\rm PPP}[ \rho  \left(  \b C_0  \right) ^{-(\a+1)} \z^{\a+1+d}, \nu^+_{1,\a}]\bigl(  \cG[1, 1]  \text{ percolates}\,\bigr )\,.
 \end{multline*}
 We now observe that $\z\lessgtr \eta$ if and only if 
% Since the value $\nu(\z/\b)$ and the probability  $\nu_{\star,  \z/\b}$ depend only on the measure $\nu$ restricted to $[-\z/\b,\z/\b]$ and since $\z/\b \leq \e$, in Case (A) we have  $\nu(\z/\b)=\nu'(\z/\b)$  and  $\nu_{\star,  \z/\b}=\nu'_{\star,  \z/\b}$ where $\nu':=\nu_{C_0,\a}^+$.  
% As  $\z/\b \leq C_0$, we can apply Lemma \ref{fissato}.  In particular, if $\z\lessgtr \eta$ and $\z \leq \min\{C_0 , \e\} \b$,  we get $\nu_{\star,  \z/\b}=\nu^+_{1,\a}$ and 
\[  \rho  \left(  \b C_0  \right) ^{-(\a+1)} \z^{\a+1+d} 
\lessgtr  \rho  \left(  \b C_0  \right) ^{-(\a+1)} \eta^{\a+1+d}=\l^+_c(\a)
\]
(the last identity follows immediately from \eqref{zetacritico}). 
  As $\l^+_c(\a)$ is the critical density for the graph  $\cG[1, 1](\o)$ with $\o$ sampled according to  ${\rm PPP}[\l,\nu^+_{1,\a}] $, we can conclude.
  %%%%%%%%%%%%%%%%%%%%%
  %
  %
  %
  %%%%%%%%%%%%%%%%%%%%%%%%%
  \section{Law of $\o_\g$ with $\o$ sampled according to $\bbP$}\label{barretta}
%%%%%%%%%%%%%%%%%%%%%%%%%%%%%%%%%%%%%%%%%%%%
In this  section we collect some properties of the law of $\o_\g$  with $\o$ sampled by $\bbP$, in addition to the ones presented in \orc{Section \ref{Tarvisio}}  \rosso{(recall from \eqref{flautino2} that $\o_\g=\{(x,E_x)\,:\, x\in \widehat{\o}\,,\; |E_x| \leq \g\}$)}. \rosso{The following result will be applied to $\bbP=\orc{\hat\bbP_\nu}$ in the next section, but the arguments used in its proof are completely general, hence we state it in full generality}.
%%%%%%%%%%%%%%%%%%%%%%%%%%%%%%%%%%%%%%%%%%%%
\begin{Lemma}\label{quinoa} Let $\bbP$ 
be the law of a stationary marked SPP with   intensity $\rho\in (0,+\infty)$ and let 
 $\vpz$ be the Palm distribution associated to    $\bbP$.  Fix $\g>0$ and 
    denote by $\bbQ$ the law of $\o_\g$ when $\o$ is sampled according to $\bbP$. 
Assume that   $\bbP( \o_\g \not =\emptyset)>0$.
Then $\bbQ $ is a stationary marked SPP  and \orc{its intensity, denoted by $\rho_\g$, is finite and positive}. Moreover, 
$\vpz(0\in \rosso{\widehat\o}_\g)=\rho_\g/\rho $ and 
 the Palm distribution $\vqz$ associated  to  $\bbQ $
equals 
  the law of $\o_\g$,  with $\o$ sampled according to $\vpz(\cdot | 0\in \rosso{\widehat\o}_\g)$.
\end{Lemma}
%%%%%%%%%%%%%%%%%%%%%%%%%%%%%%%%%%%%%%%%%%
\begin{proof}  \rosso{We recall  our notation \eqref{flautino1}:  $\widehat\o_\g:=\{x\in \widehat \o\,:\, |E_x| \leq \g\}$
(note that $\widehat{\o_\g}=\widehat{\o}_\g$).}
The stationary of $\bbQ $ is trivial. The intensity $\rho_\g$  of $\bbQ $  is finite since  upper bounded by the intensity  $\rho$ of $\bbP$. If  $\rho_\g $ was zero, then by stationarity we would have $\int_\O d\bbP(\o)  \widehat{\o}_\g (A)=0 $ for any $A\in \cB(\bbR^d)$. This would imply   that $\bbP( \o_\g =\emptyset)=1$, thus contradicting our assumption.

It remains to  prove the last statement in the lemma concerning Palm distributions.
To this aim let $h : \O_0 \to \bbR$ be a nonnegative measurable function. 
We  use Campbell's identity \eqref{campanello}  with  $\bbP$ replaced by $\bbQ$ and with test function $f(x,\o):= \mathds{1}_{[0,1]^d}(x) h(\o)$. We then get 
\be\label{sonnone1}
\begin{split}
\int _{\O_0} d \vqz (\o) h(\o) & = \frac{1}{\rho_\g} \int _\O d \bbQ (\o) \sum_{x\in \widehat{\o}\cap [0,1]^d} h(\t_x \o)\\
& = \frac{1}{\rho_\g} \int _\O d \bbP(\o) \sum_{x\in \widehat{\o}_\g \cap [0,1]^d}  h\bigl(\t_x (\o_\g)\bigr) \\
&= \frac{1}{\rho_\g} \int _\O d \bbP(\o) \sum_{x\in \widehat{\o}\cap [0,1]^d}h\bigl(\t_x (\o_\g) \bigr)\mathds{1}(|E_x|\leq \g) \,.
\end{split}
\en


We now apply Campbell's identity   \rosso{\eqref{campanello}  with} $f(x,\o)=  \mathds{1}_{[0,1]^d}(x) h(\o_\g) \mathds{1} ( 0\in \widehat{\o}_\g)$. Using that $(\t_x \o)_\g=\t_x (\o_\g)$, we get 
\begin{equation}\label{sonnone2}
%\begin{split}
   \int _{\O_0} d\vpz ( \o) h(\o_\g) \mathds{1}(0\in \widehat{\o}_\g )  =\frac{1}{\rho} \int _{\O}d \bbP(\o)  \sum_{x\in \widehat{\o}\cap [0,1]^d}   h\bigl(\t_x(\o_\g)\bigr) \mathds{1}(|E_x| \leq \g) \,.
   %\\ & =  \frac{1}{m} \int _{\O}d \bbP(\o)\int_{[0,1]^d}  d\hat{\o} (x) h\bigl(\t_x(\o_\g)\bigr) \mathds{1}(|E_x| \leq \g)  \,. \end{split}
 \end{equation}
 In the special case $h\equiv 1$  in \eqref{sonnone2} we get
 \begin{equation}\label{sonnone3}
   \vpz(0\in \widehat{\o}_\g ) 
   =  \frac{1}{\rho} \int _{\O}d \bbP(\o)  \sum_{x\in \widehat{\o}\cap [0,1]^d}   \mathds{1}(|E_x| \leq \g)=
   \frac{\bbE\bigl[ \widehat{\o}_\g([0,1]^d)\big]}{\rho}
 = \frac{\rho_\g}{\rho}\,. 
 \end{equation}
 By combining \eqref{sonnone1}, \eqref{sonnone2} and \eqref{sonnone3} we get 
 \be
 \int _{\O_0} d \vqz (\o) h(\o)=\frac{1}{ \vpz (0\in \widehat \o_\g ) }  \int _{\O_0} d\vpz ( \o) h(\o_\g) \mathds{1}(0\in \widehat \o_\g ) \,,
 \en
 thus proving our final claim in Lemma \ref{quinoa}.
\end{proof}

  
  
  
  %%%%%%%%%%%%%%%%%%%%%%%%%%%%%%%%%%%%%
%
%   Effective conductivity of the MA random resistor network
%
%
%%%%%%%%%%%%%%%%%%%%%%%%%%%%%%%%%%%%%


%%%%%%%%%%%%%%%%%%%%%%%%%%%%%%%%
%%%%%%%%%%%%%%%%%%%%%%%%%%%%%%%%
%%%%%%%%%%%%%%%%%%%%%%%%%%%%%%%%
%%%%%%%%%%%%%%%%%%%%%%%%%%%%%%%%%%%%%%
\section{Proof of Theorem \ref{parataUB} (upper bound in Mott's law)}\label{sec_UB}
Many estimates presented in the proof below work for  general stationary marked SPPs \orc{$\bbP$} with finite and positive intensity (we recall that in this case the Palm distribution is well defined, see Section \ref{palmato}). We therefore  keep a general approach as most as possible. This also will allow to make clearer  the universality behind Mott's law.
 Only at the end we will restrict to the setting of Theorem \ref{parataUB}. 


We need to upper bound 
 \begin{equation}\label{poh}
 D(\b)_{1,1} =\inf _{ f\in L^\infty(\bbP_0) } \frac{1}{2}\bbE_0 \Big[
 \sum_{x\in \widehat \o}  c_{0,x}(\o,\b) \left
 (x_1 - \nabla_x f (\o) 
\right)^2\Big]\,,
 \end{equation}
 where $\nabla_x f (\o) := f(\t_x \o) - f(\o)$.  \orc{Above, $\bbE_0$ denotes the expectation w.r.t. $\vpz$}.
 
 Due to the above variational characterization \eqref{poh} of $D(\b)_{1,1}$, for each $f\in L^\infty (\vpz)$ one has an upper bound of $D(\b)_{1,1}$. Our construction of the right functions $f$, giving efficient upper bounds, is inspired by the one in \cite{PR} (see Item (3) in
  \cite[Proof~of~Thm.~3.12]{PR}).
 
\medskip


Given a subset 
$A\subseteq\bbR^d$, we write  $\text{diam}\,A$ for the diameter of $A$ w.r.t. the uniform norm, i.e. $\text{diam}\,A:=\sup\{|x_i-y_i|\,:\,x,y\in A,\,i=1,\ldots,d\}$.  
\begin{Definition}%[Clusters $C[\zeta,\beta](x,\omega)$] 
 Given $x\in\bbR^d$, we denote by $C[\z,\b](x,\o)$   the connected component  containing  $x$ in the graph $\cG[\z,\b](\o)$.  If $x=0$, we simply  write  $C[\z,\b](\o)$.
 \end{Definition}

%
%In the proof of Theorem \ref{parataUB} the following percolation property   implied by  \cite[Theorom 1]{FagMim1} will play an important role:
%\begin{Fact}\label{patatine}\cite{FagMim1} Let $\bbP_0^*$ be the Palm distribution associated to a {\rm PPP}$(\l,\nu_{1,  \a})$ where $\l<\l_c (\a)$.
%Then
%there exists a positive constant $c=c(\l,\a)$  such that 
%\be
% \bbP^*_0\left( {\rm diam}\,C [1, 1] > n   \right) \leq e^{- c \,n } \qquad 
% \forall n \in \bbN\,.
% \en
% The same result holds if we replace $\nu_{1,  \a}, \l_c (\a)$ with $\nu^+_{1,  \a}, \l^+_c (\a)$, respectively.
% \end{Fact}
%


Given $n \in \bbN$ and given the parameters $\z,\b>0$, we consider the test function  $f_n=f_{n,\z,\b}\in L^\infty(\O_0)$ defined as \begin{equation}
\label{ftest}
f_n (\o):=
\begin{cases}
-\min \left\{ x_1:  x\in C[\z,\b](\o)\right\}  &\text{if } {\rm diam}\,C[\z,\b](\o) \leq n \,,
\\0 &\text{otherwise}\,.
\end{cases}\end{equation}
By \eqref{poh}, we have
\be\label{st1} 
2 D(\b)_{1,1} \leq \bbE_0 \Big[ \sum_{x\in \widehat \o}  c_{0,x}(\o,\b) \left
 (x_1 - \nabla_x f_n  (\o)\right)^2\Big]
 \leq A_1+A_2+A_3 \,,
\en
where
\begin{align}
& A_1:= \bbE_0 \Big[ \sum_{x\in \widehat \o}  c_{0,x}(\o,\b) \left
  (x_1 - \nabla _x f_n  (\o)\right)^2  \mathds{1} _{H_x}(\o)
  \mathds{1}_{
  \{ {\rm diam}\, C[\z,\b] \leq n \}} (\o)
  \Big]\,,\\
  & A_2:= \bbE_0 \Big[ \sum_{x\in \widehat \o}  c_{0,x}(\o,\b) \left
  (x_1 - \nabla _x f_n  (\o)\right)^2  \mathds{1} _{H_x}(\o)
  \mathds{1}_{
  \{ {\rm diam}\, C[\z,\b] > n \}}(\o)
 \Big]\,,\\
& A_3= \bbE_0 \Big[ \sum_{x\in \widehat \o}  c_{0,x}(\o,\b) \left
 (x_1 - \nabla_x f_n  (\o)\right)^2  \mathds{1} _{H^c_x}(\o) \Big]\,,
\end{align}
and $H_x:=\{\o \in \O_0\,:\,  x\in C[\z,\b](\o)\}$.


The terms $A_1$, $A_2$ and $A_3$ are similar to the terms $I_N^{(1)}(\b)$, $I_N^{(2)}(\b)$ and $I_N^{(3)}(\b)$ in the arxiv version of  \cite[Section~3]{FM}, respectively
\footnote{We point out an error in  the proof of \cite[Prop.~1]{FM}.  \cite[Eq.~(3.15)]{FM}   is valid only  by inserting in $I^{(1)}_N(\b)$ the characteristic function of the event $\{|C_0^\b(\xi)|\leq N\}$. One has then to treat another contribution given by $I^{(1)}_N(\b)$ with inside  the     characteristic function of the event $\{|C_0^\b(\xi)|> N\}$. This contribution is treated in the arxiv version of  \cite[Prop.~1]{FM}, and also in our proof of Lemma \ref{santo} (in \cite{FM} one  has  to use that $\nabla_x f_N^\b(\xi)=0$ in this case).}. By arguments similar to the ones  used there, one gets the following result:   
\begin{Lemma}\label{santo}  Fix  $\d\in (0,1)$ and a pair of conjugate exponents $p,q\in (1,+\infty)$.  Let $\bbP$ be a stationary marked SPP 
satisfying \eqref{non_periodico} and  $\bbE\bigl[ \widehat{\o}([0,1]^d) ^{\lceil q\rceil+1}\bigr]<+\infty$.
Then, for a suitable constant  $c=c( \d,p)>0$, it holds 
\be\label{francesco} D(\b)_{1,1} \leq c n^{-1/p}\bbE_0\big[ {\rm diam}\, C[\z,\b]\big]^{1/p} +
c
e^{-(1-\d)\z} (1+ \bbE_0 \bigl[ ( {\rm diam}\, C[\z,\b])^{2p}  \bigr] ^{1/p}\,. \
\en
\end{Lemma}
We sketch the proof (for completeness and due an  error in the proof of \cite[Prop.~1]{FM} already commented).
\begin{proof}
 When   $H_x $ occurs, we have    
$
C[\z,\b](\o)=C[\z,\b](x,\o)= C[\z, \b](\t_x \o) +x$
and therefore $
{\rm diam}\, C[\z,\b](\t_x\o) =  {\rm diam}\, C[\z,\b](\o) $. Hence, if in addition   ${\rm diam}\, C[\z,\b](\o) \leq n$, we get $x_1-\nabla_x f_{n}(\o) =0$, thus implying that $A_1=0$. On the other hand, if $H_x$ occurs and ${\rm diam}\, C[\z,\b](\o) > n$,  then we  have $f_n(\o)=f_n(\t_x\o)=0$, hence  $x_1-\nabla_x f_{n}(\o) =x_1$
and by H\"older's inequality 
 $A_2 \leq  \bbE_0 \big[ \bigl( \sum_{x\in \widehat \o}  e^{-|x|} x_1^2 \bigr)^q \big]^{1/q} 
 \vpz\bigl( {\rm diam}\, C[\z,\b] > n  \bigr)^{1/p}$ (the last expectation can be further bounded by Markov inequality).

 When $H_x^c$ occurs and $x\in \widehat{\o}\setminus\{0\}$, it must be   $c_{0,x}(\o,\b) < e^{-\z}$ (otherwise $\{0,x\}$ would be an edge of $\cG[\z,\b](\o)$). Since we always have $c_{0,x}(\o,\b) \leq e^{-|x|}$, we conclude that $c_{0,x}(\o,\b) \leq e^{-(1-\d)\z}e^{-\d|x|}$.
 As a consequence   we can bound
$
A_3 \leq  3 C_1+3C_2+3C_3$, where 
$C_1:=e^{-(1-\d)\z}\bbE_0 \big[ \sum_{x\in \widehat \o} e^{-\d|x|}  |x_1|^2\big]$, $C_2:=e^{-(1-\d)\z} \bbE_0 \big[ \sum_{x\in \widehat \o}  e^{-\d|x|} f_n(\t_x\o) ^2  \big]$ and 
$ C_3:=e^{-(1-\d)\z} \bbE_0 \big[ \sum_{x\in \widehat \o}  e^{-\d|x|} f_n(\o) ^2  \big]$. By \eqref{non_periodico} (and its equivalent version \eqref{non_periodico_bis}) and \cite[Lemma~1]{FM}--(i), $C_2=C_3$. To bound  $C_3$ we use H\"older's inequality and  that $|f_n|\leq  {\rm diam}\, C[\z,\b]$. We then   get
$C_3\leq e^{-(1-\d)\z}
\bbE_0 \big[ \bigl( \sum_{x\in \widehat \o}  e^{-\d|x|}\bigr) ^q  \big]^{1/q}
   \bbE_0 \bigl[ ( {\rm diam}\, C[\z,\b])^{2 p}  \bigr] ^{1/p}
   $. To conclude it is enough to use  the moment hypothesis on $\widehat{\o}([0,1]^d)$ and  \cite[Lemma~1]{FM}--(ii), which allows to bound by $c(\g,n)$ the $n$--th moment of $\sum_{x\in \widehat{\o}}c_{0,x} e^{-\g|x|}$ w.r.t. $\cP_0$ whenever the $(n+1)$--moment of  $\widehat{\o}([0,1]^d)$ w.r.t. $\cP$ is finite. In this way we can bound the expectations  $\bbE_0 \big[ \bigl( \sum_{x\in \widehat \o}  e^{-|x|} x_1^2 \bigr)^q \big] $,  $\bbE_0 \big[  \sum_{x\in \widehat \o}  e^{-|x|} x_1^2  \big] $ and $\bbE_0 \big[ \bigl( \sum_{x\in \widehat \o}  e^{-\d|x|}\bigr) ^q  \big]$.
\end{proof}

\begin{Remark}\label{bischero} Instead of the above analysis, we could have  used the bound (3.9) in \cite[Prop.~1]{FM}, applied by taking as set   \cite[Eq.~(3.1)]{FM} the edge set  of $\cG [\z,\b]$ and $\z=\ell(\b)$ for a suitable function $\ell(\b)$. On the other hand, we found that fixing now $\ell(\b)$ and introducing the notation of \cite{FM} would have been less trasparent. Note that, in any case, the final choice of the test functions differs from the one in the proof of \cite[Theorem~1]{FM} (see in particular \cite[p.~276]{FM}).
  \end{Remark}












For the next result recall  the definition \eqref{box_pokemon}  of the configuration $\o _{\z,\b}\in\O$     for $\o \in \O$  and $\z,\b>0$, i.e. 
\begin{equation*}
\begin{split}
 \o_{\z,\b}:& =\big\{ \big ( x / \z, \b E_x/\z \big) \,: \,  x\in \widehat{\o}_{\z/\b} \big \}=\big\{ \big ( x/\z, \b E_x/\z\big) \,: \, x\in\widehat \o  \,, |E_x| \leq \z /\b \big\} \,.
\end{split}
\end{equation*}
 Recall that $\o_{\z, \b}$ equals   \eqref{gandhi} 
when $\z=\ell(\b)$. \rosso{Recall also that $\nu(\g)=\nu([-\g,\g])$ and that  $\nu_{\star,\g}$ is the law of  $X/\g$ conditioned to the event $|X/\g|\leq 1$ when $X$ is a random variable with law $\nu$ (cf. respectively \eqref{alloro} and \eqref{sette2}).}


 \begin{Lemma}\label{serafico} Let $\bbP$  be a stationary marked SPP with finite and positive intensity \orc{$\rho$}.  Let  $\z,\b>0$ and let  $\bar \bbP$ be the law of the configuration  $\o_{\z,\b}$   when $\o$ is sampled according to $\bbP$. 
Suppose that $\vpz(\cD)>0$, where
\[ \cD:=\{\o\in \O_0\,:\,  0
 \in \widehat{\o}_{\z/\b} \}=\{\o\in \O_0\,:\,   |E_0|\leq \z/\b\}  \,.
 \] Then:
\begin{itemize}
\item[(i)]  $\bar\bbP$ is a stationary marked  SPP with finite and positive  intensity; 
\item[(ii)] for any  nonnegative Borel function $h$ on $\bbR$  with $h(0)=0$ it holds 
 \be\label{zoomo}
\bbE_0 \left[ \, h \left({\rm diam}\, C[ \z, \b]  \right) \,\right]=  \vpz (\cD)    \bar \bbE_0 \left[ \, h \left(\z\, {\rm diam}\,C[ 1, 1] \right)\, \right] 
  \en
  where $\bar\bbE_0$ is  the expectation w.r.t. the  Palm distribution  $\bvpz$ associated to $\bar \bbP$.
  \end{itemize}
  
  Moreover, if $\bbP=\hat\bbP_\nu$,  then $\vpz(\cD)=\nu( \z/\b)$ and $\bar \bbP$ equals 
   $(\hat \bbP_p\circ \Psi_{\z}^{-1})_{\nu_{\star,  \z/\b}}$. If in addition $\bbP={\rm PPP}[\rho,\nu]$, then $\bar \bbP$ equals $ {\rm PPP}[\rho \nu(\z/\b)\z^d   ,\nu_{\star,  \z/\b}] $.
   
      \end{Lemma}
 Note that by Item (i) the Palm distribution $\bvpz$ associated to $\bar\bbP$ and appearing in Item (ii)  is well defined.
 \begin{proof}
We first observe that 
 \be\label{ho_sete_78}
  \begin{split}
   \bbE_0 \left[ h\bigl( {\rm diam}\,  C[\z,\b] \bigr)    \right]& =
     \bbE_0 \left[ h\bigl( {\rm diam}\,  C[\z,\b] \bigr)  \rosso{\mathds{1}_\cD} \right]\\
     & = \vpz(\cD)  \bbE_0 \left[ h\bigl( {\rm diam}\, C[\z,\b] \,\bigr)  |\, \cD \right] \,.
     \end{split} 
     \en
     Indeed, the first identity follows from the hypothesis that $h(0)=0$ and the fact that  $ C[\z,\b] (\o) =\{0\}$ if   $|E_0 | >\z/\b$  (see \orc{Lemma}~\ref{fadeev_new}--(i)).
The second identity is just a  conditioning, which is allowed by the hypothesis that \orc{$\vpz(\cD)>0$}.

%$$ ***$$     
%
%By  Campbell's formula \eqref{marlena}, by conditioning on $\widehat{\o}$  and by the definition of $\nu$--randomization  we have 
% \[
% \bbP_0 ( \cD)
% %= \frac{1}{m}\bbE \Big[ \sum _{x\in \hat \o \cap [0,1]^d}  \mathds{1} ( \t_x \o \in \cD)\Big]
% =  \frac{1}{m}\bbE \Big[ \bbE\big[ \sum _{x\in \hat \o \cap [0,1]^d}  \mathds{1} ( \t_x \o \in \cD)\,|\, \hat \o \big]\Big]= \frac{\nu (\z/\b) }{m}\bbE \big[ \widehat{\o} \big( [0,1]^d\bigr) \big]= 
% \nu (\z/\b) \,. \]
% As $\bbP_0 ( \cD)=\nu (\z/\b) $, \eqref{ho_sete_78} implies  \eqref{ho_sete}, thus proving our claim.
We call 
 $\vvv{\bbQ}$ the law of 
  $\o _{\z/\b}$  when $\o$ is sampled according to $\bbP$. By Lemma 
  \ref{quinoa}  $\vvv{\bbQ}$ is a stationary marked  SPP with finite and positive  intensity. We call $\vqz$ the Palm distribution associated to $\vvv{\bbQ}$.
 By  Lemma \ref{quinoa} $\vqz$ equals the law of $\o _{\z/\b}$ when $\o$ is sampled according to $\vpz(\cdot | \orc{\cD} )$. Since in addition (by \orc{Lemma}~\ref{fadeev_new}--(i))  $\text{diam}\, C[\z,\b] (\o)= \text{diam}\, C[\z,\b] (\o_{\z/\b})$ if $\o\in \cD$, we conclude that 
  \be \label{ho_sete_1} 
 \bbE_0 \left[ h\bigl( \text{diam}\, C[\z,\b]  \bigr) | \cD \right] =\int_{\O_0} 
  d\vqz(\o) 
 h \left( \text{diam}\, C[\z ,\b](\o)\right)\,.
 \en 
Note that  $\bar \bbP$ is the image of $ \vvv{\bbQ}$ when applying the point-mark transformation $(x,E_x)\mapsto  \bigl ( x/\z, \b E_x/\z \bigr) $. It then follows that 
$\bar \bbP$ is a stationary SPP with finite and positive  intensity, and that 
$\bvpz$ is the image of $ \vqz$ by the above transformation. To conclude we observe that, by \orc{Lemma} \ref{fadeev_new}, 
the bijection 
\be\label{vettorino}
\o_{\z/\b}=\{(x,E_x)\,:\, x\in \widehat{\o} _{\z/\b} \}\mapsto \{(x/\z, \b E_x/\z)\,:\,x\in \widehat{\o} _{\z/\b}\} \in \o_{\z,\b}
\en
induces naturally a bijection between $\cG[\z,\b](\o_{\z/\b})$ and $\cG[1,1](\o_{\z,\b})$.
 By combining this bijection  with the observation that  $\bvpz$ is the image of $ \vqz$ by the  transformation $(x,E_x)\mapsto  \bigl ( x/\z, \b E_x/\z \bigr) $, we then conclude that 
 \be\label{smeraldo} 
 \int_{\O_0} 
  d\vqz\orc{(\o)}
 h \left( \text{diam}\, C[\z ,\b](\o) \right)
 = \int_{\O_0} d\bvpz(\o) 
 h \left(\z \text{diam}\, C[1 ,1](\o) \right)
 \,.
\en
By combining \eqref{ho_sete_78}, \eqref{ho_sete_1} 
and \eqref{smeraldo}  we then get  
 \eqref{zoomo}.
 % Identities \eqref{zoomo1} and \eqref{zoomo2} are an immediate application of \eqref{zoomo}.
 
  We now move to the last  statement. Let $\bbP:=\hat\bbP_\nu$. Then by  Campbell's formula \eqref{marlena}, by conditioning on $\widehat{\o}$  and by the definition of $\nu$--randomization,  we have 
 \[
 \vpz ( \cD)
 =  \frac{1}{\orc{\rho}}\bbE \Big[ \bbE\big[ \sum _{x\in \widehat \o \cap [0,1]^d}  \mathds{1} ( \t_x \o \in \cD)\,|\, \widehat \o \big]\Big]= \frac{\nu (\z/\b) }{\orc{\rho}}\bbE \big[ \widehat{\o} \big( [0,1]^d\bigr) \big]= 
 \nu (\z/\b) \,. \]
This proves that $ \vpz ( \cD)=\nu(\z/\b)$. The rest of the final statement follows  
   from \orc{Lemma \ref{fadeev_250}}.   \end{proof}



 



Let us now take $\nu \in \cvp \cup  \cv$. Recall the notational Warning \ref{pericolo} for the definition of $C_0,\e$. We also take $\z$ equal to the quantity $\ell(\b)$ defined in \eqref{canditi200}:
\be\label{zoomo3}
\z=\ell(\b):= (\l /\rho)^{\frac{1}{\a+1+d}}(C_0 \b) ^{\frac{\a+1}{\a+1+d}}\,.
\en
Note that $\z\leq \min\{C_0, \e\}\b$ for $\b $ large. In particular, by taking  $\bbP={\rm PPP}[\rho, \nu]$, for $\beta $ large we have \orc{(see~Proposition \ref{paradiso_PPP})}
\be\label{zoomo4}
\bar\bbP= {\rm PPP}[\rho \nu(\z/\b)\z^d   ,\nu_{\star,  \z/\b}]  = 
\begin{cases}
 {\rm PPP}[\l   ,\nu_{1,\a}^+]   &\text{ if } \nu\in \cvp\,,\\
{\rm PPP}[\l   ,\nu_{1,\a}]   & \text{ if } \nu \in \cv\,.
 \end{cases}
\en
%Indeed, the first identity follows from Lemma \ref{serafico}, while the  second one follows from  Lemma \ref{fissato} and our particular choice of $\z$.

We now need a fundamental ingredient from percolation:
%\verde{controllare e spiegare, penso che \cite{FagMim1} sia piu' generale}
\begin{Fact}[A.~Faggionato, A.H.~Mimun \cite{FagMim1}]\label{patatine} Let $\bvpz$ be the Palm distribution associated to $\bar \bbP:= {\rm PPP}[\l,\nu_{1,  \a}]$ with $\l<\l_c (\a)$ (subcritical case).
Then
there exists a positive constant $c=c(\l,\a)$  such that 
\be
\bvpz\left( {\rm diam}\,C [1, 1] > n   \right) \leq e^{- c \,n } \qquad 
 \forall n \in \bbN\,.
 \en
 The same result holds if we replace $\nu_{1,  \a}, \l_c (\a)$ with $\nu^+_{1,  \a}, \l^+_c (\a)$, respectively.
 \end{Fact}
The above fact follows by applying \cite[Thm.~1.4]{FagMim1} together with the observation that the critical intensities   $\l_c(\a)$, $\l_c^+(\a)$ coincide with the one in \cite[Eq.~(1.5)]{FagMim1} as proved in Appendix \ref{liberato} (see also the discussion above Theorem 1.4 in \cite{FagMim1} for the check of the assumptions in our case).
% We now comment why the assumptions in \cite[Thm.~1.4]{FagMim1}   are satisfied (the case with $\nu_{1,\a}$ is indeed already discussed in \cite{FagMim1}).  Let us take  i.i.d. random variables $E,E'$ with common distribution $\nu =\nu_{1,\a}$ or $\nu=\nu^+_{1,\a}$. Then the function $\mathfrak{h}(u):=P(|E|+|E'|+|E-E'|\leq 1-u )$ defined on $(0,+\infty)$ is positive on $(0,1)$ and is zero on $[1,+\infty)$.
%Moreover, on (0,1) it is also uniformly continuous as the product measure $\nu\otimes \nu$ has bounded density w.r.t. the Lebesgue measure. As a consequence the function $\mathfrak{h} $ is good according to \cite[Definition~1.3]{FagMim1}, which is the assumption of  \cite[Thm.~1.4]{FagMim1}. 

\smallskip

We have now all the tools to conclude the proof of Theorem \ref{parataUB}. Take \orc{$\nu\in \cvp \cup \cv$ and $\bbP= {\rm PPP}[\rho,\nu]$} as in the theorem.
Recall Definition \eqref{nuziale} of $\l_c^*(\a)$.
It is enough to prove \eqref{preUB}, since \eqref{UB} is a byproduct of \eqref{nanna} in Fact \ref{teo_eff_cond} and \eqref{preUB}.
We point out that ${\rm PPP}[\rho,\nu]$ satisfies \eqref{non_periodico} and $\widehat{\o}([0,1]^d)$ has finite moments of all orders.
We can therefore  apply Lemma \ref{santo} with $p=2$. By \eqref{zoomo} in Lemma 
\ref{serafico},  for  $s\in \{1,4\}$  we have  $
\bbE_0\big[ \bigl({\rm diam}\, C[\z,\b]\bigr)^{s} \big]  \leq   \z^{s} \bar\bbE_0\big[ \bigl({\rm diam}\, C[1,1]\bigr)^{s}  \big]$.
We take $\z=\ell(\b)$ as in \eqref{zoomo3}, with $\l<\l_c^*(\a)$.
Then, 
by \eqref{zoomo4} and  Fact \ref{patatine},  the last expectation is finite.
We conclude (see \eqref{francesco})  that for $\d\in (0,1)$ there exists $c>0$ independent from $n$ and $\b$ such that
\[
  D(\b)_{1,1} \leq c n^{-1/2}  \ell(\b)^{1/2}+ c
  e^{-(1-\d)\ell(\b)} (1+ \ell(\b)^2)\,.
  \]
By taking the limit $n\to+\infty$ we get $ D(\b)_{1,1} \leq  c
  e^{-(1-\d)\ell(\b)} (1+ \ell(\b)^2)$. Hence




% \smallskip

% We first use Lemma \ref{montemonaco1} with $p=2=q $  to bound $A_1$. We note that, for each $\d\in (0,1)$, the constants $\k_1(\d)$ and  $k_2(\d, q)$  are finite as $\hat \bbP$ is a PPP. This can be easily checked. Indeed,  to upper bound the expectations one can divide $\bbR^d$ in disjoint boxes $\L_z:=z+[-1/2,1/2)^d$, $z\in \bbZ^d$,  and estimate $\sum_{x\in \hat \o } e^{- c_1 |x|} |x|^{c_2} $ by $C \sum_{z\in\bbZ^d}  e^{- c_1 |z|} |z|^{c_2}
% \hat \o (\L_z)$ and afterwards use that $\big\{ \hat \o(\L_z) \,:\,z\in \bbZ^d\big\}$ are i.i.d. Poisson random variables with mean $\rho$.

% We take $\z=\ell(\b)$ as in \eqref{zoomo3}, with $\l<\l_c^*(\a)$. Then, by applying Lemma \ref{montemonaco1}  and afterwards \eqref{zoomo1} and $\bbP_0(\cD)\leq 1$ for the second bound, we get
% \begin{equation*}
% \begin{split}A_1 &\leq
%  3\k_1(\d) e^{-(1-\d)\ell(\b)}+ 6 \k_2(\d, q)   e^{-(1-\d)\ell(\b) }
%  \bbE_0\big[ \bigl({\rm diam}\, C[\ell(\b) ,\b]\bigr)^{2 p} \big]^{1/p}\\
%  &\leq
%  3\k_1(\d) e^{-(1-\d)\ell(\b)}+ 6 \k_2(\d, q)   e^{-(1-\d)\ell(\b) }  \ell(\b)^2
%  \bar \bbE_0\big[ \bigl({\rm diam}\, C[1 ,1]\bigr)^{2 p} \big]^{1/p}\,,
% \end{split}
% \end{equation*}
% where $\bar \bbE_0$ is the expectation w.r.t. the Palm distribution $\bvpz$ associated to $\bar \bbP$ given by \eqref{zoomo4} \orc{for $\b$ large, as we assume}. 

% Since  $\l<\l_c^*(\a) $ by Fact \ref{patatine} we get that the constant  $\bar \bbE_0\big[ \bigl({\rm diam}\, C[1 ,1]\bigr)^{2 p} \big]^{1/p}$ is finite and determined only  by $\l,\a$. 
% Hence, we have obtained that
% \be\label{soffro1}
% A_1 \leq  C_1 e^{-(1-\d)\ell(\b) } \ell(\b)^2
% \en
% where the constant $C_1$ does not depend from  $n,\beta$.

% \smallskip

% We now use Lemma \ref{montemonaco2} to bound $A_2$. We take $\z$ as above.  As \rosso{done} for $\k_1(\d)$ and $\k_2(\d,p)$, one easily \rosso{proves} that the constant $\k_3:=\bbE_0 \big[ \sum_{x\in \widehat \o} e^{-|x|}  x_1^2\big]^{1/2} $ is finite. By combining Lemma \ref{montemonaco2}, \eqref{zoomo2} and Fact \ref{patatine} (since $\l<\l_c^*(\a) $) we get 
% \be\label{batteria}
% A_2\leq  \k_3 \,
% \vpz\bigl({\rm diam}\,C[\z,\b] > n \z \bigr)^{1/2}
% \leq \k_3 \,
% \bvpz\bigl({\rm diam}\,C[1,1] >  n \bigr)^{1/2}\leq \k_3 e^{- c n}
% \,.
% \en
% By combining the above bound with \eqref{soffro1} we get $2 D(\b)_{1,1} \leq C_1 e^{-(1-\d)\ell(\b) } \ell(\b)^2+\k_3 e^{- c n}$, where the positive constants $C_1,\k_3, c$ do not depend 
%  from  $n,\beta$. By the arbitrariness of $n$ and  taking $n\to +\infty$ we then get   $2 D(\b)_{1,1} \leq C_1 e^{-(1-\d)\ell(\b) } \ell(\b)^2$. We recall that $\ell(\b)=(\l /\rho )^{\frac{1}{\a+1+d}}(C_0 \b) ^{\frac{\a+1}{\a+1+d}}$ with $\l<\l_c^*(\a)$. Then 
% \begin{multline*}
%\[
%\begin{split}
\[ \limsup _{\b\to +\infty}  \b^{- \frac{1+\a}{d+1+\a} } \ln D(\b)_{1,1} 
\leq
-(1-\d) (\l C_0^{\a+1} /\rho )^{\frac{1}{\a+1+d}} \,.
\]
By taking $\d$ and $\l$ arbitrarily near to $0$ and $ \l_c^*(\a)$ respectively, we finally get \eqref{preUB} by the expression \eqref{premium} of $\chi(\rho,\nu)$.

%%%%%%%%%%%%%%%%%%%%%%%%%%%%%%%%%%%%%%%%%%%%%%%%%%%%%%%%%%%%%%%%%%%%%%%%%%%%%%%%%%%%%%%%%%%%%%%%%%%%%%%%%%%%%%%%%%%%%%%%%%%%%%%%%%%%%%%%%%%%%%%%%%%%%%%%%%%%%%%%%%%%%%%%%%%%%%%%%%%%%%%%%%%%%%%%%%%%%%%%%%%%%%%%%%%%%%%%%%%%%%%%%%%%%%%%%%%%%%%%%%%%%%%%%%%%%%%%%%%%%%%%%%%%%%%%%%%%%%%%%%%%%%%%%%%%%%%%%%%%%%%%%%%%%%%%%%%%%%%%%%%%%%%%%%%%%%%%%%%%%%%%%%%%%%%%%%%%%%%%%%%%%%%%%%%%
\section{Proof of Theorem  \ref{parataLB} (lower bound in Mott's law)}\label{sec_LB}
Before giving the proof of Theorem \ref{parataLB}  we collect some preliminary results, keeping the discussion completely general.



We can extend the definitions introduced in  Section \ref{fumata} to general  graphs. More precisely, 
given a weighted non--oriented graph $\bbG=(\bbV,\bbE)$ with weight function $C: \bbE\to [0,+\infty)$ and vertex set $\bbV \subset \bbR^d$ and given $\ell>0$, we consider the resistor network on $S_\ell$ with nodes $x\in \bbV \cap S_\ell$, edges $\{x,y\}$ with $x,y\in \bbV \cap S_\ell$ and $\{x,y\}\cap \L_\ell \not = \emptyset$ and electrical conductivity of the edge  $\{x,y\}$ given by $C(x,y)$.
% kuka. forse non e' cosi' necessario  We assume that  $\sum_{y:\{x,y\}\in \bbE} C(x,y)<+\infty$ for all  $ x\in \bbV$.
We call $\s_\ell(\bbG, C)$ the conductivity along the first direction when the electrical potential $V$ has values $1$ and $0$ on $S_\ell^-$ and $S_\ell^+$, respectively. We point out that $V$ is not univocally defined on the connected  components of the resistor network not intersecting $S_\ell^-\cup S_\ell^+$. Indeed, there we can only say that $V$ is constant. Since the physics (and in particular $\s_\ell(\bbG, C)$) depends only  on the increments of $V$ (which are  univocally determined), without loss we can impose that  $V\equiv 0$ on the above   connected components.


As in the derivation of \eqref{papero}, one can easily proves Dirichlet's  principle (cf.~\cite[Exercise~1.3.11]{DS}):
% kuka, controllare terminologia
\be\label{papero_bis}
\s_\ell (\bbG,C )=\inf \left\{ \cD_\ell(u) \,\big{|} \, u: \bbV \cap S_\ell \to [0,1]\,, u_{|\bbV \cap S_\ell^-} \equiv 1
\,, \; u_{|\bbV \cap S_\ell^+}\equiv0  \right\}\,,
\en
where 
\begin{align*}
&  \cD_\ell(u) := \sum _{\{x,y\} \in \bbE_\ell} C(x,y) \left[ u(y)- u(x) \right]^2\,,\\
& \bbE_\ell:=\bigl\{\,\{x,y\}\in \bbE\,:\, x \in  \bbV \cap \L_\ell\,,\; y \in \bbV \cap S_\ell, \; x\not=y\bigr\}\,.
\end{align*}
As a consequence,   $\s_\ell (\bbG,C)$ satisfies  Rayleigh's monotonicity law (cf.~\cite[Section~1.4]{DS}): 
\be\label{rayleigh_mon_law}
\bbG\supset \bbG' \text{ and } C\geq C'\;\;\Rightarrow \;\; \s_\ell(\bbG, C)\geq \s_\ell(\bbG', C')\,.
\en



%%%%%%%%%%%%%%%%%%%%%%%%%%%%%%%%%%%%%%%%%

\medskip

Recall that, given a graph $\bbG=(\bbV,\bbE)$ with vertex set  $\bbV\subset \bbR^d$,   $\cN_\ell(\bbG)$ denotes  the  maximal number of  vertex-disjoint LR crossings of $\L_\ell$ in $\bbG$ (cf. Definition \ref{def_LR}).

\begin{Lemma}\label{lemma_LB_chiave}
Given a graph $\bbG=(\bbV,\bbE)$ with vertex set  $\bbV\subset \bbR^d$, we have\be\label{LB_chiave}
\s_\ell (\bbG, 1) \geq \frac{\cN_\ell(\bbG) ^2 }{2\cN_\ell(\bbG)+ |\bbV \cap \L_\ell|}\geq \frac{\cN_\ell(\bbG) ^2 }{3 |\bbV \cap \L_\ell|} \,.
\en
\end{Lemma}
% kuka \rosso{commentare existing results su cui si basa la proof. kesten? forse va precisato che le cose sono localmente finite}

\begin{proof} Let $n:= \cN_\ell (\bbG)$.
Let $u : \bbV\cap S_\ell \to [0,1]$ be a function with $u_{ |\bbV \cap S_\ell^-}\equiv 1$ and $u_{|\bbV \cap S_\ell^+}\equiv 0$. We set $\nabla_{x,y}u :=u(y)-u(x)$.  We fix a family  $\p_1,\ldots,\p_n$ of  vertex--disjoint LR crossings of $\L_\ell$ in $\bbG$. Recall that $\pi_i $ is a sequence $(x_1, x_2, \dots, x_k)$ of $k\geq 3$ distinct points in $\bbV$ such that 
$\{x_r,x_{r+1}\}\in \bbE$ for all $r\in\{1,\dots, k-1\}$, 
$x_1 \in \cS_\ell^-$, $x_k \in \cS_\ell^+$ and $x_r\in \L_\ell  $ for $1< r<k$. Below, we write $\{x,y\}\in \pi_i$ if $\{x,y\}$ is an edge of $\pi_i$, i.e. $\{x,y\}=\{x_r, x_{r+1}\}$ for some some $r$. We also set $|\p_i|:= k$. 
 

As the $\p_i$'s are vertex-disjoint, they are also edge-disjoint. Hence, by Rayleigh's monotonicity law (cf.~\eqref{rayleigh_mon_law}),  we can bound 
\begin{equation}
\label{low1}
\s_\ell (\bbG,1) \geq \sum_{\{x,y\}\in \bbE_\ell}  (\nabla_{x,y} u)^2 \geq\sum_{i=1}^n\sum_{\{x,y\}\in \p_i}(\nabla_{x,y} u)^2 
\,.
\end{equation}
By Cauchy-Schwarz's inequality we obtain  $1=\Big(\sum_{\{x,y\}\in \p_i}\nabla_{x,y} u\Big)^2$ is upper bounded by $ |\p_i|\sum_{\{x,y\}\in \p_i} (\nabla_{x,y} u)^2$.
%$$1=\Big(\sum_{\{x,y\}\in \p_i}\nabla_{x,y} u\Big)^2\leq |\p_i|\sum_{\{x,y\}\in \p_i} (\nabla_{x,y} u)^2.$$
Hence the r.h.s. in \eqref{low1} can be bounded from below as
\begin{equation}
\label{low2}
\begin{split}
\sum_{i=1}^n\sum_{\{x,y\}\in \p_i} (\nabla_{x,y} u)^2\geq  \sum_{i=1}^n\frac{1}{|\p_i|}.
\end{split}
\end{equation}
By Jensen's inequality  applied to the convex function $ x \mapsto 1/x$ for $x>0$,  we can bound from below the r.h.s. of \eqref{low2} as
\begin{equation}
\label{low3}
\begin{split}
 \sum_{i=1}^n\frac{1}{|\p_i|}&=n \cdot\frac{1}{n}\sum_{i=1}^n\frac{1}{|\p_i|}\geq n\left(\frac{1}{n}\sum_{i=1}^n|\p_i|\right)^{-1}
=\frac{n^2 }{\sum_{i=1}^n|\p_i|}
\end{split}
\end{equation}
(the above use of Jensen's inequality is not new,  see e.g.~the proof of  \cite[Prop.~3.2]{CC}).
Recall  that the $\pi_i$'s are vertex-disjoint  and  all vertexes of each $\pi_i$ are in $\bbV \cap \L_\ell$ apart the extreme ones.
As a consequence, we can bound  $\sum_{i=1}^n|\p_i|\leq  2n +| \bbV \cap \L_\ell|$.
 Combining this bound with \eqref{low1}, \eqref{low2} and \eqref{low3} we get  the first bound in \eqref{LB_chiave}. For the second one it is enough to observe  that the $\pi_i$'s are vertex-disjoint and that each $\pi_i$ has at least one point in $\bbV \cap \L_\ell$.
\end{proof}
%%%%%%%%%%%%%%%%%%%%%%%%%%%%%%%%%%%%%%%%%



%Recall the definition of $\Psi_\z \cG[\z,\b] (\o_{\z/\b})   $ in Proposition \ref{fadeev_new}--(iv).
%%%%%%%%%%%%%%%%%%%%%%%%%%%%%%%%%%
% kuka a volte ho messo $\b>0$ negli enunciati... Sotto $\o$ deve cmq essere tale che posso defire la conduttivita'

%\club We point out that the graph $\cG[\z,\b](\o)$ has local finite degree for ant $\o \in \O$. Indeed, $\o$ is a locally finite set in $\bbR^d$ and 


\begin{Lemma}\label{lasagne25}
Let $\z>0$ and $\b, \ell \geq 0$. Then  for any $\o\in \O$ it holds 
\be\label{pasta_rossa}
\ell^{2-d}\s_\ell(\o, \b)\geq  e^{-\z} \z^{2-d} L ^{2-d} \s_{L} ( \cG[1,1](\o_{\z,\b}  ),1)\,, \qquad  L:= \ell/\z\,,
\en
where $\o_{\z,\b}$ is defined in \eqref{box_pokemon}.
\end{Lemma}
%%%%%%%%%%%%%%%%%%%%%%%%%%%%%%%%%%
\begin{proof}
%%%%%%%%%%%%%%%%%%%%%%%%%%%%%%%%%%
We denote by $\cE[\z,\b]$ the edge set of the graph $\cG[\z,\b]$. By definition of $\cG[\z,\b]$ we have
\be \label{arca0}
c_{x,y}(\o, \b) \geq e^{-\z} \mathds{1} \bigl( \{x,y\} \in \cE[\z,\b](\o)\bigr)\,.
\en
We claim that 
\be \label{arca1}
\s_\ell\bigl(\o, \b\bigr) \geq e^{-\z} \s_\ell \bigl(\cG[\z,\b](\o) , 1\bigr)=   e^{-\z} \s_\ell \bigl(\cG[\z,\b](\o_{\z/\b} ) , 1\bigr) \,.
\en
Indeed, the  bound follows from Rayleigh's monotonicity law and \eqref{arca0}, while the identity follows from   \orc{Lemma} \ref{fadeev_new}--(i): points of $\widehat{\o}\setminus \widehat{\o}_{\z/\b}$ are isolated in the graph $\cG[\z,\b] $ and therefore they do not contribute to the conductivity. Recall now the  bijection between $\cG[\z,\b](\o_{\z/\b})$ and $\cG[1,1](\o_{\z,\b})$ given by \orc{Lemma} \ref{fadeev_new}--(ii). This bijection and \eqref{papero_bis} imply that 
$\s_\ell \bigl(\cG[\z,\b](\o_{\z/\b} ) , 1\bigr)= \s_{\ell/\z} \bigl(\cG[1,1](\o_{\z,\b} ) , 1\bigr)$. 
\orc{By combining  this identity with} \eqref{arca1}  we get \eqref{pasta_rossa}.
%  $\cG(\o)$ is obtained from  graph $\tilde \cG(\o)$
%by adding isolated points, which do not contribute to the conductivity. Using in addition 
%the bijection given in Proposition \ref{fadeev_new}--(ii)  where $\Psi(x):= x/\z$, we conclude that 
%\be \label{arca2}
% \s_\ell \bigl(\cG(\o) , 1\bigr)= \s_\ell \bigl(\tilde \cG(\o) , 1\bigr)=\s_{\ell/\z} \bigl( \bbG_\o,1 \bigr)\,, \qquad \bbG_\o:=\Psi_* \tilde \cG(\o)\,,
% \en
%and the law of $\bbG_\o$ equals the law of  
% $\cG\left [1, 1,\l(\z/\b, \rho,\nu)\z^d   ,\nu_{\star,  \z/\b}\right](\o)$ by sampling $\o$ under $\bbP$.
%%%%%%%%%%%%%%%%%%%%%%%%%%%%%%%%%%
\end{proof}
%%%%%%%%%%%%%%%%%%%%%%%%%%%%%%%%%%




%As a combination of Lemma \ref{poissoniano}, Lemma \ref{frittella} and Fact \ref{cortocircuito} we get the following:
\rosso{Using also  Fact \ref{cortocircuito} (which states that ${\rm PPP}[\l, \nu_{1,\a}^+]$ has a good LR crossing distribution), we can  finally give the proof of Theorem  \ref{parataLB}.}
%%%%%%%%%%%%%%%%%%%%%%%%%%%%%%
\subsection{Proof of Theorem \ref{parataLB}} \rosso{Take $\nu \in \cvp\cup \cv$.
 Due to \eqref{nanna} it is enough to prove the bound \eqref{LB}.}
 
We take $\z$ as in \eqref{zoomo3}, i.e.~$\z=\z(\b):=(\l /\rho)^{\frac{1}{\a+1+d}}(C_0 \b) ^{\frac{\a+1}{\a+1+d}}$ but, differently from the proof of Theorem \ref{parataUB}, we now take $\l>\l_c^*(\a)$. Recall that $\bbP={\rm PPP}[\rho, \nu]$.
 Let $\bar\bbP$ be   the law of $\o_{\z,\b}$ when $\o$ is sampled according to $\bbP$ and let  $\bar \bbE$ be the expectation associated to $\bar \bbP$.
As already observed in \eqref{zoomo4},  for $\b$ large
  $\bar \bbP$ equals
$ {\rm PPP}[\l   ,\nu_{1,\a}^+]  $ if $ \nu\in \cvp$ and equals ${\rm PPP}[\l   ,\nu_{1,\a}]  $ if $ \nu \in \cv$.


 By Lemmas  \ref{lemma_LB_chiave}  and  \ref{lasagne25}, setting  $ L:= \ell/\z$ we have 
 \be\label{pasta_rossa_bis}
 \begin{split}
\ell^{2-d}\s_\ell(\o, \b)& \geq  e^{-\z } (\z L) ^{2-d}  \s_{L} ( \cG[1,1](\o_{\z,\b}  ),1)
\geq  e^{-\z }(\z L) ^{2-d} \frac{\cN_L\left( \cG[1,1](\o_{\z ,\b}  )\right)^2}{ 3 | \widehat{\o_{\z  ,\b} } \cap \L_L|}\,.
\end{split}
\en



We first take the limit $\ell \to +\infty$ thinking $\b$ as fixed (trivially, $L\to +\infty$).  Since $\l>\l_c^*(\a)$, by Fact \ref{cortocircuito} concerning $\bar \bbP={\rm PPP}[\l   ,\nu^+_{1,\a}]  $  (in case $\nu\in \cvp$) and by the assumption that $\bar \bbP={\rm PPP}[\l   ,\nu_{1,\a}]  $ has a good LR crossing distribution (in case $\nu\in \cv$), we conclude that for $\b$ fixed  $\bbP$--a.s.
$\cN_L\left( \cG[1,1](\o_{\z ,\b}  \right) \geq c L^{d-1}$ for $L$ large and for a suitable constant $c>0$ (in the case $\nu\in \cv$ it could be $c=c(\o)$ a priori, but this is not a problem). On the other hand, by the ergodic theorem we have that $(2L)^{-d} | \widehat{\o_{\z ,\b} } \cap \L_L|$ converges to the intensity $\l$ of $\bar \bbP$   as $\ell \to +\infty$.
We then get that $\bbP$--a.s. 
\be\label{pasta_bianca}
\liminf_{\ell\to +\infty} L ^{2-d} \frac{\cN_L\left( \cG[1,1](\o_{\z ,\b}  )\right)^2}{ 3 | \widehat{\o_{\z ,\b} } \cap \L_L|}\geq \liminf_{L\to +\infty} L ^{2-d} \frac{c^2 L^{2d-2}}{ 3 \l  (2L)^d }=\frac{c^2}{3 \l 2^d}\,.
\en
By combining \eqref{pasta_rossa_bis} and \eqref{pasta_bianca} we get 
\be\label{rispetto}
\s(\b)=\lim_{\ell\to +\infty}\ell^{2-d}\s_\ell(\o, \b) \geq e^{-\z  } \z ^{2-d}  c^2 (3 \l 2^d)^{-1}\,.
\en
Since $\z=(\l /\rho)^{\frac{1}{\a+1+d}}(C_0 \b) ^{\frac{\a+1}{\a+1+d}}$
from \eqref{rispetto} we get 
$
\liminf _{\b \uparrow \infty}  \b^{-\frac{\a+1}{\a+1+d}} \ln \s(\b)  \geq -(\l  C_0^{\a+1}/\rho)^{\frac{1}{\a+1+d}}$, thus concluding the proof of Theorem \ref{parataLB}.


%%%%%%%%%%%%%%%%%%%%%%%%%%%%%%%%%%%
\appendix

%\section{Lemma \ref{transitare}}\label{proof_transitare}
%%%%%%%%%%%%%%%%%%%%%%%%%%%%%%%%%%%%
%%%%%%%%%%%%%%%%%%%%%%%%%%%%%%%%%%%%
%%%%%%%%%%%%%%%%%%%%%%%%%%%%%%%%%%%%
%\begin{Lemma}\label{transitare} For $d\geq 2$ 
%there exists a critical threshold  $\z_c=\z_c(\b,\rho,\nu)\in (0, +\infty) $  for which \eqref{critico} holds.
%\end{Lemma}
%%%%%%%%%%%%%%%%%%%%%%%%%%%%%%%%%%%%
%\begin{proof} The proof follows by coupling arguments and comparison with the Poisson Boolean model \cite{MR}. To simplify the notation we write  $\theta(\z)$ instead of $\theta(\z,\b, \rho, \nu)$ since $\b, \rho,\nu$ are fixed. We have already observed that $\theta(\z)\in \{0,1\}$.
%By a standard coupling argument one gets that $\theta(\z) \leq \theta (\z')$ if $\z \leq \z'$. Hence it is enough to show that (i) $\theta(\z)=0 $ for some $\z>0$ and that (ii) $\theta(\z)=1 $ for some $\z>0$.  Below $\o$ is sampled according to ${\rm PPP}[\rho, \nu]$.
%
%We prove Item (i). Due to \eqref{connectcont} the graph $\cG:=  \cG[\z, \b](\o)$ is contained in the graph with vertex set $\widehat \o$ and   edges given by the pairs $\{x,y\} \subset \widehat \o$ such that $x\not =y$ and $|x-y|\leq \z$. Hence $\theta(\z)=0 $  if the Boolean model  with  radius $\z/2$  on  a ${\rm PPP}[\rho]$    a.s. does not percolate. By spatial rescaling (map $x$ into $x/(\z/2)$) this takes place  if and only if the Boolean model with radius $1$ on a ${\rm PPP}[\l]$, where $\l= \rho (\z/2)^d $,  a.s. does not percolate.  It is known \cite{MR} that the last property is verified whenever  $\l < \l_c$,  $\l_c$ being  a  positive  constant depending only on the dimension $d\geq 2$. We therefore conclude that $\theta (\z)=0$ for $\z < 2(\l_c / \rho )^{1/d}$. 
%
%We prove Item (ii). We fix $\ell $ large enough to have  $p:=\nu (  [- \ell/\b, \ell/\b] )>0$. 
%The graph $\cG$ has  subgraph $G= ( \widehat\o_{\ell/\b},  \cE)$, where $ \widehat\o_{\ell/\b}:=\{x\in \widehat\o\,:\, |E_x| \leq \ell/\b\}$ and  $\cE$ is given by the pairs $\{x,y\} \subset \widehat \o_{\ell/\b}$ satisfying \eqref{connectcont}, i.e.~$|x-y|+\beta(|E_x|+|E_y|+|E_x-E_y)|)\leq \z$.  We now take $\z$ large enough to have $\z/2 \geq 4  \ell$. Then, given $x\not =y$ in  $\widehat\o_{\ell/\b}$  with $|x-y|\leq \z/2$, 
%by  
%\eqref{connectcont} we get that $\{x,y\}\in \cE$.   Hence, $G$ contains as a subgraph a Boolean model on $ \widehat\o_{\ell/\b}$ with radius $\z/4$. Since $\o$ is sampled according to ${\rm PPP}[\rho,\nu]$,  $\widehat\o_{\ell/\b}$ has law ${\rm PPP}[p\rho]$.  At this point, to prove that $\cG$ percolates a.s. for $\z$ large, it is enough to show that the Boolean model on a PPP with density $p \rho$ and radius $\z/4$ percolates for $\z$ large. By spatial rescaling (map $x$ into $x/(\z/4)$) this occurs if and only if the Boolean model with radius $1$ on a PPP with density $\l=  p \rho (\z/4)^d $ percolates a.s. This holds if $\l>\l_c$ \cite{MR}, i.e. if $\z > 4
%(\l_c / p \rho )^{1/d}$.
%\end{proof}

%%%%%%%%%%%%%%%%%%%%%%%%%
%
%  	VECCHIA PARTE SU Z^D
%
%%%%%%%%%%%%%%%%%%%%%%%%%%
%\section{Proof of Theorem \ref{teo3}}\label{proof_teo3}
%
% \begin{Definition}
% Given $\theta \in \bbR$ we define $\bbG(\theta)=\bigl( \bbZ^d, \bbE(\theta) \bigr)$ as the non--oriented graph with vertex set $\bbZ^d$ and  edges given by the bonds $\{x,y\}\in \bbE^d$ such that $E_{x,y} \leq  \theta $.
% \end{Definition}
% 
% Note that $\bbG(\theta)$ is simply the subgraph of $\bbZ^d$ supported by edges $\{x,y\}$ with conductivity $c_{x,y}(\o,\b)$ lower bounded by $e^{-\b \theta}$.
%
%
%\medskip
%
%We write  $N_\ell(\theta)$ for the number of edge--disjoint crossings of $\bbB_\ell$ along the first direction, in the graph $\bbG(\theta)$.
%  We write $C(\theta)$ for the connected component of $\bbG(\theta)$ containing the origin.
%The following lemma is the key tool to prove a scaling behavior for $D(\b)$:
%%%%%%%%%%%%%%%%%%%%%%%%%%%%%%%%%%%%%%%
%\begin{Proposition}\label{nocciola}   
%The following holds:
%\begin{itemize}
%\item[(i)]
%%$\bullet$ 
%Suppose  that $\bbE[  c_{0,e_i}(\o,\b)^2 ]<+\infty$  for all $i=1,2,\dots, d$, for $\b$ large enough.  In addition, 
%suppose that there exists $\theta^*$ such that, for $\theta >\theta^*$, for some constant  $c(\o), \ell(\o)>0$ and for all  $\o$ varying in a set  of positive $\bbP$--probability\footnote{The constants $c(\o), \ell(\o)$ and the set for which \eqref{legna} holds can depend from $\theta$}
% we 
%have \be\label{legna}
%N_\ell(\theta)[\o]  \geq c(\o)   \ell^{d-1} \text{ for } \ell \geq \ell(\o)\,.
%\en 
%Then  we have 
%\be\label{LB_rc}
%\liminf_{\b \to +\infty} \frac{1}{\b} \ln D_{1,1}(\b) \geq  - \theta ^*\,.
%\en
%
%\item[(ii)]
%%$\bullet$  
%Suppose that  there exists $c_0\in \bbR$ such that $\bbP\left( E_{0,e_i} \geq c_0\right) =1$ for all $i=1,2,\dots, d$. In addition, 
%suppose that there exists $\theta_*$ such that, for $\theta <\theta_*$, we can find  $c_1>0$ satisfying
%\be\label{cappuccino}\text{ $\,\;\;$ }\text{ $\,\;$ }
%\bbP\left( \exists x \in C(\theta) \text{ with }|x_1|\geq n  \right) \leq e^{-c_1\, n }\;\; \;\text{for $n\in \bbN$ large enough}\,.
%\en
%Then  we have 
%\be\label{UB_rc}
%\limsup_{\b \to + \infty} \frac{1}{\b} \ln D_{1,1}(\b) \leq - \theta _*\,.
%\en
%\end{itemize}
%\end{Proposition}
%We observe that, since $c_{0,e_i}(\o,\b)=e^{ - \b E_{0,e_i} (\o)}$, 
% the condition $\bbP\left( E_{0,e_i} \geq c_0\right) =1$   is equivalent to the condition that $\bbE[  c_{0,e_i}(\o,\b)^2 ] \leq c  e^{2 c_0 \b}$ for some $c>0$ and  for all $\b >0$.  %%%%%%%%%%%%%%%%%%%%%%%%%%%%%%%%%%%%
%\begin{proof}[Proof of Proposition \ref{nocciola}]  We first prove the lower bound \eqref{LB_rc}. To this aim we fix  $\theta >\theta^*$ and we take $\o $ satisfying \eqref{legna}.
%% Since each edge of $\bbG(\theta)$ has conductivity lower bounded by $e^{-\beta \theta}$, 
%%by Thompson's principle and 
%By \eqref{legna} and  Lemma \ref{lemma_LB_chiave} we then have 
%\be 
%\s_\ell ( \bbG(\theta)[\o], 1)\geq  \frac{ \cN_\ell(\o)^2}{3( \ell+1)^d}\geq \frac{ c(\o) \ell^{2d-2}}{3( \ell+1)^d}  \qquad \forall \ell \geq \ell(\o)\,.
%\en
%Since in the original resistor network  each edge of $\bbG(\theta)$ has conductivity   lower bounded by $e^{-\beta \theta}$, by Thompson's principle we get 
%\be 
%\s_\ell (\o, \b) \geq e^{-\beta \theta} \frac{ c(\o) \ell^{2d-2}}{3( \ell+1)^d} \qquad \forall \ell \geq \ell(\o)\,.
%\en
%Hence $\ell^{2-d} \s_\ell (\o, \b) \geq e^{-\beta \theta}   c(\o)c(d) $ for $\ell\geq \ell(\o)$. By \eqref{paloma} and since \eqref{legna}  holds on a set of positive $\bbP$--probability,  we then have that $D(\beta) \geq  e^{-\beta \theta}   c(\o)c(d) $, which trivially implies that $\liminf_{\b \to +\infty} \b^{-1} \ln D_{1,1}(\b) \geq  - \theta$.  By the arbitrariness of $\theta>\theta^*$ we get  \eqref{LB_rc}.
%
%We now prove the upper bound \eqref{UB_rc}. Let us fix $\theta <\theta_*$. 
%Given a subset $A\subset \bbZ^d$ we define $\text{diam}_1 (A):=\sup\{|x_1-y_1|\,:\, x, y \in A$.
%Given $n \in \bbN$  we consider the test function  $f_n \in L^\infty(\O)$ given by  
%\begin{equation}
%\label{ftest25}
%f_n:=
%\begin{cases}
%-\min \left\{ x_1:  x\in C(\theta) \right\}  &\text{if } {\text{diam}}_1(C(\theta))\leq n \b,
%\\0 &\text{otherwise}\,.
%\end{cases}\end{equation}
%By \eqref{poh} we have
%$2 D_{1,1}(\b) 
% \leq A_1+A_2$, where
%\begin{align}
%& A_1:= \bbE  \Big[ \sum_{x:|x|=1 }  c_{0,x}(\o,\b) \left
% (x_1 - \nabla_x f_n  (\o)\right)^2  \mathds{1} _{H_x^c}(\o)\Big]\,, \label{telefono125}\\
%& A_2:= \bbE  \Big[ \sum_{x:|x|=1}  c_{0,x}(\o,\b) \left
% (x_1 - \nabla_x f_n  (\o)\right)^2  \mathds{1} _{H_x}(\o) \Big]\,,\label{telefono225}
%\end{align}
%and $H_x $ denotes the event that $\{0,x\}$ is an edge of $\bbG(\theta)$.
%
%
%Let us bound $A_1$. If $H_x$ does not occur, then $c_{0,x}(\o, \b) < e^{-\b \theta }$. Using also that  $\|f_n\|_\infty \leq n\b$ we can bound $A_1 \leq e^{-\b \theta }(1+2n \b)^2$.
%We now consider $A_2$. If $H_x$ occurs then $0$ and $x$ both belong to $C(\theta)$. Hence $C(\theta)[\t_x \o]= C(\theta)-x$. In particular we have $\nabla_x f_n(\o)= x_1$ if 
% $\text{diam}_1(C(\theta))\leq n\b$ and $\nabla_x f_n(\o)= 0$ otherwise.  Due to the above observation and by Schwartz \club inequality we can bound
% \be
% \begin{split}
% A_2 &\leq \sum_{x:|x|=1} \bbE\left[ c_{0,x}(\o,\b)  \mathds{1}( \text{diam}_1(C(\theta))> n\b )\right ]\\
% & \leq \sum_{x:|x|=1} \bbE\left[ c_{0,x}(\o,\b) ^2\right]^{1/2}\bbP\left(
% \text{diam}_1(C(\theta))> n\b
% \right) ^{1/2} 
% \leq 2d e^{- c_0 \b} e^{- c_1  n\b/2}\,.
%\end{split}
% \en
% Note that in the last inequality we have used our two assumptions.
%  We now fix $n$ such that $c_1n/2 +c_0>\theta_*$. Then, as function of $\b$, we have $A_2 \ll A_1$.  Recalling that $2 D_{1,1}(\b) 
% \leq A_1+A_2$ and that $A_1 \leq e^{-\b \theta }(1+2n \b)^2$, we conclude that 
% $\limsup _{\b \to +\infty} \b^{-1} \ln D_{1,1}(\b) \leq - \theta$. By the arbitrariness of  $\theta <\theta_*$ we get \eqref{UB_rc}.
%\end{proof}
%
%
%
%\subsection{Proof of Theorem \ref{teo3}}
%%%%%%%%%%%%%%%%%%%%%%%%%%%%%%%%%%%%%%%%%%%%
% We observe that the random set of edges of the graph  $\bbG(\theta)$ equals the random set of edges in  the bond percolation on $\bbZ^d$ with parameter $p(\theta)$. Moreover, the function $p(\cdot)$ is weakly increasing.  In particular, if $p(\theta)> p_c$  (which holds if  $\theta >\theta^*$), then \eqref{legna} is satisfied by 
%\cite[Theorem~(7.68), p.~180]{G}  (for $d\geq 3$), \cite[Lemma~(11.22), p.~295]{G} (for $d=2$) and Borel--Cantelli lemma.  By applying Proposition \ref{nocciola}--(i), we then get \eqref{gelato1}. To get \eqref{gelato2} we observe that,  if $p(\theta)<p_c$ (which holds if  $\theta <\theta_*$), then \eqref{cappuccino} holds by the exponential tail decay of the radius in the subcritical Bernoulli bond percolation (cf. \cite[Theorem~(5.4), p.~88]{G}).
%
%Finally, the limit \eqref{gelato3} is an immediate consequence of the first part.
%
%
%

%%%%%%%%%%%%%%%%%%%%%%%%%%%%%%%%%%%%%%%%%%%%%%%%%%%%%%%%%%%%%%%%%%%%%%%%%%%%%%%%%%%%%%%%%%%%%%%%%%%%%%%%%%%%%%%%%%%%%%%%%%%%%%%%%%%%%%%%%%
 %%%%%%%%%%%%%%%%%%%%%%%%%%%%%%%%%%%%%%%%%%%%%%%%%%%%%%%%%%%%%%%%%%%%%%%%%%%%%%%%%%%%%%%%%%%%%%%%%%%%%%%%%%%%%%%%%%%%%%%%%%%%%%%%%%%%%%%%%%%%%%%%%%%%%%%%%%%%%%%%%%%%%%%%%%%%%%%%%%%%%%%%%%%%%%%%%%%%%%%%%%%%%%%%%%%%%%%%%%%%%%%%%%%%%%%%%%%%%%%%%%%%%%%%%%%%%%%%%%%%%%%%%%%%%%%%%%%%%%%%%%%%%%%%%%%%%%%%%%%%%%%%%%%%%%%%%%%%%%%%%%%%%%%%%%%%%%%%%%%%%%%%%%%%%%%%%%%%%%%%%%%%%%%%
%%%%%%%%%%%%%%%%%%%%%%%%%%%%%%%%%%%%%%%%%%%%
%\section{Law of $\o_\g$ with $\o$ sampled according to $\bbP$}\label{barretta}
%%%%%%%%%%%%%%%%%%%%%%%%%%%%%%%%%%%%%%%%%%%%%
%In this appendix we collect some properties of the law of $\o_\g$  with $\o$ sampled by $\bbP$, in addition to the ones presented in Lemma \ref{lemma_2021} for the case $\bbP=\hat\bbP_\nu$ \rosso{(recall from \eqref{flautino2} that $\o_\g=\{(x,E_x)\,:\, x\in \widehat{\o}\,,\; |E_x| \leq \g\}$)}. 
%%%%%%%%%%%%%%%%%%%%%%%%%%%%%%%%%%%%%%%%%%%%%
%\begin{Lemma}\label{quinoa} Let $\bbP$ 
%be the law of a stationary marked SPP with   intensity $\rho\in (0,+\infty)$ and let 
% $\bbP_0$ be the Palm distribution associated to    $\bbP$.  Fix $\g>0$ and 
%    denote by $\bbQ$ the law of $\o_\g$ when $\o$ is sampled according to $\bbP$. 
%Assume that   $\bbP( \o_\g \not =\emptyset)>0$.
%Then $\bbQ $ is a stationary marked SPP with  intensity $\rho_\g\in (0,+\infty)$. Moreover, 
%$\bbP_0(0\in \rosso{\widehat\o}_\g)=\rho_\g/\rho $ and 
% the Palm distribution $\vqz$ associated  to  $\bbQ $
%equals 
%  the law of $\o_\g$,  with $\o$ sampled according to $\vpz(\cdot | 0\in \rosso{\widehat\o}_\g)$.
%\end{Lemma}
%%%%%%%%%%%%%%%%%%%%%%%%%%%%%%%%%%%%%%%%%%%
%\begin{proof} 
%The stationary of $\bbQ $ is trivial. The intensity $\rho_\g$  of $\bbQ $  is finite since  upper bounded by the intensity  $\rho$ of $\bbP$. If  $\rho_\g $ was zero, then by stationarity we would have $\int_\O d\bbP(\o)  \widehat{\o}_\g (A)=0 $ for any $A\in \cB(\bbR^d)$. This would imply   that $\bbP( \o_\g =\emptyset)=1$, thus contradicting our assumption.
%
%It remains to  prove the last statement in the lemma concerning Palm distributions.
%To this aim let $h : \O_0 \to \bbR$ be a nonnegative measurable function. 
%We  use Campbell's identity \eqref{campanello}  with  $\bbP$ replaced by $\bbQ$ and with test function $f(x,\o):= \mathds{1}_{[0,1]^d}(x) h(\o)$. We then get 
%\be\label{sonnone1}
%\begin{split}
%\int _{\O_0} d \bbQ_0 (\o) h(\o) & = \frac{1}{\rho_\g} \int _\O d \bbQ (\o) \sum_{x\in \widehat{\o}\cap [0,1]^d} h(\t_x \o)\\
%& = \frac{1}{\rho_\g} \int _\O d \bbP(\o) \sum_{x\in \widehat{\o}_\g \cap [0,1]^d}  h\bigl(\t_x (\o_\g)\bigr) \\
%&= \frac{1}{\rho_\g} \int _\O d \bbP(\o) \sum_{x\in \widehat{\o}\cap [0,1]^d}h\bigl(\t_x (\o_\g) \bigr)\mathds{1}(|E_x|\leq \g) \,.
%\end{split}
%\en
%
%
%We now apply Campbell's identity   as  in  \eqref{campanello}  with test function $f(x,\o)=  \mathds{1}_{[0,1]^d}(x) h(\o_\g) \mathds{1} ( 0\in \widehat{\o}_\g)$. Using that $(\t_x \o)_\g=\t_x (\o_\g)$, we get 
%\begin{equation}\label{sonnone2}
%%\begin{split}
%   \int _{\O_0} d\bbP_0 ( \o) h(\o_\g) \mathds{1}(0\in \widehat{\o}_\g )  =\frac{1}{\rho} \int _{\O}d \bbP(\o)  \sum_{x\in \widehat{\o}\cap [0,1]^d}   h\bigl(\t_x(\o_\g)\bigr) \mathds{1}(|E_x| \leq \g) \,.
%   %\\ & =  \frac{1}{m} \int _{\O}d \bbP(\o)\int_{[0,1]^d}  d\hat{\o} (x) h\bigl(\t_x(\o_\g)\bigr) \mathds{1}(|E_x| \leq \g)  \,. \end{split}
% \end{equation}
% In the special case $h\equiv 1$  in \eqref{sonnone2}, we get
% \begin{equation}\label{sonnone3}
%   \bbP_0  (0\in \widehat{\o}_\g ) 
%   =  \frac{1}{\rho} \int _{\O}d \bbP(\o)  \sum_{x\in \widehat{\o}\cap [0,1]^d}   \mathds{1}(|E_x| \leq \g)=
%   \frac{\bbE\bigl[ \widehat{\o}_\g([0,1]^d)\big]}{\rho}
% = \frac{\rho_\g}{\rho}\,. 
% \end{equation}
% By combining \eqref{sonnone1}, \eqref{sonnone2} and \eqref{sonnone3} we get 
% \be
% \int _{\O_0} d \bbQ_0 (\o) h(\o)=\frac{1}{ \bbP_0  (0\in \widehat \o_\g ) }  \int _{\O_0} d\bbP_0 ( \o) h(\o_\g) \mathds{1}(0\in \widehat \o_\g ) \,,
% \en
% thus proving our final claim in Lemma \ref{quinoa}.
%\end{proof}
%%%%%%%%%%%%%%%%%%%%%%%%%%%%%%%%%%%%%%%%%%%%
%%%%%%%%%%%%%%%%%%%%%%%%%%%%%%%%%%%%%%%%%%%%%%%%%%%%%%%%%%%%%%%%%%%%%%%%%%%%%%%%%%%%%%%%%%%%%%%%%%%%%%%%%%%%%%%%%%%%%%%%%%%%%%%%%%%%%%%%%%%%%%%%%%%%%%%%%%%%%%%%%%%%%%%%%%%%%%%%%%%%%%%%%%%%%%%%%%%%%%%%%%%%%%%%%%%%%%%%%%%%
\section{Critical values $\z_c (\b,\r,\nu)$, $\l_c(\a)$ and  $\l^+_c(\a)$}\label{guizzo}
%
%In this appendix we show  that Definition  \ref{fiordo} is well posed and that  $\l^+_c(\a)$ and  $\l_c(\a)$ coincide with the critical percolation intensities  when considering the cluster at origin and the Palm distribution.

Given $r>0$ and  $\hat \o \in \hat \O$, the Boolean model $G_r(\hat \o)$ with determinist radius $r$ is the graph with vertex set $\hat \o$ and edges $\{x,y\}$ with $0<|x-y|\leq 2 r$. When $\hat \o$ is sampled according to ${\rm PPP}(\l)$, this random graph is called Poisson Boolean model. In this case 
it is known that, for a suitable $\l_c\in (0,+\infty)$, $G_r(\hat \o)$  a.s. does not percolate for $\l<\l_c$ and $G_r(\hat \o)$  a.s.   percolates for $\l>\l_c$  (cf.~\cite{MR}).

 \rosso{As in \eqref{flautino1}, given $\o\in \O$ and $\g>0$, we set $\widehat{\o}_\g:=\{x\in \widehat{\o}\,:\, |E_x|\leq \g\}$.}

\subsection{Critical threshold $\z_c (\b,\r,\nu)$}\label{marino} We now derive the existence of $\z_c (\b,\r,\nu)$ in $ (0,+\infty)$ satisfying \eqref{critico}.  Sample $\o$ according to ${\rm PPP}[\r,\nu]$.
 Since  $\cG [\z,\b](\o)$ has vertex set $\widehat \o$ and edges $\{x,y\}$ with $x\not=y$ and $|x-y|+\b( |E_x|+|E_y|+|E_x-E_y|) \leq \z$,  $\cG [\z,\b](\o)$
 is contained in $G_{\z/2}(\widehat{\o})$ and contains  $G_{\z/10}(\widehat{\o}_{\z/5\b})$. 
 By varying $\o$ the above random graphs are Poisson Boolean models with deterministic radii $\z/2$, $\z/10$ and intensities $\rho$, $\rho \nu(\z/5\b)$, respectively.
  Moreover, they are mapped by the homothety  $\bbR^d\ni x\mapsto x/\z$ into Poisson Boolean models with deterministic radii $1/2$, $1/10$ and intensities $\rho \z^d $, $\rho \nu(\z/5 \b)\z^d$, respectively. Hence, by the above  discussed phase transition of the Poisson Boolean model, a.s. $G_{\z/2}(\widehat{\o})$ does not percolate for $\z$ small and  a.s.  $G_{\z/10}(\widehat{\o}_{\z/5\b})$ percolates for $\z$ large. 
 This proves that $\theta(\z,\b, \rho, \nu)$ 
 equals $0$  for $\z$ small and $1$ for $\z$ large.  
 Since  $\theta(\cdot,\b, \rho, \nu)$ is trivially non-decreasing and, as already mentioned, has value in $\{0,1\}$, we conclude that there exists  $\z_c (\b,\r,\nu)$ in $ (0,+\infty)$ satisfying \eqref{critico}. 
 



\subsection{Critical intensities $\l_c(\a)$ and  $\l^+_c(\a)$}\label{sfondo} We now prove \orc{the} existence of values $\l_c(\a) $ and $\l_c^+(\a)$ as in Definition  \ref{fiordo}.
Let $\o$ be sampled according to $\text{PPP}[\l, \nu]$, where $\nu=\nu_{1,\a}^+$ or $\nu=\nu_{1,\a}$.
%The event $\cC:=\{ \cG[1,1](\o)\text{ percolates}\}$ is translation invariant, therefore $\text{PPP}[\l, \nu](\cC)\in \{0,1\}$ by the ergodicity of the  PPP \cite{MR}. Since by stochastic domination $\text{PPP}[\cdot, \nu](\cC)$ is  non-decreasing,  to conclude the proof  it is enough to show that $\text{PPP}[\cdot, \nu](\cC)$ takes the values $0,1$.
%To this aim we observe that,
As discussed in Section \ref{marino},  $\cG[1,1](\o)$ is contained in $G_{1/2}(\widehat{\o})$ and contains  $G_{1/10}(\widehat{\o}_{1/5})$, which are 
  Poisson Boolean models with deterministic radii $1/2$ and $1/10$  and intensities $\l$ and  $\l \nu(1/5)$ respectively. Since $\nu(1/5)>0$, the conclusion follows from the   phase transition of the Poisson Boolean model and the property that  \orc{$ \theta(1,1, \cdot, \nu)$} is non-decreasing.
  
    
  
   
   
%   (which is the so-called Poisson Boolean model with intensity $\l$ and deterministic radius $1/2$).  It is know that, when $\l$ is small enough, a.s.  $G(\o)$ does not percolate\cite{MR}. Hence  $\text{PPP}[\l, \nu](\cC)=0 $ for $\l$ small.
%
%On the other hand, $\cG[1,1](\o)$ contains the graph  $G'(\o)$  built on $\widehat{\o}_{1/5}$ by putting an edge between distinct vertexes $x,y$ if $|x-y|\leq 1/5$. Since $\widehat{\o}_{1/5}$  is a ${\rm PPP}[ \nu(1/5)\l ]$ and $\nu(1/5)>0$,  the random graph $\cG'[\o]$  is a Poisson Boolean model with intensity $\l \nu(1/5)$ and deterministic radius $1/10$. Since the latter percolates a.s. for large intensities \cite{MR}, we get that 
% $\text{PPP}[\l, \nu](\cC)=1$ for $\l$ large.


\subsection{Comparison of $\l_c(\a)$ and  $\l^+_c(\a)$ with the critical intensity  in \cite{FagMim1}} \label{liberato}
In order to apply the results of \cite{FagMim1} we show that 
\be\label{rinoceronte}
 {\rm PPP}^0(\l, \nu_{1,\a})\left( 0\leftrightarrow \infty\text{ in } \cG[1,1]\right) 
\begin{cases}
>0 &\text{ if } \l > \l_c(\a)\,,\\
=0 &\text{ if } \l < \l_c(\a)\,.
\end{cases} \en
Above ${\rm PPP}^0(\l, \nu_{1,\a})$ denotes the Palm distribution of  ${\rm PPP}(\l, \nu_{1,\a})$.
Moreover, a similar expression (with similar derivation) holds when replacing  $\nu_{1,\a}$ and $\l_c(\a)$ by $\nu^+_{1,\a}$ and $ \l^+_c(\a)$. As a consequence,  $\l_c(\a)$, $\l_c^+(\a)$ coincide with the critical intensity in \cite[Eq.~(1.5)]{FagMim1}. 

To shorten the notation let $\bbP:=  {\rm PPP}(\l, \nu_{1,\a})$, $\cA:=\{0 \leftrightarrow \infty\}$  and  $\cC:=\{ \cG[1,1](\o)\text{ percolates}\}$.  When 
$\l < \l_c(\a)$ we have  $\bbP(\cC)=0$. Since $\cC$ is translation invariant, by \cite[Lemma~7.1]{Fhom1}  we conclude that $\vpz(\cC)=0$, thus implying that $\vpz(\cA)=0$. If $\l > \l_c(\a)$, then $\bbP(\cC)=1$, thus implying that $\bbP(\cB)>0$ where $\cB:=\{\exists x\in \widehat{\o}\cap[0,1]^d\,:\, x\leftrightarrow \infty\}$ (since $\cC$ is a countable union of translations of $\cB$).
By Campbell's identity \eqref{campanello} with $f(x,\o):=\orc{\mathds{1}_{[0,1]^d}(x) \mathds{1}_\cA (\o )}$ and since $\bbP(\cB)>0$, we get
$
\vpz(\cA)=\frac{1}{\l} \int _\O d\bbP(\o) \sum_{x\in [0,1]^d} \mathds{1}(x \leftrightarrow \infty)\geq \frac{1}{\l}\bbP(\cB)>0$.


 %We focus on $\l_c(\a)$, the same arguments apply to $\l_c^+(\a)$. 


%%%%%%%%%%%%
\section{Proof of Fact \ref{cortocircuito}}\label{pinco}
%The definition of LR crossing in \cite[Section~2]{FagMim2} does not require that it has at least $3$ points, differently from our Definition \ref{def_LR}. On the other hand, since $|x-y|\leq 1$ for any edge $\{x,y\}$ of $\cG[1,1](\o)$, for $\ell$ large enough the LR crossings of $\L_\ell$ must have at least 3 points.
To get Fact \ref{cortocircuito}
it is enough apply  \cite[Thm.~1]{FagMim2} with structural function $h(a,b)=1-(|a|+|b|+|a-b|)$, probability  $P$ given by ${\rm PPP}[\l,\nu^+_{1,\a}]$ and distribution $\nu =\nu^+_{1,\a}$. To this aim we need to   check the  validity of the Assumptions (A1),...,(A5) in \cite[Section~2]{FagMim2}.  Since $\nu$ has support $[0,1]$ and for $a,b\in [0,1]$ it holds $h(a,b)=1-2 \max\{a,b\}$, the function $h$ is decreasing in $a,b$ varying in the support of $\nu$. This proves (A5). As discussed before Theorem 1 in \cite{FagMim2}, since $h$ is continuous, $\nu$ has bounded support and (A5)  is satisfied, (A4) is automatically satisfied.
(A1) and (A3) are trivially satisfied. The core is to check (A2), and our proof below is inspired by \cite[Section~8]{FagMim2}.  In our language, we need to show that there exist $\g\in (0,1)$ and \orc{$\l_*<\l $} such that $\cG[\g,1](\o)$ percolates for ${\rm PPP}[\l_*, \nu]$--a.a.~$\o$ \orc{(we recall that the $1$ in $\cG[\g,1](\o)$ is the inverse temperature)}.

By definition of $\l_c^+(\a)$, for any fixed $m> \l_c^+(\a)$ the graph $\cG[1, 1] (\o)$ percolates for  ${\rm PPP}[m,\nu]$--a.a. $\o$. 
Let us fix also  $\g\in (0,1)$. If $X$ is a random variable with law $\nu$, then   $\g X$ has law  $\nu_\g$ because of  the special form of $\nu_{1,\a}^+$ (see Definition \ref{san_benedetto}). Hence the marked SPP given by $\tilde \o:=\{ (\g x, \g E_x)\,:\, x\in \widehat \o\}$  with $\o$ sampled according to ${\rm PPP}[m,\nu]$ has law ${\rm PPP}[ \g^{-d} m, \nu_\g]$. On the other hand, when $\cG[1, 1] (\o)$ percolates, also $\cG[\g, 1](\tilde \o)$ percolates since  the map $(x,E_x) \mapsto (\g x, \g E_x)$ induces a vertex bijection between the two graphs
 and \orc{since, given $x\not= y$ in $\widehat{\o}$,   $\{x,y\}$ is an edge of of $\cG[1, 1] (\o)$ if and only if}  it holds 
%\be\label{tyler}
$|\g x- \g y| + |\g E_x|+|\g E_y|+|\g E_x- \g E_y| \leq \g $.
%\en
Hence, by the above observations, \orc{$\cG[\g, 1] (\o)$} percolates a.s. when \orc{$\o$}  is sampled according to ${\rm PPP}[ \g^{-d} m, \nu_\g]$. \orc{In conclusion, given $m'>0$, $\cG[\g, 1] (\o)$ with $\o$ sampled by ${\rm PPP}[m',\nu_\g]$ percolates a.s. if $\g^d m'>\l_c^+(\a)$.}


Let us fix $\l_*\in (\l_c^+(\a),\l)$.
By Lemma \ref{lemma_2021}, if $\o $ is sampled by ${\rm PPP}[\l_*, \nu]$, then $\o_\g:=\{(x,E_x)\,: \, x\in \widehat\o,\, |E_x|\leq \g\}$
 has law ${\rm PPP} [\nu(\g) \l_*, \nu_\g]$.  \orc{In particular, since $\cG[\g,1](\o)$  contains $\cG[\g,1](\o_\g)$ and by the previous conclusion,
 to prove that  $\cG[\g,1](\o)$ a.s. percolates with $\o$ sampled by ${\rm PPP}[\l_*, \nu]$ it is enough to have  $\g^d \nu(\g) \l_*>\l_c^+(\a)$. As $ \l_*>\l_c^+(\a)$ and $\nu(\g)=\g^{\a+1}$, it is enough to take $\g$ enough close to $1$. This proves the validity of  (A2).}
% We claim that there exists  $\g\in (0,1)$ and $m>\l^+_c(\a)$ such that 
%$  \g^{-d} m \geq \nu(\g) \l_* $. Indeed, since both  $\nu(\g) =\g^{\a+1}$ and $\g^{-d}$ approach $1$ when $\g \uparrow 1$, it is enough to take $m \in (\l_*, \l^+_c(\a))$ and then to take $\g <1$ but near enough to $1$.
%Let us fix $\g,m$ as in our claim. Then $\cG[\g,1](\o_\g)$ contains a   subgraph distributed as the above $\cG[\g, 1] (\o')$ with $\o'$  sampled according to ${\rm PPP}[ \g^{-d} m, \nu_\g]$. Since the latter percolates a.s.,
% $\cG[\g,1](\o_\g)$ itself percolates for ${\rm PPP}[\l_*, \nu]$--a.a.~$\o$.
%  Since the graph $\cG[\g,1](\o_\g)$ is included in the graph $\cG[\g,1](\o)$, we then conclude that  $\cG[\g,1](\o)$ percolates for ${\rm PPP}[\l_*, \nu]$--a.a.~$\o$, thus proving the validity of (A2).
%
%

%%%%%%%%%%%%%%%%%%%%%%%%%%%%%%%%%%%%%%%%%%%%%%%%%%%%%%%%%%%%%%%%%%%%%%%%%%%%%%%%%%%%%%%%%%%%%%%%%%%%%%%%%%%%%%%%%%%%%%%%%%%%%%%%%%%%%%%%%%
 \section{Failure of the FKG inequality in the MA random resistor network}\label{app_FKG}
%%%%%%%%%%%%%%%%%%%%%%%%%%%%%%%%%%%%%%%%%%%%%%%%%%%%%%%%%%%%%%%%%%%%%%%%%%%%%%%%%%%%%%%%%%%%%%%%
The  results on LR crossings  of \cite{FagMim2}  are derived under the assumption that the FKG inequality holds for the random graphs treated there. In the specific case of the MA random resistor network, 
the most general form of the FKG inequality   would be the following one. 
First we call an event $A\in \cB(\O)$  \emph{increasing} if, given $\o\in \O$, it holds  $\o\in A$  whenever $\cG[ \z, \b](\o')\subset \cG[\z,\b](\o)$  for some $\o'\in A$. The above inclusion means that the vertex set and edge  set of  $\cG[ \z, \b](\o')$ are included rispectively in the 
 vertex set and edge  set of  $\cG[ \z, \b](\o)$.


 
 
 We claim that  the  above general FKG does not always hold when  the energy marks in the MA resistor network appear with different sign. To explain the reason let us consider first  the following toy model. Consider the points  $0=(0,0)$, $e_1=(1,0)$, $e_2=(0,1)$ in $\bbZ^2$. We take   independent random variables  $E_0,E_{e_1},E_{e_2}$  \orc{uniformly distributed on $[-1,1]$}. We consider the random graph $  G$ with vertex set  $ V=\{0,e_1,e_2\}$ and edges  given by the pairs $\{x,y\}$ with $x\not =y $ in $ V$ 
such that  $|x-y|+|E_x|+|E_y|+|E_x-E_y| \leq 4$.
Let us consider the decreasing (with respect to the partial order given by inclusion of graphs)  events  $A,B$ defined as $A:=\big\{ \text{neither } \{0,e_1\} \text{ nor } \{e_1,e_2\} \text{ is  an edge} \big\}$ and $B:=\big \{ \{0,e_2\}  \text{ is not an  edge}\big\}$. We claim that both $A$ and $B$ have positive probability to take place but  $A\cap B$ is impossible, hence $A$ and $B$ cannot be positively correlated (thus violating the FKG inequality). Let us prove our claim.  We observe that   $w(a,b):=|a|+|b|+|a-b|$ equals $2 \max\{|a|,|b|\}$ if $ab\geq 0$ and it equals  $2 |a-b|$ if $ab \leq 0$. Hence if $w(a,b)>2$
 with $a,b\in [-1,1]$, then $a,b$ must have different signs.
 The event $A$ corresponds to the requirement that   $ w(E_0,E_{e_1})>3$ and  $ w(E_{e_1}, E_{e_2})>4-\sqrt{2}>2$. In particular, the event  $A$ has positive probability and  implies that $E_0, E_{e_1}$ have different signs and that  $E_{e_1}, E_{e_2}$ have different signs. 
Similarly, the event  $B$ has positive probability and implies that  $E_0$, $E_{e_2}$  have different signs. The constraint on the signs makes the event $A\cap B$ impossible.

Considering the MA random resistor network  $\cG [ \z,\b]=\cG[4,1]$ and 
by playing with the above toy example, one can easily  show that the FKG inequality fails when the energy marks \orc{uniformly distributed on} $[-1,1]$. 
%%%%%%%%%%%%%%%%%%%%%%%%%%%%%%%%%%%%%%%%%%%%%%%%%%%%%%%%%%%%%%%%%%%%%%%%%%%%%%%%%%%%%%%%%%%%%%%%%%%%%%%%%%%%%%%%%%%%%%%%%%%%%%%%%%%%%%%%%%%%%%%%%%%%%%%%%%%%%%%%%%%%%%%%%%%%%%%%%%%%%%%%%%%%
%Since both  $\nu(\g) =\g^{\a+1}$ and $\g^{-d}$ approach $1$ when $\g \uparrow 1$ it is enough to take $m \in (\l_*, \l^+_c(\a))$ and then to take $\g <1$ but near enough to $1$.

%$$ *****$$   
%  
% 
% 
% and we do it by    using  ideas from  the proof of \cite[Corollary~2.5]{FagMim2}. In our language, (A2) means that there exist $\l_*<\l$ and $\e>0$ small such that $\cG[1-\e, 1] (\o)$ (note: $\z=1-\e$, $\b=1$)  percolates a.s.  when $\o$ is sampled according to ${\rm PPP}[\l_*,\nu]$. We fix $\l_*\in (\l_c^+(\a),\l)$. Then the graph $\cG[1, 1] (\o)$ percolates a.s.  when $\o$ is sampled according to ${\rm PPP}[\l_*,\nu]$. 
%Given $\g\in (0,1)$,   $\{x,y\}$ is an edge of $\cG[1, 1] (\o)$ if and only if $x\not= y$, $x,y\in \widehat \o$ and  $|x-y| +  |E_x|+|E_y|+|E_x-E_y| \leq 1$. Hence
%
%\be\label{tyler}
%|\g x- \g y| + |E_x|+|E_y|+|E_x-E_y| \leq \g +(1-\g) ( |E_x|+|E_y|+|E_x-E_y|) \,.
%\en
%Apart an event of zero probability, we can restrict to energy marks in $[0,1]$. Then the r.h.s. of 
%\eqref{tyler} is bounded from above by $\g+2(1-\g)$. This proves that  the image of $\cG[1,1](\o)$ induced by the point transform $(x,E_x)\mapsto ( \g x, E_x)$ is contained in $\cG[\g+2(1-\g),1](\tilde \o)$ where $\tilde \o:=\{(\g x, E_x)\,:\,x\in \widehat \o\}$.
%\club



\bigskip

\smallskip \noindent {\bf Acknowledgments}:  I thank Hlafo Alfie Mimun for useful simulations concerning  the violation of the FKG inequality in the MA resistor network. 
%Most of this work has been written at  home (in Codroipo and in Rome), I thank my family  for the nice atmosphere.  
% I thank my sons Lorenzo and Pierpaolo for their support on 1st July 2021.
%The derivation of Mott's law uses also  preliminary results contained (together with other results) in  \cite{Fhom1}, \cite{Fhom2}, \cite{FagMim1} and \cite{FagMim2}.   \cite{Fhom1}, \cite{Fhom2} have a singular story and I thank again my family for the support along that story.

% which is a 
 %%%%%%%%%%%%%%%%%%%%%%%%%%%%%%%%%%%%%%%%%
%%%%%%%%%%%%%%%%%%%%%%%%%%%%%%%%%%%%%%%%%
%%%%%%%%%%%%%%%%%%%%%%%%%%%%%%%%%%%%%%%%%
%%%%%%%%%%%%%%%%%%%%%%%%%%%%%%%%%%%%%%%%%
%%%%%%%%%%%%%%%%%%%%%%%%%%%%%%%%%%%%%%%%%
%%%%%%%%%%%%%%%%%%%%%%%%%%%%%%%%%%%%%%%%%
%%%%%%%%%%%%%%%%%%%%%%%%%%%%%%%%%%%%%%%%%
\begin{thebibliography}{2}



\bibitem{AHL} V. Ambegoakar, B.I. Halperin, J.S. Langer; \emph{Hopping conductivity in disordered systems}. Phys, Rev B {\bf 4}, 2612-2620 (1971).



\bibitem{ABS}
G.~Androulakis, J.~Bellissard, C.~Sadel;
\emph{Dissipative Dynamics in Semiconductors at Low Temperature}. J. Stat. Phys. {\bf 147}, 
Issue 2,  448--486 (2012). 



\bibitem{BRSW}  J.~Bellissard , R.~Rebolledo , D.~Spehner , W.~Von~Waldenfels;
\emph{The Quantum Flow Of Electronic Transport I: The Finite Volume Case.} Unpublished. Available online.



\bibitem{Bi} P.~Billingsley; \emph{Convergence of Probability Measures}.  Second ed., in: Wiley Series in Probability and Statistics: Probability and Statistics, John Wiley \& Sons, Inc., New York, 1999.


\bibitem{Bis} M.~Biskup; \emph{Recent progress on the random conductance model}.  
Probability Surveys, Vol.~8,  294-373
(2011).


%\bibitem{Br}  H.~Brezis; \emph{Functional Analysis, Sobolev Spaces and Partial Differential Equations}.  New York, Springer Verlag, 2010. 



 \bibitem{CF1}  P.~Caputo, A.~Faggionato; \emph{Diffusivity of 1--dimensional
generalized Mott variable range hopping}. Ann. Appl. Probab. {\bf
19}, 1459--1494 (2009).	


\bibitem{CFP}
P.~Caputo, A.~Faggionato, T.~Prescott; \emph{Invariance principle for Mott variable range hopping and other walks on point processes}. Ann.~Inst.~H.~Poincar\'{e} Probab.~Statist.~{\bf 49}, 654--697 (2013).
%
%

\bibitem{Ji} J.~Cerny;  \emph{Critical path analysis for continuum percolation}.  Ann.~Inst.~H.~Poincar\'{e} Probab.~Statist.~{\bf 40}, 661--675 (2004).

\bibitem{CC} J.T.~Chayes, L.~Chayes; \emph{Bulk transport properties and exponent inequalities for random resistor and flow  networks}. Commun.~Mat.~Phys.~{\bf 105}, 133--152 (1986).

    
\bibitem{DV} D.J.~Daley, D.~Vere-Jones;  \emph{An Introduction to the Theory of Point Processes}. New York, Springer Verlag, 1988.




  \bibitem{demasi} A.\ De Masi, P.A.\ Ferrari, S.\ Goldstein, W.D.\ Wick; \emph{An invariance principle for reversible Markov processes. Applications to random motions in random environments.}  J. Stat. Phys. {\bf  55}, 787--855 (1989).



\bibitem{DS} P.G.~Doyle, J.L.~Snell; \emph{Random walks and electric networks}. Carus Mathematical Monographs, Mathematical Association of America  (1984). 




 \bibitem{Fhom1} A.~Faggionato; \emph{Stochastic homogenization of random walks on  point processes}. \rosso{Ann.~Inst.~H.~Poincar\'{e} Probab.~Statist. {\bf 59}, 662--705 (2023).}
 

 
  

%\bibitem{F_sep} A.~Faggionato; \emph{Hydrodynamic limit of simple exclusion processes in symmetric random environments via duality and homogenization}. arXiv:2011.11361 (2020).

\bibitem{Fhom2} A.~Faggionato; \emph{Scaling limit of the \rosso{directional} conductivity of random resistor networks on point processes}. arXiv:2108.11258 (2021).



   \bibitem{FM}
 A.\ Faggionato, P.\ Mathieu; \emph{Mott law as upper bound for a random
walk in a random environment}. Comm. Math. Phys. {\bf 281}, 263--286
(2008). 


\bibitem{FagMim1}
 A.~Faggionato, H.A.~Mimun; \emph{Connection probabilities in the Poisson Miller--Abrahams random resistor network and other Poisson random graphs with bounded edges}.  ALEA, Lat. Am. J. Probab. Math. Stat. {\bf 16}, 463-486 (2019).

\bibitem{FagMim2} A.~Faggionato, H.A.~Mimun;
\emph{Left-right crossings in the Miller-Abrahams random resistor network and in generalized Boolean models}.   Stoch. Process. Their Appl. {\bf 137}, 62--105 (2021).




\bibitem{FSS} A.\ Faggionato, H.\ Schulz--Baldes, D.\ Spehner; {\em
     Mott law as lower bound for a random walk in a random
     environment.} Comm.\ Math.\ Phys., {\bf 263},
     21--64 (2006).


\bibitem{Fel2} W. Feller; \emph{An introduction to probability theory and its applications},  Vol.~II,  Second edition. John Wiles \& Sons, 1970. % kuka : controllare che sia second edition



\bibitem{FKAS} P.~Franken,  D.~K\"{o}nig, U.~Arndt, V.~Schmidt, \emph{Queues
and Point Processes}.
John Wiley and Sons, Chichester (1982).



%\bibitem{G} G. Grimmett;  \emph{Percolation}. 	Die Grundlehren der mathematischen Wissenschaften {\bf 321}. Second edition, Springer Verlag, Berlin, 1999.



\bibitem{Kal}  O.~Kallenberg. \emph{Foundations of modern probability}. Second edition. Probability and its Applications,  Springer Verlag, Berlin, 2002.



\bibitem{Ke} H.~Kesten; \emph{Percolation theory for mathematicians}. Birkh\"auser, Boston, 1982.


\bibitem{Kur} J.~Kurkij\"arvi; \emph{Hopping conductivity in one dimension}. Phys. Rev. B {\bf 8}, 922--924.
(1973).

\bibitem{LP} G.~Last, M.~Penrose; \emph{Lectures on the Poisson Process}.  IMS Textbook,  Cambridge University Press, Cambridge, 2017.

\bibitem{MR}  R. Meester, R. Roy;  \emph{Continuum percolation}. Cambridge Tracts in Mathematics {\bf 119}. First edition, Cambridge University Press, Cambridge, 1996.

\bibitem{MA}  A. Miller, E. Abrahams; \emph{Impurity Conduction at Low Concentrations}. Phys. Rev. {\bf 120}, 745--755 (1960).



\bibitem{Mott}  N.F.  Mott:  J. Non-Crystal. Solids {\bf 1}, 1 (1968);  Phil. Mag {\bf 19}, 835 (1969).

\bibitem{Mott_Nob} N.F.  Mott; \emph{Electrons in glass}. 
Nobel Lecture, 8 December, 1977. Available online at https://www.nobelprize.org/uploads/2018/06/mott-lecture.pdf




\bibitem{MD} N.F. Mott, E.A. Davis; \emph{Electronic processes in non-crystalline materials}.  Oxford Classic Texts in the Physical Sciences, OUP Oxford, Oxford, 2012.


\bibitem{PR} A.~Piatnitski, E~Remy; \emph{Homogenization of elliptic difference operators}.
SIAM J.~Math. Anal. {\bf 33}, 53-83 (2001).


\bibitem{Po}  M.~Pollak; \emph{A percolation treatment of dc hopping conduction}. Journal of Non-Crystall. Solids {\bf 11}, 1, 1--24 (1972).


\bibitem{POF} M.~Pollak, M.~Ortu{\~n}o, A. Frydman; \emph{The electron glass}. First edition, Cambridge University Press, United Kingdom, 2013.



\bibitem{Sa} M.~Sahimi; \emph{Applications of percolation theory}. Taylor \& Francis,
CRC Press, 1994.



\bibitem{SE}
S.~Shklovskii, A.L.~Efros; \emph{Electronic Properties of Doped Semiconductors}. Springer Verlag, Berlin, 1984.





\bibitem{T} A.A.~Tempel'man; \emph{Ergodic theorems for general dynamical systems}. Trudy Moskov. Mat. Obsc. {\bf 26}, 95--132 (1972) [Translation in Trans. Moscow Math. Soc. {\bf 26}, 94--132, (1972)]


\end{thebibliography}



\end{document}

%%%%%%%%%%%%%%%%%%%%%%%%%%%%%%%%%
%%%%%%%%%%%%%%%%%%%%%%%%%%%%%%%%%
%%%%%%%%%%%%%%%%%%%%%%%%%%%%%%%%%
%%%%%%%%%%%%%%%%%%%%%%%%%%%%%%%%%
%%%%%%%%%%%%%%%%%%%%%%%%%%%%%%%%%
%%%%%%%%%%%%%%%%%%%%%%%%%%%%%%%%%
%%%%%%%%%%%%


\end{document}
\section{Progresses}

{\bf Warning}: $\bbP_\rho$ is the law of the marked Poisson point process where the iid marks have law $\nu$ and the PPP $\xi$ has density $\rho$. We write $\o = \left(\xi , (E_x)_{x\in \xi}\right)$ for a generic realization of the marked simple point process. 



\bigskip

For $n\in\bbN$ define the following sets:
\begin{itemize}
\item $\L_n:=[-n,n]^d$;
\item $A_n:=\{x\in\bbR^d\,|\,x_1\leq -n,\,|x_i|\leq n\,\forall i=2,\ldots,d\}$;
\item $B_n:=\{x\in\bbR^d\,|\,x_1\geq n,\,|x_i|\leq n\,\forall i=2,\ldots,d\}$;
\item $D_n=\{x\,:\,|x_i|\leq n,\,\forall i=2,\ldots,d\}=A_n\cup\L_n\cup B_n$.
\end{itemize}
Let  $u:\xi\rgh\bbR$ be the potential function:
\begin{itemize}
\item $u(x)=0$ if $x\in A_n\cap\xi$;
\item $u(x)=1$ if $x\in B_n\cap\xi$;
\item $\forall\,x\in\L_n\cap \xi$ there holds $\sum_{y\in \xi} c_{x,y}(u(y)-u(x))=0$.
\end{itemize}
We define the effective conductivity 
\begin{equation}
\label{condcompl}
\s_n(\b,\o):=\frac{1}{2}\sum_{x \in D_n \cap \xi }\sum_{y \in D_n \cap \xi  }c_{x,y}(\nabla_{x,y} u)^2,
\end{equation}
where $\nabla_{x,y} u:=u(x)-u(y)$. Since we are interested at the behaviour for of $n^{2-d}\s_n(\b,\o)$ for $n\rgh\infty$, by lemma \ref{ext} we can restrict to consider a new conductivity, that we still denote by $\s_n(\b,\o)$, defined as
$$\s_n(\b,\o)=\sum_{x\in \L_n\cap \xi}\sum_{\substack{y\in D_n\cap \xi:\\x\neq y}}c_{x,y}(\nabla_{x,y}u)^2.$$



\bigskip

Consider the proof of Theorem \ref{teo1}. We have seen that points $\{ x\in \xi\,:\, |E_x| > \z/\b\}$ are isolated in the graph $\cG_{\z, \b,\rho}$, and in particular we have
\begin{equation}
\cG_{\z,\b,\rho}= \tilde{\cG}_{\z,\b,\rho} \sqcup \{ x\in \xi\,:\, |E_x| > \z/\b\}\,.
\end{equation}
Consider the map $\phi:\bbR^d \to \bbR^d$ defined as $\phi(x)= x/\z$.  
We have then showed that $\phi\bigl( \tilde \cG_{\z,\b,\rho}\bigr)$ has the same law of $ \cG_{1,1,\l}$ where 
$\l:=(\z^{\a+1+d}/\b^{\a+1})\rho$.
Take $\z:=\z_c(1+\e)$ with $\z_c=\z_c(\b,\rho)$ defined in \eqref{zetacritico}.  Then
\begin{equation}
\l = (1+\e) ^{\a+1+d} (\z_c ^{\a+1+d}/\b^{\a+1})\rho=(1+\e) ^{\a+1+d} \l_c\,.
\end{equation}
Note that $\l$ \underline{ does not depend } from $\rho,\beta$ and moreover  $\l>\l_c$.

\medskip

Note that all edges in $\cG_{\z,\b,\rho}$, and therefore in $ \tilde{\cG}_{\z,\b,\rho}$, have conductivity at least $e^{-\z_c(1+\e)}$.

\medskip

We write  $\s_n(\b, \o)$ for the effective conductance of the MA resistor network on $\L_n$ with potential $0$ on $A_n$ and potential $1$ on $B_n$. Given a graph $G$ in $\bbR^{d}$   we write $\s_n^{*}(G)$ for the effective conductivity on $\L_n$ with potential $0$ on $A_n$ and potential $1$ on $B_n$ of the resistor network obtained from $G$ by assigning conductivity $1$ to any edge in $G$.  Then, by the Dirichlet principle, we get
\begin{equation}\label{lenticchie}
\boxed{
\s_n(\b, \o) \geq e^{-\z_c(1+\e)} \s_n^{*}\bigl(\tilde  \cG_{\z,\b,\rho} (\o)\bigr) \,.}
\end{equation}
On the other hand, by the above discussion concerning the proof of Theorem \ref{teo1}, the random sequence  (under $\bbP_\rho$)
\begin{equation}\label{gaeta1}
\left( n ^{2-d} \s_n^{*}\bigl( \tilde \cG_{\z,\b,\rho}(\o) \bigr) \right)_{n \geq 1}
\end{equation}
has the same law of the random sequence (under $\bbP_\l$)
\begin{equation}\label{gaeta2}
\left( n^{2-d}\s_{n/\z}^{*}\bigl( \cG_{1,1,\l}(\o) \bigr) \right)_{n \geq 1}
\end{equation}
We want to prove that $\bbP_\l$--a.s. it holds  
\begin{equation}\label{sio}
\boxed{\k_0:=\varliminf _{n \to \infty}   n^{2-d}\s_{n}^{*}\bigl( \cG_{1,1,\l}(\o)\bigr)>0\,.}
\end{equation}
Note that $\k_0$ depends only on $\l$, in particular it is 
$\b$--independent.
Then, by the equivalence in law of \eqref{gaeta1} and \eqref{gaeta2}, we have $\bbP_\rho$--a.s.
\begin{equation}\label{sio1}
\varliminf _{n \to \infty}  n^{2-d} \s_n^{*}\bigl( \tilde\cG_{\z,\b,\rho}(\o) \bigr) \geq \k _0\z ^{2-d}= \k_0 \z_c  ^{2-d} (1+\e) ^{2-d}\,.
\end{equation}

By combining \eqref{lenticchie} with \eqref{sio1}  and using that 
\[ \beta ^{-\frac{\a+1}{\a+1+d}}\z_c(1+\e)= (1+\e) (\l_c /\rho) ^{\frac{1}{\a+1+d}}
\,,\]
we have $\bbP_\rho$--a.s.
\begin{equation}
\begin{split}
\varliminf_{\b \to \infty} \varliminf _{n \to \infty} \beta ^{-\frac{\a+1}{\a+1+d}}\log [ n ^{2-d} \s_n (\b,\o)] \geq 
&-(1+\e) (\l_c /\rho) ^{\frac{1}{\a+1+d}}\\
&+ \varliminf_{\b \to \infty} \beta ^{-\frac{\a+1}{\a+1+d}} \log\left[  \k_0 \z_c (\b,\rho)  ^{2-d} (1+\e) ^{2-d} \right]\,,
\end{split}
\end{equation}
thus implying  $\bbP_\rho$--a.s. the asymptotic lower bound 
\begin{equation*}\varliminf_{\b \to \infty} \varliminf _{n \to \infty} \beta ^{-\frac{\a+1}{\a+1+d}}\log [ n ^{2-d}  \s_n (\b,\o)] \geq  -(1+\e) (\l_c /\rho) ^{\frac{1}{\a+1+d}}\,.
\end{equation*}
By the arbitrariness of $\e$ we have that $\bbP_\rho$--a.s. the asymptotic lower bound 
\begin{equation}\label{arrivo}
\boxed{\varliminf_{\b \to \infty} \varliminf _{n \to \infty} \beta ^{-\frac{\a+1}{\a+1+d}}\log [ n ^{2-d}  \s_n (\b,\o)] \geq  - (\l_c /\rho) ^{\frac{1}{\a+1+d}}\,.}
\end{equation}
It remains to prove \eqref{sio}. This is the content of Prop. \ref{prop_sio}


\begin{Proposition}\label{prop_sio}
Given $\l>\l_c$  the following holds $\bbP_\l$--a.s.
\begin{equation}
\varliminf _{n \to \infty}   n^{2-d}\s_{n}^{*}\bigl( \cG_{1,1,\l}(\o)\bigr)>0\,.
\end{equation} 
\end{Proposition}
This is based on the following conjecture-lemma. In order to state it, we need to define before a left-right crossing of $\L_n$ of the graph  $\cG_{1,1,\l}$.  We recall that  $\cE_{1,1,\l}(\o)$ is the set of edges,
we denote by % $\cE_\l$ the edge set of $\cG_{1,1,\l}(\o)$ and by
 $\xi_\l$ the set of  nodes. A \textit{left-right crossing} (LR crossing for short) of $\L_n$ (in $\cG_{1,1,\l}$) is a sequence of distinct points $x^{(1)},\ldots,x^{(j)}$ such that
\begin{itemize}
\item $x^{(i)}\in\xi_\l\cap D_n$ for all $i=1,\ldots,j$;
\item $\{x^{(i)},x^{(i+1)}\}\in\cE_{1,1,\l}$ for all $i=1,\ldots,j-1$;
\item $(x^{(1)})_1\leq -n$ and $(x^{(j)})_1\geq n$;
\item $-n<(x^{(k)})_1<n$ for all $k=2,\ldots,j-1$.
\end{itemize}
We denote by $M_n(\o)$ the maximum number of edge-disjoint LR crossing of $\L_n$ in $\cG_{1,1,\l}(\o)$. We recall that two LR crossings are edge--disjoint if they have no common edge, on the other hand they could have some common vertex. 
\begin{Lemma}
\label{maxLR}
For any $\l>\l_c$  there exist a positive constant $C$ and a non-negative sequence $(\a_n)_{n\in\bbN}$ such that $\sum_{n\in\bbN}\a_n<+\infty$ and
$$\bbP_\l(M_n\leq Cn^{d-1})\leq \a_n\quad\forall n\geq 1.$$
\end{Lemma}
In order to prove proposition \ref{prop_sio}, we exhibit a lower bound for $\s_n^*(\cG_{1,1,\l}(\o))$. 
\\By lemma \ref{maxLR} and Borel-Cantelli's lemma, there exists $C=C(\l)>0$ such that the following holds: there exists an event $\cH$ such that $\bbP_{\l}(\cH)=1$ and for all $\o\in\cH$ there exists $n_0(\o)\in \bbN$ such that
$$M_n(\o)\geq Cn^{d-1}\quad\forall\,n>n_0(\o).$$
We set $m:=\lceil Cn^{d-1}\rceil$. Assuming $\o\in\cH$ and $n\geq n_0(\o)$, we call $\p_1,\ldots,\p_m$ some edge-dijoint open LR crossings of $\L_n$ (referred to the graph $\cG_{1,1,\l}(\o)$).
\\Then
\begin{equation}
\label{low1}
\begin{split}
\s_n^*(\cG_{1,1,\l}(\o))&\geq \sum_{\substack{\{x,y\}\in \cE_\l:\\ x,y\in D_n\cap\xi_\l}} c_{x,y}(\nabla_{x,y} u)^2 \geq\sum_{i=1}^m\sum_{\{x,y\}\in \p_i} c_{x,y}(\nabla_{x,y} u)^2 
\\ &\geq e^{-1}\sum_{i=1}^m\sum_{\{x,y\}\in \p_i} c_{x,y}(\nabla_{x,y} u)^2.
\end{split}
\end{equation}
By Cauchy-Schwartz's inequality we obtain
$$1=\left(\sum_{\{x,y\}\in \p_i}\nabla_{x,y} u\right)^2\leq |\p_i|\sum_{\{x,y\}\in \p_i} c_{x,y}(\nabla_{x,y} u)^2.$$
Then r.h.s. in \eqref{low1} can be bounded from below in the following way
\begin{equation}
\label{low2}
\begin{split}
e^{-1}\sum_{i=1}^m\sum_{\{x,y\}\in \p_i} c_{x,y}(\nabla_{x,y} u)^2\geq  e^{-1}\sum_{i=1}^m\frac{1}{|\p_i|}.
\end{split}
\end{equation}
By Jensen's inequality we can bound from below the r.h.s. of \ref{low2}
\begin{equation}
\label{low3}
\begin{split}
 e^{-1}\sum_{i=1}^m\frac{1}{|\p_i|}&=me^{-1}\cdot\frac{1}{m}\sum_{i=1}^m\frac{1}{|\p_i|}\geq me^{-1}\left(\frac{1}{m}\sum_{i=1}^m|\p_i|\right)^{-1}=
\\ &=\frac{m^2e^{-1}}{\sum_{i=1}^m|\p_i|}.
\end{split}
\end{equation}
Observe that 
$$\sum_{i=1}^m|\p_i|\leq \sum_{x\in \xi_\l\cap \L_n}\sum_{\substack{y\in\xi_\l\cap D_n,\\0<|x-y|<1}} 1.$$ 
Then r.h.s. of \eqref{low3} can be bounded from below in the following way
\begin{equation}
\label{low4}
\begin{split}
\frac{m^2e^{-1}}{\sum_{i=1}^m|\p_i|}\geq\frac{m^2e^{-1}}{\sum_{x\in \xi_\l\cap \L_n}\sum_{\substack{y\in\xi_\l\cap D_n,\\0<|x-y|<1}} 1}.
\end{split}
\end{equation}
Since $m\geq Cn^{d-1}$, we get from \eqref{low1}, \eqref{low2}, \eqref{low3} and \eqref{low4} that
\begin{equation}
\label{low5}
\begin{split}
\s_n^*(\cG_{1,1,\l}(\o))\geq \frac{Cn^{2d-2}}{\sum_{x\in \xi_\l\cap \L_n}\sum_{\substack{y\in\xi_\l\cap D_n,\\0<|x-y|<1}} 1}=\frac{Cn^{d-2}}{\frac{1}{n^d}\sum_{x\in \xi_\l\cap \L_n}\sum_{\substack{y\in\xi_\l\cap D_n,\\0<|x-y|<1}} 1}
\end{split}
\end{equation}
Observe that
\begin{equation}
\label{low6}
\begin{split}
\frac{1}{n^d}\sum_{x\in \xi_\l\cap \L_n}\sum_{\substack{y\in\xi_\l\cap D_n,\\0<|x-y|<1}} 1=\frac{\xi_\l(\L_n)}{n^d}\frac{1}{\xi_\l(\L_n)}\sum_{x\in\xi_\l\cap\L_n} f(\tau_x\xi_\l),
\end{split}
\end{equation}
where $f(\xi_\l)=\xi_\l(B_1(0)\setminus\{0\})$ and $B_1(0):=\{x\in\bbR^d\,|\,|x|<1\}$. Moreover $f\in L^1(\bbP_{0,\l})$, since
$$\bbE_{0,\l}[f(\xi_\l)]=c,$$
where $c>0$ is a constant depending only on $\l$.
Then there holds
\begin{equation}
\label{low7}
\lim_{n\rgh\infty}\frac{\xi_\l(\L_n)}{n^d}=\tilde c\quad \bbP_{0,\l}-\text{a.s.}
\end{equation}
and
\begin{equation}
\label{low8}
\lim_{n\rgh\infty}\frac{1}{\xi_\l(\L_n)}\sum_{x\in\xi_\l\cap\L_n} f(\tau_x\xi_\l)=\bbE_{0,\l}[f(\xi_\l)]=\bbE_{0,\l}[\xi_\l(B_1(0)\setminus\{0\})]=c\quad \bbP_{0,\l}-\text{a.s.},
\end{equation}
where $\tilde c$ and $c$ are positive constants depending only on $\l$.
\\So, by \eqref{low5}, \eqref{low6}, \eqref{low7} and \eqref{low8}, we get that
$$\varliminf_{n\rgh\infty}n^{2-d}\s_n^*(\cG_{1,1,\l}(\o))\geq c,$$
where $c>0$ is a constant depending only on $\l$. This concludes the proof of proposition \ref{prop_sio}.

%%%%%%%%%%%%%%%%%%%%%%%%%%%%%%%%%%%%%%%%%
%%%%%%%%%%%%%%%%%%%%%%%%%%%%%%%%%%%%%%%%%
%%%%%%%%%%%%%%%%%%%%%%%%%%%%%%%%%%%%%%%%%
%%%%%%%%%%%%%%%%%%%%%%%%%%%%%%%%%%%%%%%%%
%%%%%%%%%%%%%%%%%%%%%%%%%%%%%%%%%%%%%%%%%

\section{Upper bound}
\begin{Proposition}\label{luce1} For any $\rho >0$  the following holds $\bbP_\rho$--a.s.
\begin{equation}\label{arrivo_upper}
\varlimsup_{\b \to \infty} \varlimsup _{n \to \infty} \beta ^{-\frac{\a+1}{\a+1+d}}\log [ n ^{2-d}  \s_n (\b,\o)] \leq  - (\l_c /\rho) ^{\frac{1}{\a+1+d}}\,.
\end{equation}
\end{Proposition}

We first consider a similar bound for the diffusion matrix of Mott random walk. To this aim let $D(\b)$ be the $d\times d$ symmetric matrix such that
\begin{equation}\label{volinf}
(a \cdot D(\b) a)= \inf _{ f\in L^\infty(\bbP_{0,\rho})} \bbE_{0,\rho} \Big[\int \xi(dx) c_{0,x}\bigl( a\cdot x- \nabla_x f(\o) \bigr)^2
\Big]\,,\qquad a\in \bbR^d\,,
\end{equation}
where $\nabla_{x} f(\o):= f( \t_x \o)- f(\o)$.  For isotropy $D(\b)$ is a multiple of $\bbI$, i.e. $D(\b)= (e_1 \cdot D(\b) e_1)\bbI$.
\begin{Proposition}\label{luce2} For any $\rho >0$  it holds
\begin{equation}\label{arrivo_upper}
\varlimsup_{\b \to \infty}  \beta ^{-\frac{\a+1}{\a+1+d}}\log  (e_1 \cdot D(\b, \rho) e_1)\leq  - (\l_c /\rho) ^{\frac{1}{\a+1+d}} \,.
\end{equation}
\end{Proposition}

\subsection{Proof of Proposition \ref{luce2}}
In order to prove Proposition \ref{luce2} we consider the graph $\cG_{\z,\b,\r}$ with $\z=\z_c(1-\e)$. By Theorem \ref{teo1} a.s. the graph $\cG_{\z,\b,\r}$ does not percolate. As we have done before we can split such a graph as
\begin{equation}
\cG_{\z,\b,\rho}= \tilde{\cG}_{\z,\b,\rho} \sqcup \{ x\in \xi\,:\, |E_x| > \z/\b\}\,.
\end{equation}
We remember that we denote by $\tilde\cE$ the edge set of the graph $\tilde \cG_{\z,\b,\r}$. Moreover, if $\{x,y\}\in \tilde\cE$ if and only if  $c_{x,y}\geq e^{-\z}=e^{-\z_c(1-\e)}$.


Fixed $m\in\bbN_+$ we define the test function $f_{\z m}:\O\rgh\bbR$ in this way
\begin{equation}
\label{ftest}
f_{\z m}(\o)=
\begin{cases}
-\min_{x\in C_{\z,\b,\r}(0,\o)} x_1 &\text{if diam}(C_{\z,\b,\r}(0,\o))\leq \z m,
\\0 &\text{otherwise},
\end{cases}\end{equation}
where for $x\in\bbR^d$ we denote by $C_{\z,\b,\r}(x,\o)$ the cluster which contains $x$ in the graph $\cG_{\z,\b,\r}(\o)$ (if $x$ is not a vertex of $\cG_{\z,\b,\r}(\o)$ we set $C_{\z,\b,\r}(x,\o):=\emptyset$) and for every $A\subseteq\bbR^d$ we define $\text{diam}(A):=\max\{|x_i-y_i|\,:\,x,y\in A,\,i=1,\ldots,d\}$.
\\In order to compute the l.h.s. of \eqref{arrivo_upper}, we look for an upper bound for the following expectation
\begin{equation}
\label{up1}
\begin{split}
&\bbE_{0,\r}\left[\int_{\bbR^d} \xi(dx)c_{0,x}(x_1-\nabla_x f_{\z m}(\o))^2\right]=
\\ &=\bbE_{0,\r}\left[\int_{\bbR^d\setminus\{0\}} \xi(dx)c_{0,x}(x_1-\nabla_x f_{\z m}(\o))^2\right]=\cC_1+\cC_2,
\end{split}
\end{equation}
where
\begin{equation}
\label{cc12}
\begin{split}
&\cC_1:=\bbE_{0,\r}\left[\int_{\bbR^d\setminus\{0\}}\xi(dx)c_{0,x}(x_1-\nabla_x f_{\z m}(\o))^2\mathds{1}_{H^c_x}\right],
\\ &\cC_2:=\bbE_{0,\r}\left[\int_{\bbR^d\setminus\{0\}}\xi(dx)c_{0,x}(x_1-\nabla_x f_{\z m}(\o))^2\mathds{1}_{H_x}\right]
\end{split}
\end{equation}
and, for $x\in\bbR^d$, $H_x$ is the event that $\{0,x\}$ is an element of $\tilde \cE$.
\medskip
\\Before starting with the computation of the previous terms we want to draw attention to the following estimates that we will find many times in the computation of many terms. 

Define $\l:=(\z^{\a+1+d}/\b^{\a+1})\r= (1-\e) ^{\a+1+d} (\z_c ^{\a+1+d}/\b^{\a+1})\rho=(1-\e) ^{\a+1+d} \l_c$ and observe that $\l$ depends only on $\l_c$ and $\e$. Moreover, since $\l<\l_c$, then the graph $\cG_{1,1,\l}$ does not percolate.

\begin{Remark}
Observe that $\r,\l_c,\e$ (and consequently $\l$) are fixed parameters, so in the following computations we will not stress the dependence of some constants on them.
\\We will denote by $c(\b)$ a positive constant that has a polynomial-dependence on $\b$, that is, there exist two positive constants $k_1$ and $k_2$ such that $c(\b)\leq k_1\b^{k_2}$.
\\Observe that $\z$ is a positive constant that has a polynomial-dependence on $\b$, since $\z=(1-\e)\z_c$ and $\z_c$ is defined in \eqref{zetacritico}.
\end{Remark}

$\bullet$ Fixed a positive constant $a$, since $f_{\z m}(\o)=0$ if $|E_0|> \z/\b$, we get
\begin{equation}
\label{int3}
\begin{split}
\bbE_{0,\r}[|f_{\z m}(\o)|^a]&=\bbE_{0,\r}\left[|f_{\z m}(\o)|^a\mathds{1}_{\{|E_0|\leq \z/\b\}}\right]=
\\ &=\bbE_{0,\r}\left[|f_{\z m}(\o)|^a\,|\,|E_0|\leq \z/\b\right]\nu\left(\left[-\frac{\z}{\b},\frac{\z}{\b}\right]\right)=
\\ &=c(\b)\bbE_{0,\l}\left[|f_{ m}(\o)|^a\right]\leq c(\b)\bbE_{0,\l}\left[| \text{diam}(C_{1,1,\l}(0,\o)|^a\right]
\\ &=c(\b)\int_0^\infty \bbP_{0,\l}(\text{diam}(C_{1,1,\l}(0,\o))>s^{\frac{1}{a}})ds
\\ &\leq  c(\b)\int_0^{+\infty}e^{-c(\l) s^{\frac{1}{a}}}ds=: \tilde c(\b),
\end{split}
\end{equation}
% kokeshi: spiegare quarta identita'
where $\tilde c(\b)$ is a positive constant that has a polynomial-dependence on $\b$, while $c(\l)$ is a positive constant depending only on $\l$. Observe that the last inequality is due to Theorem 1 in \cite{FAMH}, in fact
\begin{equation}
\label{stimadiam}
\bbP_{0,\l}(\text{diam}(C_{1,1,\l}(0,\o))>n)\leq \bbP_\l(0\lrgh S_n)\leq e^{-c(\l)n},
\end{equation}
where $S_n:=\{x\in\bbR^d\,:\,||x||_\infty=n\}$.
\bigskip
\\\indent $\bullet$ Fixed two positive constants $a$ and $b$, we compute the following integral applying \eqref{int3} and \eqref{f1} with $g(x,\xi)=e^{-a|x|}|f_{\z m}(\tau_x\o)|^b$ 
\begin{equation}
\label{int4}
\begin{split}
&\bbE_{0,\r}\left[\int_{\bbR^d\setminus\{0\}}\xi(dx)e^{-a|x|}| f_{\z m}(\tau_x\o))|^b\right]=\r\int_{\bbR^d}dxe^{-a|x|}\bbE_{0,\r}[|f_{\z m}|^b]\leq c(\b).
\end{split}
\end{equation}
\bigskip
\\\indent $\bullet$ Fixed two positive constants $a$ and $b$, we compute the following integral applying Cauchy-Schwartz's inequality, \eqref{int3} and \eqref{int2} with $h(x)=e^{-a|x|}$
\begin{equation}
\label{int4.5}
\begin{split}
&\bbE_{0,\r}\left[\int_{\bbR^d\setminus\{0\}}\xi(dx)e^{-a|x|}| f_{\z m}(\o))|^b\right]\leq\left(\bbE_{0,\r}\left[\left(\int_{\bbR^d\setminus\{0\}}\xi(dx)\,e^{-a|x|}\right)^2\right]\right)^\frac{1}{2}\left(\bbE_{0,\r}\left[| f_{\z m}(\o))|^{2b}\right]\right)^\frac{1}{2}\\ &\leq c(\b).
\end{split}
\end{equation}
\bigskip
\\ \indent $\bullet$Fixed two positive constants $a$ and $b$, we compute the following integral applying Cauchy-Schwartz's inequality and \eqref{int3}
\begin{equation}
\label{int5}
\begin{split}
\bbE_{0,\r}\left[\int_{\bbR^d\setminus\{0\}}dx\,e^{-a|x|}| f_{\z m}(\o))|^b\right]
&\leq\left(\bbE_{0,\r}\left[\left(\int_{\bbR^d\setminus\{0\}}dx\,e^{-a|x|}\right)^2\right]\right)^\frac{1}{2}\left(\bbE_{0,\r}\left[| f_{\z m}(\o))|^{2b}\right]\right)^\frac{1}{2}
\\&\leq c(\b).
\end{split}
\end{equation}
\bigskip
\\Now we can start with the computation of $\cC_1$
\begin{equation}
\label{up2}
\begin{split}
&\cC_1=\bbE_{0,\r}\left[\int_{\bbR^d\setminus\{0\}}\xi(dx)c_{0,x}(x_1-\nabla_x f_{\z m}(\o))^2\mathds{1}_{H^c_x}\right]=
\\ & =\bbE_{0,\r}\left[\int_{\bbR^d\setminus\{0\}}\xi(dx)c_{0,x}(x_1-\nabla_x f_{\z m}(\o))^2\mathds{1}_{H^c_x}\mathds{1}_{\{|E_x|> \z/\b\}}\right]+
\\ &+\bbE_{0,\r}\left[\int_{\bbR^d\setminus\{0\}}\xi(dx)c_{0,x}(x_1-\nabla_x f_{\z m}(\o))^2\mathds{1}_{H^c_x}\mathds{1}_{\{|E_x|\leq \z/\b\}}\right]
\end{split}
\end{equation}
Consider the first term in the sum in the r.h.s. of \eqref{up2}
\begin{equation}
\label{up3}
\begin{split}
\bbE_{0,\r}\left[\int_{\bbR^d\setminus\{0\}}\xi(dx)c_{0,x}(x_1-\nabla_x f_{\z m}(\o))^2\mathds{1}_{H^c_x}\mathds{1}_{\{|E_x|> \z/\b\}}\right]
\end{split}
\end{equation}
Observe that if $|E_x|>\z/\b$, then $f(\tau_x\o)=0$ and 
$$(x_1-\nabla_xf_{\z m}(\o))^2\leq 2|x_1|^2+2|f_{\z m}(\o)|^2,$$
and 
$$c_{0,x}\leq e^{-|x|-\z}.$$
Then
\begin{equation}
\label{up4}
\begin{split}
&\bbE_{0,\r}\left[\int_{\bbR^d\setminus\{0\}}\xi(dx)c_{0,x}(x_1-\nabla_x f_{\z m}(\o))^2\mathds{1}_{H^c_x}\mathds{1}_{\{|E_x|> \z/\b\}}\right]
\\ &\leq2e^{-\z}\left(\bbE_{0,\r}\left[\int_{\bbR^d\setminus\{0\}} \xi(dx)e^{-|x|}|x_1|^2\right]+\bbE_{0,\r}\left[\int_{\bbR^d\setminus\{0\}} \xi(dx)e^{-|x|}|f_{\z m}(\o)|^2\right]\right)
\end{split}
\end{equation}
So, by \eqref{int4.5} and \eqref{int1} we can bound from above the r.h.s. of \eqref{up4} in the following way
\begin{equation}
\label{up7}
\bbE_{0,\r}\left[\int_{\bbR^d\setminus\{0\}}\xi(dx)c_{0,x}(x_1-\nabla_x f_{\z m}(\o))^2\mathds{1}_{H^c_x}\mathds{1}_{\{|E_x|> \z/\b\}}\right]\leq e^{-\z} c(\b).
\end{equation}


Consider now the second term in the sum in the r.h.s. of \eqref{up2}.
\begin{equation}
\label{up8}
\begin{split}
\bbE_{0,\r}\left[\int_{\bbR^d\setminus\{0\}}\xi(dx)c_{0,x}(x_1-\nabla_x f_{\z m}(\o))^2\mathds{1}_{H^c_x}\mathds{1}_{\{|E_x|\leq \z/\b\}}\right]
\end{split}
\end{equation}
Observe that for all $\d\in[0,1]$ 
$$c_{0,x}\leq e^{-(1-\d)\z}e^{-\d|x|}.$$
Then \eqref{up8} can be bounded from above in the following way
\begin{equation}
\label{up9}
\begin{split}
&\bbE_{0,\r}\left[\int_{\bbR^d\setminus\{0\}}\xi(dx)c_{0,x}(x_1-\nabla_x f_{\z m}(\o))^2\mathds{1}_{H^c_x}\mathds{1}_{\{|E_x|\leq \z/\b\}}\right]
\\ &\leq 2e^{-(1-\d)\z}\bbE_{0,\r}\left[\int_{\bbR^d\setminus\{0\}}\xi(dx)e^{-\d|x|}|x_1|^2\right]+2e^{-(1-\d)\z}\bbE_{0,\r}\left[\int_{\bbR^d\setminus\{0\}}\xi(dx)e^{-\d|x|}| f_{\z m}(\o))|^2\right]+
\\ &+  2e^{-(1-\d)\z}\bbE_{0,\r}\left[\int_{\bbR^d\setminus\{0\}}\xi(dx)e^{-\d|x|}| f_{\z m}(\tau_x\o))|^2\right].
\end{split}
\end{equation}
By \eqref{int4}, \eqref{int4.5} and \eqref{int1} we get
\begin{equation}
\label{up13}
\begin{split}
&\bbE_{0,\r}\left[\int_{\bbR^d\setminus\{0\}}\xi(dx)c_{0,x}(x_1-\nabla_x f_{\z m}(\o))^2\mathds{1}_{H^c_x}\mathds{1}_{\{|E_x|\leq \z/\b\}}\right]
\leq e^{-(1-\d)\z}c(\b).
\end{split}
\end{equation}
\medskip
\\So 
\begin{equation}
\label{up14}
\cC_1\leq (e^{-\z}+e^{-(1-\d)\z})c(\b).
\end{equation}
\bigskip
\\ \indent Consider now $\cC_2$. Observe that if $\{0,x\}\in \tilde\cE$, then $C_{\z,\b,\r}(0,\o)=C_{\z,\b,\r}(x,\o)$ and 
$$x_1-\nabla_x f_{\z m}(\o) =
\begin{cases}
0 &\text{if }\text{diam}(C_{\z,\b,\r}(0,\o))\leq \z m,
\\x_1 &\text{otherwise}.
\end{cases}$$
Moreover
\begin{equation*}
\mathds{1}_{\{\text{diam} (C_{\z,\b\r}(0,\o))>\z m\}}\mathds{1}_{\{|E_0|>\z/\b\}}=0.
\end{equation*}
Then, using the Cauchy-Schwarz's inequality and \eqref{int2}, we get
\begin{equation}
\label{up15}
\begin{split}
&\cC_2=\bbE_{0,\r}\left[\int_{\bbR^d\setminus\{0\}}\xi(dx)c_{0,x}(x_1-\nabla_x f_{\z m}(\o))^2\mathds{1}_{H_x}\right]
\\ &\leq\bbE_{0,\r}\left[\int_{\bbR^d\setminus\{0\}}\xi(dx)e^{-|x|}|x_1|^2\mathds{1}_{\{\text{diam} (C_{\z,\b\r}(0,\o))>\z m\}}\right]
\\ &\leq \left(\bbE_{0,\r}\left[\left(\int_{\bbR^d\setminus\{0\}}\xi(dx)e^{-|x|}|x_1|^2\right)^2\right]\right)^{\frac{1}{2}}\left(\bbE_{0,\r}\left[\mathds{1}_{\{\text{diam} (C_{\z,\b\r}(0,\o))>\z m\}}\right]\right)^{\frac{1}{2}}
\\ &\leq c\r\left(\bbE_{0,\r}\left[\mathds{1}_{\{\text{diam} (C_{\z,\b\r}(0,\o))>\z m\}}\mathds{1}_{\{|E_0|\leq\z/\b\}}\right]\right)^{\frac{1}{2}}
\\ &\leq c\r\left(\bbP_{0,\l}(\text{diam} (C_{1,1,\l}(0,\o))>m)\nu\left(\left[-\frac{\z}{\b},\frac{\z}{\b}\right]\right)\right)^{\frac{1}{2}}\leq e^{-c(\l) m}c(\b)
\end{split}
\end{equation}
where $c$ is a positive constant and $c(\l)$ is a positive constant depending only on $\l$. Note that in the last inequality we have used Theorem 1 in \cite{FAMH} as observed in \eqref{stimadiam}.
\bigskip
\\\indent So we get finally
\begin{equation}
\label{up16}
\begin{split}
\bbE_{0,\r}\left[\int_{\bbR^d} \xi(dx)c_{0,x}(x_1-\nabla_x f_{\z m}(\o))^2\right]\leq ( e^{-\z}+e^{-(1-\d)\z}+e^{-c(\l)m})c(\b).
\end{split}
\end{equation}
\\Now let $m$ go to $\infty$ and then do the same for $\b$. We get
\begin{equation}
\begin{split}
& \varlimsup_{\b\rgh\infty}\varlimsup_{m\rgh\infty}\b^{-\frac{\a+1}{\a+1+d}}\log\bbE_{0,\r}\left[\int_{\bbR^d} \xi(dx)c_{0,x}(x_1-\nabla_x f_{\z m}(\o))^2\right]
\\ &\leq \varlimsup_{\b\rgh\infty}\b^{-\frac{\a+1}{\a+1+d}}\log(e^{-\z}+e^{-(1-\d)\z}).
\end{split}
\end{equation}
Letting now $\d$ go to 0, we get
$$\varlimsup_{\b\rgh\infty}\varlimsup_{m\rgh\infty}\b^{-\frac{\a+1}{\a+1+d}}\log\bbE_{0,\r}\left[\int_{\bbR^d} \xi(dx)c_{0,x}(x_1-\nabla_x f_{\z m}(\o))^2\right]\leq -(1-\e)\left(\frac{\l_c}{\r}\right)^{\frac{1}{\a+1+\d}}.$$
Finally we can let $\e$ go to 0 and we obtain
$$\varlimsup_{\b\rgh\infty}\varlimsup_{m\rgh\infty}\b^{-\frac{\a+1}{\a+1+d}}\log\bbE_{0,\r}\left[\int_{\bbR^d} \xi(dx)c_{0,x}(x_1-\nabla_x f_{\z m}(\o))^2\right]\leq -\left(\frac{\l_c}{\r}\right)^{\frac{1}{\a+1+\d}}.$$
Therefore we have the thesis of proposition \ref{luce2}.
%%%%%


\subsection{Proof  Proposition \ref{luce1}}
 Fixed $m\in\bbN_+$, we consider the test function $u_{\z m}:\xi\times\O\rgh[0,1]$ defined as 
\begin{equation}
\label{testfun}
u_{\z m}(x,\o)=
\begin{cases}
0, &\text{if }x_1-f_{\z m}(\tau_x\o)<-n,
\\\frac{n+x_1-f_{\z m}(\tau_x \o)}{2n}, &\text{if }-n\leq x_1-f_{\z m}(\tau_x\o)\leq n,
\\1, &\text{if }x_1-f_{\z m}(\tau_x\o)>n,
\end{cases}
\end{equation}
where $f_{\z m}$ is the function defined in \eqref{ftest}.
\\In order to prove proposition \ref{luce1}, we have to study
\begin{equation}
\label{uc3}
\begin{split}
&n^{2-d}\sum_{x\in \L_n\cap \xi}\sum_{\substack{y\in D_n\cap\xi:\\ x\neq y}}c_{x,y}(\nabla_{x,y}u_{\z m})^2=\frac{1}{n^d}\sum_{x\in \L_n\cap \xi}\sum_{\substack{y\in D_n\cap\xi:\\x\neq y}}c_{x,y}(2n\cdot\nabla_{x,y}u_{\z m})^2=
\\ &=\frac{\xi(\L_n)}{n^d}\cdot \frac{1}{\xi(\L_n)}\sum_{x\in \L_n\cap \xi}\sum_{\substack{y\in D_n\cap\xi:\\x\neq y}}c_{x,y}(2n\cdot\nabla_{x,y}u_{\z m})^2.
\end{split}
\end{equation}
Since 
$$\lim_{n\rgh \infty}\frac{\xi(\L_n)}{n^d}=c(\r)\qquad \bbP_{0,\r}-\text{a.s.}$$
where $c(\r)>0$ is a constant depending only on $\r$, is sufficient to look for an upper bound for
\begin{equation}
\label{uc3}
\frac{1}{\xi(\L_n)}\sum_{x\in \L_n\cap \xi}\sum_{\substack{y\in D_n\cap\xi:\\x\neq y}}c_{x,y}(2n\cdot\nabla_{x,y}u_{\z m})^2.
\end{equation}
To simplify the notation, we will write the second sum without specifying $x\neq y$. 
\\Define the event $H_{xy}:=\{\{x,y\}\in\tilde \cE\}$. 
\\Observe that for every $x,y\in D_n$ such that $H_{xy}$ does not occur
\begin{equation}
\label{stimagrad1}
2n\cdot|\nabla_{x,y}u_{\z m}|\leq |x_1-y_1|+|f_{\z m}(\tau_x\o)|+|f_{\z m}(\tau_y\o)|,
\end{equation}
and then
\begin{equation}
\label{stimagrad2}
4n^2\cdot|\nabla_{x,y}u_{\z m}|^2\leq 2(|x_1-y_1|^2+|f_{\z m}(\tau_x\o)|^2+|f_{\z m}(\tau_y\o)|^2).
\end{equation}
Otherwise, for every $x,y\in D_n$ such that $H_{xy}$ occurs
\begin{equation}
\label{stimagrad3}
2n\cdot|\nabla_{x,y}u_{\z m}|\leq |x_1-y_1|\mathds{1}_{\{\text{diam}(C_{\z,\b,\r}(x,\o))>\z m\}}.
\end{equation}
So \eqref{uc3} can be bounded from above by
\begin{equation}
\label{uc5}
\begin{split}
&\frac{1}{\xi(\L_n)}\sum_{x\in \L_n\cap \xi}\sum_{y\in D_n\cap\xi}c_{x,y}|x_1-y_1|^2\mathds{1}_{\{\text{diam}(C_{\z,\b,\r}(x,\o))>\z m\}}\mathds{1}_{H_{xy}}+\\ &+\frac{2}{\xi(\L_n)}\sum_{x\in \L_n\cap \xi}\sum_{y\in D_n\cap\xi}c_{x,y}(|x_1-y_1|^2+|f_{\z m}(\tau_x\o)|^2+|f_{\z m}(\tau_y\o)|^2)\mathds{1}_{H_{xy}^c}
\end{split}
\end{equation}
Let us study the first sum in \eqref{uc5}
\begin{equation}
\label{uc6}
\begin{split}
&\frac{1}{\xi(\L_n)}\sum_{x\in \L_n\cap \xi}\sum_{y\in D_n\cap\xi}c_{x,y}|x_1-y_1|^2\mathds{1}_{\{\text{diam}(C_{\z,\b,\r}(x,\o))>\z m\}}\mathds{1}_{H_{xy}}
\\ & \leq \frac{1}{\xi(\L_n)}\sum_{x\in \L_n\cap \xi}\sum_{y\in D_n\cap\xi}e^{-|x-y|}|x_1-y_1|^2\mathds{1}_{\{\text{diam}(C_{\z,\b,\r}(x,\o))>\z m\}}
\\ &\leq\frac{1}{\xi(\L_n)}\sum_{x\in \L_n\cap \xi}\sum_{y\in \xi}e^{-|x-y|}|x_1-y_1|^2\mathds{1}_{\{\text{diam}(C_{\z,\b,\r}(x,\o))>\z m\}}=\frac{1}{\xi(\L_n)}\sum_{x\in \L_n\cap \xi}h(\tau_x\xi),
\end{split}
\end{equation}
where $h(\xi)=\sum_{y\in \xi}e^{-|y|}|y_1|^2\mathds{1}_{\{\text{diam}(C_{\z,\b,\r}(0,\o))>\z m\}}$.
\\By the ergodic theorem, we get that
\begin{equation}
\label{c1-4}
\lim_{n\rgh\infty}\frac{1}{\xi(\L_n)}\sum_{x\in \L_n\cap \xi}h(\tau_x\xi)=\bbE_{0,\r}[h(\xi)]\qquad \bbP_{0,\r}-\text{a.s.}
\end{equation}
By Cauchy-Schwarz's inequality and \eqref{int2}, we obtain
\begin{equation}
\label{uc7}
\begin{split}
&\bbE_{0,\r}[h(\xi)]= \bbE_{0,\r}\left[\int_{\bbR^d\setminus\{0\}}\xi(dy)e^{-|y|}|y_1|^2\mathds{1}_{\{\text{diam}(C_{\z,\b,\r}(0,\o))>\z m\}}\right]
\\ &\leq\left(\bbE_{0,\r}\left[\left(\int_{\bbR^d\setminus\{0\}}\xi(dy)e^{-|y|}|y_1|^2\right)^2\right]\right)^{\frac{1}{2}}\bbP_{0,\r}(\text{diam}(C_{\z,\b,\r}(0,\o))>\z m)^\frac{1}{2}
\\ &\leq c(\r)\bbP_{0,\r}(\text{diam}(C_{\z,\b,\r}(0,\o))>\z m\,|\,|E_0|\leq\z/\b)^\frac{1}{2}\nu\left(\left[-\frac{\z}{\b},\frac{\z}{\b}\right]\right)^\frac{1}{2}=
\\&=c(\r,\z,\b)\bbP_{0,\l}(\text{diam}(C_{1,1,\l}(0,\o))> m)^\frac{1}{2}\leq c(\r,\z,\b)e^{-c(\l)m},
\end{split}
\end{equation}
where $c(\r),c(\r,\z,\b)$ and $c(\l)$ are positive constants that have a polynomial-dependence on their arguments.
\\So 
\begin{equation}
\label{uc8}
\varlimsup_{m\rgh\infty}\varlimsup_{n\rgh\infty}\frac{1}{\xi(\L_n)}\sum_{x\in \L_n\cap \xi}\sum_{y\in D_n\cap\xi}c_{x,y}|x_1-y_1|^2\mathds{1}_{\{\text{diam}(C_{\z,\b,\r}(x,\o))>\z m\}}=0\qquad\bbP_{0,\r}-\text{a.s.}
\end{equation}
Consider now the second term of \eqref{uc5}
\begin{equation}
\label{uc4}
\begin{split}
&\frac{2}{\xi(\L_n)}\sum_{x\in \L_n\cap \xi}\sum_{y\in D_n\cap\xi}c_{x,y}(|x_1-y_1|^2+|f_{\z m}(\tau_x\o)|^2+|f_{\z m}(\tau_y\o)|^2)\mathds{1}_{H_{xy}^c}
\\ & \leq\frac{2}{\xi(\L_n)}\sum_{x\in \L_n\cap \xi}\sum_{y\in D_n\cap\xi}c_{x,y}(|x_1-y_1|^2+|f_{\z m}(\tau_x\o)|^2+|f_{\z m}(\tau_y\o)|^2)=
\\&= \frac{2}{\xi(\L_n)}\sum_{x\in \L_n\cap \xi}\sum_{y\in D_n\cap\xi}c_{x,y}|x_1-y_1|^2+\frac{2}{\xi(\L_n)}\sum_{x\in \L_n\cap \xi}\sum_{y\in D_n\cap\xi}c_{x,y}|f_{\z m}(\tau_x\o)|^2
\\ &+\frac{2}{\xi(\L_n)}\sum_{x\in \L_n\cap \xi}\sum_{y\in D_n\cap\xi}c_{x,y}|f_{\z m}(\tau_y\o)|^2=\cC_1+\cC_2+\cC_3,
\end{split}
\end{equation}
where
$$\cC_1=\frac{2}{\xi(\L_n)}\sum_{x\in \L_n\cap \xi}\sum_{y\in D_n\cap\xi}c_{x,y}|x_1-y_1|^2,$$
$$\cC_2=\frac{2}{\xi(\L_n)}\sum_{x\in \L_n\cap \xi}\sum_{y\in D_n\cap\xi}c_{x,y}|f_{\z m}(\tau_x\o)|^2,$$
$$\cC_3=\frac{2}{\xi(\L_n)}\sum_{x\in \L_n\cap \xi}\sum_{y\in D_n\cap\xi}c_{x,y}|f_{\z m}(\tau_y\o)|^2.$$
In order to compute an upper bound for these terms we define the event $T_x:=\{|E_x|\leq \z/\b\}$.
\\We start giving an upper bound for $\cC_1$
\begin{equation}
\label{c1-1}
\begin{split}
\cC_1=\frac{2}{\xi(\L_n)}\sum_{x\in \L_n\cap \xi}\sum_{y\in D_n\cap\xi}c_{x,y}|x_1-y_1|^2\mathds{1}_{T_x}+\frac{2}{\xi(\L_n)}\sum_{x\in \L_n\cap \xi}\sum_{y\in D_n\cap\xi}c_{x,y}|x_1-y_1|^2\mathds{1}_{T_x^c}.
\end{split}
\end{equation}
Observe that if $T_x$ occurs, then $c_{x,y}\leq e^{-(1-\d)\z}e^{-\d|x-y|}$, for any $\d\in[0,1]$. Then
\begin{equation}
\label{c1-3}
\begin{split}
&\frac{2}{\xi(\L_n)}\sum_{x\in \L_n\cap \xi}\sum_{y\in D_n\cap\xi}c_{x,y}|x_1-y_1|^2\mathds{1}_{T_x}\leq 
\frac{2e^{-(1-\d)\z}}{\xi(\L_n)}\sum_{x\in \L_n\cap \xi}\sum_{y\in D_n\cap\xi}e^{-\d|x-y|}|x_1-y_1|^2=
\\&\leq \frac{2e^{-(1-\d)\z}}{\xi(\L_n)}\sum_{x\in \L_n\cap \xi}\sum_{y\in \xi}e^{-\d|x-y|}|x_1-y_1|^2=\frac{2e^{-(1-\d)\z}}{\xi(\L_n)}\sum_{x\in \L_n\cap \xi}h(\tau_x\xi),
\end{split}
\end{equation}
where $h(\xi)=\sum_{y\in \xi}e^{-\d|y|}|y_1|^2$. 
Since, by \eqref{int1}, we get
\begin{equation}
\label{c1-5}
\bbE_{0,\r}[h(\xi)]=\bbE_{0,\r}\left[\int_{\bbR^d\setminus\{0\}}\xi(dy)e^{-\d|y|}|y_1|^2\right]=c(\r,\d)
\end{equation}
where $c(\r,\d)$ is a constant that has a polynomial-dependence on its argumets, then by \eqref{c1-4} we obtain
\begin{equation}
\label{c1-6}
\varlimsup_{n\rgh\infty}\frac{2}{\xi(\L_n)}\sum_{x\in \L_n\cap \xi}\sum_{y\in D_n\cap\xi}c_{x,y}|x_1-y_1|^2\mathds{1}_{T_x}\leq c(\r,\d)e^{-(1-\d)\z}.
\end{equation}
As far as the other sum in the r.h.s. of \eqref{c1-1}, since $T^c_x$ occurs, then $c_{xy}\leq e^{-|x-y|-\z}$. So we can proceed as follows
\begin{equation}
\label{c1-7}
\begin{split}
&\frac{2}{\xi(\L_n)}\sum_{x\in \L_n\cap \xi}\sum_{y\in D_n\cap\xi}c_{x,y}|x_1-y_1|^2\mathds{1}_{T^c_x}\leq
\frac{2e^{-\z}}{\xi(\L_n)}\sum_{x\in \L_n\cap \xi}\sum_{y\in D_n\cap\xi}e^{-|x-y|}|x_1-y_1|^2
\\ &\leq \frac{2e^{-\z}}{\xi(\L_n)}\sum_{x\in \L_n\cap \xi}\sum_{y\in\xi}e^{-|x-y|}|x_1-y_1|^2 =\frac{2e^{-\z}}{\xi(\L_n)}\sum_{x\in \L_n\cap \xi}h(\tau_x\xi),
\end{split}
\end{equation}
where $h(\xi)=\sum_{y\in \xi}e^{-|y|}|y_1|^2$. By \eqref{c1-4} and reasoning as in \eqref{c1-5}, we get that $\bbE_{0,\r}[h(\xi)]=c(\r)$, where $c(\r)$ is a positive constant that has a polynomial-dependence on $\r$. So by \eqref{c1-4} and \eqref{c1-7} we get
\begin{equation}
\label{c1-8}
\varlimsup_{n\rgh\infty}\frac{2}{\xi(\L_n)}\sum_{x\in \L_n\cap \xi}\sum_{y\in D_n\cap\xi}c_{x,y}|x_1-y_1|^2\mathds{1}_{T^c_x}\leq c(\r)e^{-\z}.
\end{equation}
Then
\begin{equation}
\label{c1-9}
\varlimsup_{n\rgh\infty}\cC_1\leq c(\r,\d)(e^{-(1-\d)\z}+e^{-\z}).
\end{equation}
Consider now $\cC_2$
\begin{equation}
\label{c2-1}
\cC_2=\frac{2}{\xi(\L_n)}\sum_{x\in \L_n\cap \xi}\sum_{y\in D_n\cap\xi}c_{x,y}|f_{\z m}(\tau_x\o)|^2\mathds{1}_{T_x}+\frac{2}{\xi(\L_n)}\sum_{x\in \L_n\cap \xi}\sum_{y\in D_n\cap\xi}c_{x,y}|f_{\z m}(\tau_x\o)|^2\mathds{1}_{T^c_x}.
\end{equation}
Observe that if $T_x$ does not occur, then $f_{\z m}(\tau_x \o)=0$. So
\begin{equation}
\label{c2-2}
\begin{split}
&\cC_2=\frac{2}{\xi(\L_n)}\sum_{x\in \L_n\cap \xi}\sum_{y\in D_n\cap\xi}c_{x,y}|f_{\z m}(\tau_x\o)|^2\mathds{1}_{T_x}\leq\frac{2e^{-(1-\d)\z}}{\xi(\L_n)}\sum_{x\in \L_n\cap \xi}\sum_{y\in D_n\cap\xi}e^{-\d|x-y|}|f_{\z m}(\tau_x\o)|^2
\\ &\leq \frac{2e^{-(1-\d)\z}}{\xi(\L_n)}\sum_{x\in \L_n\cap \xi}\sum_{y\in\xi}e^{-\d|x-y|}|f_{\z m}(\tau_x\o)|^2=\frac{2e^{-(1-\d)\z}}{\xi(\L_n)}\sum_{x\in \L_n\cap \xi}h(\tau_x\xi),
\end{split}
\end{equation}
where $h(\xi)=\sum_{y\in \xi}e^{-\d|y|}|f_{\z m}(\o)|^2$. 
\\By \eqref{int5} and \eqref{c1-4}, we get
\begin{equation}
\label{c2-3}
\lim_{n\rgh\infty}\frac{2e^{-(1-\d)\z}}{\xi(\L_n)}\sum_{x\in \L_n\cap \xi}h(\tau_x\xi)=\bbE_{0,\r}[h(\xi)]=c(\l,\z,\b,\d,\r)e^{-(1-\d)\z}\qquad \bbP_{0,\r}-\text{a.s.}
\end{equation}
and then
\begin{equation}
\label{c2-4}
\varlimsup_{n\rgh\infty}\cC_2\leq c(\l,\z,\b,\d,\r)e^{-(1-\d)\z}.
\end{equation}
Consider now $\cC_3$
\begin{equation}
\label{c3-1}
\cC_3=\frac{2}{\xi(\L_n)}\sum_{x\in \L_n\cap \xi}\sum_{y\in D_n\cap\xi}c_{x,y}|f_{\z m}(\tau_y\o)|^2\mathds{1}_{T_y}+\frac{2}{\xi(\L_n)}\sum_{x\in \L_n\cap \xi}\sum_{y\in D_n\cap\xi}c_{x,y}|f_{\z m}(\tau_y\o)|^2\mathds{1}_{T^c_y}.
\end{equation}
Observe that if $T_y$ does not occur, then $f_{\z m}(\tau_y\o)=0$. So
\begin{equation}
\label{c3-1}
\begin{split}
&\cC_3=\frac{2}{\xi(\L_n)}\sum_{x\in \L_n\cap \xi}\sum_{y\in D_n\cap\xi}c_{x,y}|f_{\z m}(\tau_y\o)|^2\mathds{1}_{T_y}
\\ &\leq \frac{2e^{-(1-\d)\z}}{\xi(\L_n)}\sum_{x\in \L_n\cap \xi}\sum_{y\in D_n\cap\xi}e^{-\d|x-y|}|f_{\z m}(\tau_y\o)|^2.
\end{split}
\end{equation}
Define the family of sets $\L_n^i:=[-n+2in,n+2in]\times[-n,n]^{d-1}$ for $i\in\bbZ$. Then $D_n=\cup_{i\in\bbZ}\L_n^i$ and we can rewrite the r.h.s. of \eqref{c3-1} as
\begin{equation}
\label{c3-2}
\begin{split}
&2e^{-(1-\d)\z}\Big( \frac{1}{\xi(\L_n)}\sum_{x\in \L_n\cap \xi}\sum_{i\in\bbZ\setminus\{0\}}\sum_{y\in\L_n^i\cap\xi}e^{-\d|x-y|}|f_{\z m}(\tau_y\o)|^2+
\\ &+ \frac{1}{\xi(\L_n)}\sum_{x\in \L_n\cap \xi}\sum_{y\in \L_n\cap\xi}e^{-\d|x-y|}|f_{\z m}(\tau_y\o)|^2\Big)
\end{split}
\end{equation}
We consider separately the two sums in \eqref{c3-2}
\begin{equation}
\label{c3-3}
\begin{split}
&\frac{1}{\xi(\L_n)}\sum_{x\in \L_n\cap \xi}\sum_{i\in\bbZ\setminus\{0\}}\sum_{y\in\L_n^i\cap\xi}e^{-\d|x-y|}|f_{\z m}(\tau_y\o)|^2=
\\ &=\frac{\xi(\L_n^i)}{\xi(\L_n)}\frac{1}{\xi(\L_n^i)}\sum_{i\in\bbZ\setminus\{0\}}\sum_{y\in\L_n^i\cap\xi}\sum_{x\in \L_n\cap \xi}e^{-\d|x-y|}|f_{\z m}(\tau_y\o)|^2.
\end{split}
\end{equation}
Observe that since $x\in \L_n$ and $y\in\L_n^i$, then for all $\g\in[0,1]$ there holds 
$$|x-y|\geq(1-\g)  (|i|-1)+\g|x-y|\qquad\forall\,i\in\bbZ\setminus\{0\}.$$
So the r.h.s. of \eqref{c3-3} can be bounded from above by
\begin{equation}
\label{c3-4}
\begin{split}
&\frac{\xi(\L_n^i)}{\xi(\L_n)}\frac{1}{\xi(\L_n^i)}\sum_{i\in\bbZ\setminus\{0\}}e^{-\d(1-\g)  (|i|-1)}\sum_{y\in\L_n^i\cap\xi}\sum_{x\in \xi}e^{-\d \g|x-y|}|f_{\z m}(\tau_y\o)|^2=
\\ &=\frac{\xi(\L_n^i)}{\xi(\L_n)}\frac{1}{\xi(\L_n^i)}\sum_{i\in\bbZ\setminus\{0\}}e^{-\d(1-\g)  (|i|-1)}\sum_{y\in\L_n^i\cap\xi}h(\tau_y\o),
\end{split}
\end{equation}
where $h(\o):=\sum_{x\in \xi}e^{-\d\g|x|}|f_{\z m}(\o)|^2$.
By \eqref{c1-4}, since as $n\rgh\infty$
$$\frac{\xi(\L_n^i)}{\xi(\L_n)}\rgh 1\quad \bbP_{0,\r}-\text{a.s.},$$
we get
\begin{equation}
\label{c3-5}
\begin{split}
&\lim_{n\rgh\infty}\frac{\xi(\L_n^i)}{\xi(\L_n)}\frac{1}{\xi(\L_n^i)}\sum_{i\in\bbZ\setminus\{0\}}e^{-\d(1-\g) (|i|-1)}\sum_{y\in\L_n^i\cap\xi}\sum_{x\in\xi}e^{-\d\g|x-y|}|f_{\z m}(\tau_y\o)|^2=
\\&=\sum_{i\in\bbZ\setminus\{0\}}e^{-\d(1-\g) (|i|-1)}\bbE_{0,\r}\left[h(\o)\right]\qquad\bbP_{0,\r}-\text{a.s.}
\end{split}
\end{equation}
But by \eqref{int3} and \eqref{int5} we get
\begin{equation}
\label{c3-6}
\begin{split}
&\bbE_{0,\r}[h(\o)]=\bbE_{0,\r}\left[ \int_{\bbR^d}\xi(dx)e^{-\d\g|x|}|f_{\z m}(\o)|^2\right]=
\\ &=\bbE_{0,\r}\left[ \int_{\bbR^d\setminus\{0\}}\xi(dx)e^{-\d\g|x|}|f_{\z m}(\o)|^2\right]+\bbE_{0,\r}[|f_{\z m}(\o)|^2]
\\ &\leq c(\r,\l,\d,\g,\z,\b),
\end{split}
\end{equation}
where $c(\r,\l,\d,\g,\z,\b)$ is a positive constant that has a polynomial dependence on its arguments.
\\So the r.h.s. of \eqref{c3-5} can be bounded from above by 
$$c(\r,\l,\d,\g,\z,\b)\sum_{i\in\bbZ\setminus\{0\}}e^{-\d(1-\g) (|i|-1)}<\infty$$
and then 
\begin{equation}
\label{c3-7}
\varlimsup_{n\rgh\infty}\frac{1}{\xi(\L_n)}\sum_{x\in \L_n\cap \xi}\sum_{i\in\bbZ\setminus\{0\}}\sum_{y\in\L_n^i\cap\xi}e^{-\d|x-y|}|f_{\z m}(\tau_y\o)|^2\leq c(\r,\l,\d,\g,\z,\b).
\end{equation}
We can remove the $\g$-dependence of $c(\r,\l,\d,\g,\z,\b)$ by taking $\g=1$.
\\Consider now the second sum in \eqref{c3-2}
\begin{equation}
\label{c3-8}
\begin{split}
&\frac{1}{\xi(\L_n)}\sum_{x\in \L_n\cap \xi}\sum_{y\in \L_n\cap\xi}e^{-\d|x-y|}|f_{\z m}(\tau_y\o)|^2=\frac{1}{\xi(\L_n)}\sum_{y\in \L_n\cap\xi}\sum_{x\in \L_n\cap \xi}e^{-\d|x-y|}|f_{\z m}(\tau_y\o)|^2
\\ &\leq \frac{1}{\xi(\L_n)}\sum_{y\in \L_n\cap\xi}\sum_{x\in \xi}e^{-\d|x-y|}|f_{\z m}(\tau_y\o)|^2=\frac{1}{\xi(\L_n)}\sum_{y\in \L_n\cap\xi}h(\tau_y\o),
\end{split}
\end{equation}
where $h(\o)=\sum_{x\in  \xi}e^{-\d|x|}|f_{\z m}(\o)|^2$. By \eqref{c1-4} and \eqref{int5} we get that 
\begin{equation}
\label{c3-9}
\lim_{n\rgh\infty}\frac{1}{\xi(\L_n)}\sum_{y\in \L_n\cap\xi}h(\tau_y\o)=\bbE_{0,\r}[h(\o)]\leq c(\z,\b,\r,\l,\d).
\end{equation}
Then 
\begin{equation}
\label{c3-10}
\varlimsup_{n\rgh\infty}\cC_3\leq e^{-(1-\d)\z}c(\z,\b,\r,\l,\d),
\end{equation}
where $c(\z,\b,\r,\l,\d)$ is a positive constant that has a polynomial-dependence on its arguments.
\\Then we get that
\begin{equation}
\label{upfin1}
\varlimsup_{n\rgh\infty}\cC_1+\cC_2+\cC_3\leq (e^{-\z}+e^{-(1-\d)\z})c(\z,\b,\d,\r,\l).
\end{equation}
Since $\d$ can be choosen arbitrarily in $[0,1]$, we take $\d=1$ and, by \eqref{uc3}, \eqref{uc5}, \eqref{uc8}, \eqref{uc4} and \eqref{upfin1}, we get
\begin{equation}
\label{upfin2}
\varlimsup_{m\rgh\infty}\varlimsup_{n\rgh\infty}n^{2-d}\sum_{x\in\L_n\cap\xi}\sum_{y\in D_n\cap\xi}c_{x,y}(\nabla_{x,y}u_{\z m})^2\leq e^{-\z}c(\z,\b,\r,\l)\qquad \bbP_{0,\r}-\text{a.s.}
\end{equation}
where we remember that $c(\z,\b,\r,\l)$ is a positive constant that has a polynomial-dependence on its arguments. Then $\bbP_{0,\r}$-a.s.
\begin{equation}
\label{upfin3}
\begin{split}
&\varlimsup_{\b \to \infty} \varlimsup _{n \to \infty} \beta ^{-\frac{\a+1}{\a+1+d}}\log [ n ^{2-d}  \s_n (\b,\o)]
\\&\leq\varlimsup_{\b\rgh\infty}\varlimsup_{m\rgh\infty}\varlimsup_{n\rgh\infty}\b^{-\frac{\a+1}{\a+1+d}}\log\left[n^{2-d}\sum_{x\in\L_n\cap\xi}\sum_{y\in D_n\cap\xi}c_{x,y}(\nabla_{x,y}u_{\z m})^2\right]
\\&\leq -(1-\e)\left(\l_c/\r\right)^\frac{1}{\a+1+d}.
\end{split}
\end{equation}
Letting $\e\rgh0$ we obtain
\begin{equation}
\label{upfin4}
\varlimsup_{\b \to \infty} \varlimsup _{n \to \infty} \beta ^{-\frac{\a+1}{\a+1+d}}\log [ n ^{2-d}  \s_n (\b,\o)]\leq-\left(\l_c/\r\right)^\frac{1}{\a+1+d},
\end{equation}
and therefore we get the thesis of proposition \ref{luce1}.

\appendix
\section{Computation of some terms in the conductivity fromula \eqref{condcompl}}
We remember that $u_{\z m}\Big|_{A_n}=0$ and $u_{\z m}\Big|_{B_n}=1$. Then
$$\sum_{x\in A_n\cap\xi}\sum_{y\in B_n\cap\xi}c_{x,y}(\nabla_{x,y} u_{\z m})^2=\sum_{x\in A_n\cap\xi}\sum_{y\in B_n\cap\xi}c_{x,y}.$$
The following lemma gives the asymptotic behaviour of this term for $n\rgh\infty$.
\begin{Lemma}
\label{ext}
$$\lim_{n\rgh\infty} n^{2-d}\sum_{x\in A_n\cap\xi}\sum_{y\in B_n\cap\xi}c_{x,y}=0,\quad \bbP_{0,\r}-\text{a.s.}$$
\end{Lemma}
\begin{proof}
For $i=0,1,2,\ldots$ define the sets
$$A_n^i=\{x\in A_n\,|\,-n-i-1\leq x_1\leq-n-i\},\quad B_n^i=\{x\in B_n\,|\,n+i\leq x_1\leq n+i+1\}$$
Then
\begin{equation}
\label{negl1}
\sum_{x\in A_n\cap\xi}\sum_{y\in B_n\cap\xi}c_{x,y}=\sum_{i=0}^\infty \sum_{j=0}^\infty\sum_{x\in A_n^i\cap\xi}\sum_{y\in B_n^j\cap\xi} c_{x,y}
\end{equation}
Observe that if $x\in A_n^i\cap\xi$ and $y\in B_n^j\cap\xi$, then $|x-y|\geq 2n+i+j$ and consequentely $c_{x,y}\leq e^{-2n-i-j}$. Then r.h.s in \eqref{negl1} can be bounded from above by
\begin{equation}
e^{-2n}\sum_{i=0}^\infty \sum_{j=0}^\infty|\xi(A_n^i)|\cdot|\xi(B_n^j)|e^{-(i+j)}=e^{-2n}\left(\sum_{i=0}^\infty|\xi(A_n^i)|e^{-i}\right)\left(\sum_{j=0}^\infty|\xi(B_n^j)|e^{-j}\right).
\end{equation}
Define
$$\cC_n:=n^{2-d}e^{-2n}\left(\sum_{i=0}^\infty|\xi(A_n^i)|e^{-i}\right)\left(\sum_{j=0}^\infty|\xi(B_n^j)|e^{-j}\right).$$
By Markov's inequality we get for $s\in(0,+\infty)$
\begin{equation}
\label{negl2}
\bbP_{0,\r}\left(\cC_n>\frac{1}{s}\right)\leq s\cdot\bbE_{0,\r}\left[\cC_n\right]
\end{equation}
Let us compute $\bbE_{0,\r}[\cC_n]$. Since $A^i_n$ e $B^j_n$ are disjoint for all $i,j=0,1\ldots$, then
$$\bbE_{0,\r}[\cC_n]=e^{-2n}(c^2n^{d}\r^2).$$
where $0<c=\sum_{i=0}^\infty e^{-i}<\infty$. So
\begin{equation}
\label{dec}
\bbP_{0,\r}\left(\cC_n>\frac{1}{s}\right)\leq s\cdot e^{-2n}(c^2n^{d}\r^2).
\end{equation}
Then by \eqref{negl1} and \eqref{negl2}
$$\bbP_{0,\r}\left(\sum_{x\in A_n\cap\xi}\sum_{y\in B_n\cap\xi}c_{x,y}>\frac{1}{s}\right)=\bbP_{0,\r}\left(\cC_n>\frac{1}{s}\right)\leq s\cdot e^{-2n}(c^2n^{d}\r^2).$$
Let $n$ go to infinity, we get
$$\bbP_{0,\r}\left(\lim_{n\rgh\infty}\sum_{x\in A_n\cap\xi}\sum_{y\in B_n\cap\xi}c_{x,y}>\frac{1}{s}\right)=0\quad \forall \,s\in(0,+\infty),$$
that is
$$\bbP_{0,\r}\left(\lim_{n\rgh\infty}\sum_{x\in A_n\cap\xi}\sum_{y\in B_n\cap\xi}c_{x,y}=0\right)=1.$$
\end{proof}



\section{Some computation rules for Poisson point process}
Let $g:\bbR^d\times\cN\rgh\bbR$ be a measurable function, where $\cN$ is the space of all the possible realizations for the point process $\xi$. 
%Then, formula (13.1.5) in \cite{DVJ} applied in our case tells that
%\begin{equation}
%\label{f1}
%\bbE_\r\left[\int_{\bbR^d}\xi(dx)g(x,\xi)\right]=\r\int_{\bbR^d}dx E_x[g(x,\xi)],
%\end{equation}
%where $E_x$ is the average computed with respect to the probability measure $P_x:=\bbP_\r*\d_x$ and $\d_x$ is the Dirac's delta distribution with centered at $x\in\bbR^d$.
\\By formulas (12.1.6) and formula (12.1.12) in \cite{DVJ} we get
\begin{equation}
\label{f1}
\bbE_\r\left[\int_{\bbR^d}g(x,\xi)\xi(dx)\right]=\r\int_{\bbR^d}dx\, \bbE_\r[g(x,\xi\cup\{x\})].
\end{equation}
We want to apply \eqref{f1} for the computation of some integrals that will recur many times.
\medskip
\\\textbf{First application:} $g(x,\xi)=h(x)$, where $h:\bbR^d\rgh\bbR$
\begin{equation}
\label{int1}
\bbE_\r\left[\int_{\bbR^d}h(x)\xi(dx)\right]=\r\int_{\bbR^d}dx\, h(x).
\end{equation}
\medskip
\textbf{Second application:} $g(x,\xi)=\left(\int_{\bbR^d}\xi(dy)h(y)\right)h(x)$, where $h:\bbR^d\rgh\bbR$
\begin{equation}
\label{int2}
\begin{split}
&\bbE_\r\left[\left(\int_{\bbR^d}h(x)\xi(dx)\right)^2\right]=\bbE_\r\left[\int_{\bbR^d}g(x,\xi)\xi(dx)\right]=
\\ &=\r\int_{\bbR^d}dx\bbE_\r[g(x,\xi\cup\{x\})]=\r\int_{\bbR^d}dx\bbE_\r\left[\left(h(x)+\int_{\bbR^d}\xi(dy)h(y)\right)h(x)\right]=
\\ &=\r\int_{\bbR^d}dx\, h^2(x)+\r\int_{\bbR^d}dx\, h(x)\bbE_\r\left[\int_{\bbR^d}\xi(dy)g(y)\right]=
\\ &=\r\int_{\bbR^d}dx\, h^2(x)+\r^2\left(\int_{\bbR^d}dx\,h(x)\right)^2.
\end{split}
\end{equation}
 

\end{document}
